% Leçon 22 — Expression écrite : culture chinoise et camerounaise (你喜欢喀麦隆还是中国艺术？为什么？) — Chinois 1ère A — Cours signé OnBuch+
% Programme officiel MINESEC. Compiler avec : tectonic ZHA22-ecrit-culture-chinoise-et-camerounaise.tex
\newcommand{\DOCMATIERE}{Chinois}
\newcommand{\DOCNIVEAU}{1ère A}
\newcommand{\DOCMODULE}{Module 3 : Citoyenneté et ouverture au monde}
\newcommand{\DOCLECON}{ZHA22}
\newcommand{\DOCTITRE}{Leçon 22 — Expression écrite : culture chinoise et camerounaise (你喜欢喀麦隆还是中国艺术？为什么？)}
% OnBuch+ — préambule « cours signés » ENRICHI (compilé avec tectonic / XeLaTeX)
% Système visuel « Pop » d'OnBuch+ : fond crème (jamais blanc), bordures encre
% épaisses, ombres « dures » décalées, titres Archivo Black en capitales.
% Version MATHÉMATIQUES : graphes (pgfplots/tikz), boîtes pédagogiques
% (définition, propriété, méthode, attention, le savais-tu, exercice résolu,
% exercice, TP, bilan), page de garde et signature OnBuch+ ancrée en bas.
%
% Macros attendues dans main.tex AVANT \input{preamble} :
%   \DOCMATIERE  \DOCNIVEAU  \DOCMODULE  \DOCLECON (numéro)  \DOCTITRE
\documentclass[11pt]{article}

\usepackage{fontspec}
\usepackage[french]{babel} % français = langue principale ; le chinois via \zh{..} / textebox (xeCJK)
\usepackage{xeCJK}
\usepackage{pifont} % \ding{51} \ding{55} utilisés par les modèles
\usepackage[a4paper,margin=2cm,top=3.4cm,bottom=2.8cm,headheight=1.4cm,footskip=1.3cm]{geometry}
\usepackage{amsmath,amssymb,amsfonts}
\usepackage{mathtools} % \xmapsto, \coloneqq... souvent utilisés par les modèles
\usepackage{cancel} % simplifications barrées (bilans de forces)
% \abs{z} (module, valeur absolue) et \norm{u} (norme), utilisés par les modèles
\providecommand{\abs}{}\let\abs\relax\DeclarePairedDelimiter{\abs}{\lvert}{\rvert}
\providecommand{\norm}{}\let\norm\relax\DeclarePairedDelimiter{\norm}{\lVert}{\rVert}
\usepackage[table]{xcolor} % [table] : fournit aussi \rowcolors (alternance de lignes)
\usepackage{pagecolor}
\usepackage{tikz}
\usetikzlibrary{arrows.meta,positioning,calc,shapes.geometric,shapes.misc,shapes.arrows,decorations.pathmorphing,decorations.pathreplacing,decorations.markings,patterns,shadows,mindmap,trees,fit,backgrounds,matrix,intersections,angles,quotes}
\usepackage{pgfplots}
\usepackage[version=4]{mhchem} % équations nucléaires \ce{^{235}_{92}U}
\usepackage[european,siunitx]{circuitikz} % schémas électriques
\pgfplotsset{compat=1.18}
\usepgfplotslibrary{fillbetween,groupplots} % aires entre courbes (calcul intégral)
\usepackage{enumitem}
\usepackage{fancyhdr}
% [most] charge toutes les bibliothèques tcolorbox (listings, minted, raster,
% documentation...) et gonfle inutilement la mémoire TeX sur un document long
% (~300 boîtes) : on ne charge que ce qui est réellement utilisé.
\usepackage[skins,breakable]{tcolorbox}
\usepackage{graphicx}
\usepackage{colortbl}
\usepackage{booktabs}
\usepackage{tabularx}
\usepackage{diagbox} % case d'en-tête diagonale des tableaux à double entrée
% Colonnes extensibles centrée / gauche / droite (C, L, R), souvent utilisées par les modèles
\newcolumntype{C}{>{\centering\arraybackslash}X}
\newcolumntype{L}{>{\raggedright\arraybackslash}X}
\newcolumntype{R}{>{\raggedleft\arraybackslash}X}
\usepackage{array}
\usepackage{multicol}
\usepackage{multirow}
\usepackage{siunitx}
% parse-numbers=false : accepte aussi \SI{2,5 \times 10^{-3}}{...}
\sisetup{locale=FR,output-decimal-marker={,},group-minimum-digits=4,parse-numbers=false}
\DeclareSIUnit\ohm{\ensuremath{\symup{\Omega}}} % U+2126 absent de la police
\DeclareSIUnit\division{div} % réglages d'oscilloscope (ms/div, V/div)
\DeclareSIUnit\tr{tr} % tours (vitesse de rotation, ex. \SI{1500}{\tr\per\minute})
\DeclareSIUnit\molar{M} % concentration molaire (ex. \SI{3}{\molar}, \SI{1}{\micro\molar})
\DeclareSIUnit\an{an} % année (le modèle confond parfois avec \year, un compteur TeX) — ex. \SI{5}{\kilo\meter\per\an}
\DeclareSIUnit\FCFA{FCFA} % franc CFA (ex. \SI{500}{\FCFA\per\kilo\gram})
\DeclareSIUnit\degreeBrix{\textdegree Brix} % teneur en sucre d'un sirop/jus (réfractométrie)
\DeclareSIUnit\month{mois} % le modèle utilise parfois \month, un compteur TeX, comme unité de durée
\usepackage{qrcode}
% \degree : commande fréquemment utilisée par les modèles pour un angle ou une
% température en dehors de \SI{}{} (non fournie par les paquets chargés).
\providecommand{\degree}{^{\circ}}
% \qed (fin de démonstration), utilisé par les modèles sans amsthm
\providecommand{\qed}{\hfill$\square$}
% \barre : double barre des valeurs interdites dans un tableau de variations (tkz-tab)
\providecommand{\barre}{\ensuremath{\|}}
% environnement proof (amsthm non chargé) utilisé par les modèles
\ifdefined\proof\else\newenvironment{proof}[1][Démonstration]{\par\noindent\textit{#1.}\ }{\qed\par}\fi
\usepackage{lastpage}
\usepackage{hyperref}
\hypersetup{hidelinks}

% ============================================================
% Palette « Pop » OnBuch+ — mêmes hex que la classe P (plus_ui.dart)
% ============================================================
\definecolor{poporange}{HTML}{F59321}
\definecolor{popdark}{HTML}{E07A0C}
\definecolor{popcream}{HTML}{FFF6E8}
\definecolor{popink}{HTML}{1C1714}
\definecolor{poppurple}{HTML}{7A5AE0}
\definecolor{popgreen}{HTML}{2E9E6A}
\definecolor{poppink}{HTML}{E0457B}
\definecolor{popblue}{HTML}{2F7DE1}
\definecolor{popgold}{HTML}{C98A12}
\definecolor{popmuted}{HTML}{8A7F76}
\definecolor{popline}{HTML}{EADFD2}
\definecolor{popgray}{HTML}{8A7F76}
\colorlet{popgrey}{popgray}
% Alias tolérants (le modèle nomme parfois les couleurs différemment)
\colorlet{popwhite}{white}
\colorlet{popblack}{popink}
\colorlet{popred}{poppink}
\colorlet{popyellow}{popgold}
% Teintes claires (fonds de tableaux/encadrés) — le modèle nomme parfois
% les couleurs avec un suffixe L (ex. popblueL) sans les avoir définies.
\colorlet{poporangeL}{poporange!20}
\colorlet{popdarkL}{popdark!20}
\colorlet{poppurpleL}{poppurple!20}
\colorlet{popgreenL}{popgreen!20}
\colorlet{poppinkL}{poppink!20}
\colorlet{popblueL}{popblue!20}
\colorlet{popgoldL}{popgold!20}
\colorlet{popredL}{popred!20}
\colorlet{popgrayL}{popgray!20}
% Teintes claires pour fonds de tableaux / zones
\colorlet{popblueL}{popblue!10!white}
\colorlet{poporangeL}{poporange!14!white}
\colorlet{popgreenL}{popgreen!12!white}
\colorlet{poppurpleL}{poppurple!12!white}
\colorlet{poppinkL}{poppink!12!white}
\colorlet{popgoldL}{popgold!14!white}

\pagecolor{popcream}

% ============================================================
% Typographie
% ============================================================
\defaultfontfeatures{Ligatures=TeX}
\setmainfont{PlusJakartaSans-Medium.ttf}[Path=../../../fonts/,BoldFont=PlusJakartaSans-Bold.ttf]
% Symboles, API et emojis absents de la police principale : repli sur DejaVu Sans (caractères actifs)
\newfontface\symfont{DejaVuSans.ttf}[Path=../../../fonts/]
\makeatletter
\newcommand\symactive[1]{\catcode#1=13 \begingroup\lccode`\~=#1 \lowercase{\endgroup\def~}{{\symfont\char#1}}}
\newcommand\symblank[1]{\catcode#1=13 \begingroup\lccode`\~=#1 \lowercase{\endgroup\def~}{}}
\newcommand\symhyph[1]{\catcode#1=13 \begingroup\lccode`\~=#1 \lowercase{\endgroup\def~}{\mbox{-}}}
\newcount\symc
\newcommand\symrange[2]{\symc=#1 \@whilenum\symc<#2 \do{\expandafter\symactive\expandafter{\the\symc}\advance\symc by 1 }}
\newcommand\symblankrange[2]{\symc=#1 \@whilenum\symc<#2 \do{\expandafter\symblank\expandafter{\the\symc}\advance\symc by 1 }}
\makeatother
\symrange{"0250}{"02FF}\symactive{"2713}\symactive{"2714}\symactive{"2717}\symactive{"2718}\symactive{"2610}\symactive{"2611}\symactive{"2612}
\symrange{"2600}{"27BF}
\catcode"2705=13 \begingroup\lccode`\~="2705 \lowercase{\endgroup\def~}{{\symfont\char"2713}}\symblank{"26BD}\symblank{"26A1}
\symrange{"2460}{"257F}\symactive{"223C}\symactive{"2248}\symactive{"2260}
\symactive{"01F8}\symactive{"01F9}\symactive{"1E3E}\symactive{"1E3F}
\symhyph{"2011}\symhyph{"2E1A}\symblankrange{"1F300}{"1FAFF}\symblank{"FE0F}
% Caractères chinois : WenQuanYi Zen Hei (sous-ensemble GB2312 + pinyin), gras/italique simulés
\setCJKmainfont{WenQuanYiZenHei-subset.ttf}[Path=../../../fonts/,AutoFakeBold=3,AutoFakeSlant=0.2]
\xeCJKsetup{CJKspace=false}
% Pinyin : voyelles à caron (ǎ ǐ ǒ ǔ ǖ ǘ ǚ ǜ) absentes de la police principale -> caractères actifs
\newfontface\pyfont{WenQuanYiZenHei-subset.ttf}[Path=../../../fonts/]
\newcommand\pyactive[1]{\catcode#1=13 \begingroup\lccode`\~=#1 \lowercase{\endgroup\def~}{{\pyfont\char#1}}}
\pyactive{"01CD}\pyactive{"01CE}\pyactive{"01CF}\pyactive{"01D0}\pyactive{"01D1}\pyactive{"01D2}\pyactive{"01D3}\pyactive{"01D4}\pyactive{"01D5}\pyactive{"01D6}\pyactive{"01D7}\pyactive{"01D8}\pyactive{"01D9}\pyactive{"01DA}\pyactive{"01DB}\pyactive{"01DC}
% Maths en sans-serif (Fira Math) avec chiffres et lettres droites de Plus
% Jakarta, pour que les formules se fondent dans le texte.
\usepackage[math-style=ISO,bold-style=ISO]{unicode-math}
\setmathfont{FiraMath-Regular.otf}[Path=../../../fonts/]
\setmathfont{PlusJakartaSans-Medium.ttf}[Path=../../../fonts/,range={up/num,up/{Latin,latin},\mathup}]
\usepackage{icomma} % virgule décimale sans espace en mode math
% Caractères mathématiques Unicode absents des polices de texte (Plus Jakarta,
% Archivo Black) : rendus par leur équivalent mathématique, en texte comme en
% formule — sinon ils s'affichent en carré vide (ex. « Arithmétique dans ℤ »).
\usepackage{newunicodechar}
\newunicodechar{ℝ}{\ensuremath{\mathbb{R}}}
\newunicodechar{ℤ}{\ensuremath{\mathbb{Z}}}
\newunicodechar{ℂ}{\ensuremath{\mathbb{C}}}
\newunicodechar{ℕ}{\ensuremath{\mathbb{N}}}
\newunicodechar{ℚ}{\ensuremath{\mathbb{Q}}}
\newunicodechar{𝔻}{\ensuremath{\mathbb{D}}}
\newunicodechar{α}{\ensuremath{\alpha}}
\newunicodechar{β}{\ensuremath{\beta}}
\newunicodechar{γ}{\ensuremath{\gamma}}
\newunicodechar{δ}{\ensuremath{\delta}}
\newunicodechar{ε}{\ensuremath{\varepsilon}}
\newunicodechar{θ}{\ensuremath{\theta}}
\newunicodechar{λ}{\ensuremath{\lambda}}
\newunicodechar{μ}{\ensuremath{\mu}}
\newunicodechar{σ}{\ensuremath{\sigma}}
\newunicodechar{τ}{\ensuremath{\tau}}
\newunicodechar{φ}{\ensuremath{\varphi}}
\newunicodechar{ψ}{\ensuremath{\psi}}
\newunicodechar{ω}{\ensuremath{\omega}}
\newunicodechar{Σ}{\ensuremath{\Sigma}}
% majuscules produites par \MakeUppercase dans les titres (α-aminés -> Α)
\newunicodechar{Α}{\ensuremath{\alpha}}
\newunicodechar{Β}{\ensuremath{\beta}}
\newunicodechar{Γ}{\ensuremath{\Gamma}}
\newunicodechar{Δ}{\ensuremath{\Delta}}
\newunicodechar{Ε}{\ensuremath{\varepsilon}}
\newunicodechar{Θ}{\ensuremath{\Theta}}
\newunicodechar{Λ}{\ensuremath{\Lambda}}
\newunicodechar{Μ}{\ensuremath{\mu}}
\newunicodechar{Τ}{\ensuremath{\tau}}
\newunicodechar{Φ}{\ensuremath{\Phi}}
\newunicodechar{Ψ}{\ensuremath{\Psi}}
\newunicodechar{Ω}{\ensuremath{\Omega}}
\newunicodechar{↦}{\ensuremath{\mapsto}}
\newunicodechar{∈}{\ensuremath{\in}}
\newunicodechar{∉}{\ensuremath{\notin}}
\newunicodechar{∧}{\ensuremath{\wedge}}
\newunicodechar{⇔}{\ensuremath{\Leftrightarrow}}
\newunicodechar{⟺}{\ensuremath{\Longleftrightarrow}}
\newunicodechar{⇒}{\ensuremath{\Rightarrow}}
\newunicodechar{⟹}{\ensuremath{\Longrightarrow}}
\newunicodechar{∘}{\ensuremath{\circ}}
\newunicodechar{∪}{\ensuremath{\cup}}
\newunicodechar{∩}{\ensuremath{\cap}}
\newunicodechar{≡}{\ensuremath{\equiv}}
\newunicodechar{⊂}{\ensuremath{\subset}}
\newunicodechar{⊥}{\ensuremath{\perp}}
\newunicodechar{∃}{\ensuremath{\exists}}
\newunicodechar{∀}{\ensuremath{\forall}}
\newunicodechar{∛}{\ensuremath{\sqrt[3]{\,}}}
\newunicodechar{∜}{\ensuremath{\sqrt[4]{\,}}}
\newunicodechar{⃗}{\ensuremath{{}^{\to}}}
\newunicodechar{ⁿ}{\ifmmode^{n}\else\textsuperscript{\ensuremath{n}}\fi}
\newunicodechar{ˣ}{\ifmmode^{x}\else\textsuperscript{\ensuremath{x}}\fi}
\newunicodechar{ᵇ}{\ifmmode^{b}\else\textsuperscript{\ensuremath{b}}\fi}
\newunicodechar{ʳ}{\ifmmode^{r}\else\textsuperscript{\ensuremath{r}}\fi}
\newunicodechar{ᵘ}{\ifmmode^{u}\else\textsuperscript{\ensuremath{u}}\fi}
\newunicodechar{ʸ}{\ifmmode^{y}\else\textsuperscript{\ensuremath{y}}\fi}
\newunicodechar{ᵃ}{\ifmmode^{a}\else\textsuperscript{\ensuremath{a}}\fi}
\newunicodechar{ᵉ}{\ifmmode^{e}\else\textsuperscript{\ensuremath{e}}\fi}
\newunicodechar{ᵏ}{\ifmmode^{k}\else\textsuperscript{\ensuremath{k}}\fi}
\newunicodechar{ᵖ}{\ifmmode^{p}\else\textsuperscript{\ensuremath{p}}\fi}
\newunicodechar{ᴺ}{\ifmmode^{N}\else\textsuperscript{\ensuremath{N}}\fi}
\newunicodechar{ᶜ}{\ifmmode^{c}\else\textsuperscript{\ensuremath{c}}\fi}
\newunicodechar{ʰ}{\ifmmode^{h}\else\textsuperscript{\ensuremath{h}}\fi}
\newunicodechar{ᵠ}{\ifmmode^{q}\else\textsuperscript{\ensuremath{q}}\fi}
\newunicodechar{ₙ}{\ifmmode_{n}\else\textsubscript{\ensuremath{n}}\fi}
\newunicodechar{ₐ}{\ifmmode_{a}\else\textsubscript{\ensuremath{a}}\fi}
\newunicodechar{ᵢ}{\ifmmode_{i}\else\textsubscript{\ensuremath{i}}\fi}
\newunicodechar{ᵣ}{\ifmmode_{r}\else\textsubscript{\ensuremath{r}}\fi}
\newunicodechar{ₖ}{\ifmmode_{k}\else\textsubscript{\ensuremath{k}}\fi}
\newunicodechar{ᵦ}{\ifmmode_{\beta}\else\textsubscript{\ensuremath{\beta}}\fi}
\newunicodechar{ₓ}{\ifmmode_{x}\else\textsubscript{\ensuremath{x}}\fi}
\newfontfamily\popdisplay{ArchivoBlack-Regular.ttf}[Path=../../../fonts/]
\newfontfamily\poplogo{SpaceGrotesk-Bold.ttf}[Path=../../../fonts/]
\newfontfamily\popsemi{PlusJakartaSans-SemiBold.ttf}[Path=../../../fonts/]
\newfontfamily\popxbold{PlusJakartaSans-ExtraBold.ttf}[Path=../../../fonts/]
\newfontfamily\popxboldls{PlusJakartaSans-ExtraBold.ttf}[Path=../../../fonts/,LetterSpace=3.0]

\setlist[enumerate]{leftmargin=1.7em,itemsep=5pt,topsep=4pt}
\setlist[itemize]{leftmargin=1.7em,itemsep=4pt,topsep=4pt,label={\color{poporange}\textbullet}}
\setlength{\parindent}{0pt}
\setlength{\parskip}{5pt}
\renewcommand{\arraystretch}{1.25}

% pgfplots : style Pop par défaut
\pgfplotsset{
  every axis/.append style={axis line style={popink,line width=1pt},tick style={popink},
    grid style={popline,line width=0.5pt},label style={font=\small},tick label style={font=\footnotesize}},
  popcurve/.style={poporange,line width=1.8pt,no markers,smooth},
}
% Même style disponible dans un \draw TikZ simple (hors pgfplots)
\tikzset{popcurve/.style={poporange,line width=1.8pt}}
\tikzset{popgrid/.style={popline,very thin}}
\colorlet{popcurve}{poporange} % le nom du style est parfois utilisé comme couleur

% ============================================================
% Pill / squiggle
% ============================================================
\newcommand{\poppill}[3]{%
  \tikz[baseline=-0.55ex]\node[draw=popink,line width=1.2pt,rounded corners=2.6pt,%
    fill=#2,inner xsep=6.5pt,inner ysep=3pt]{%
    \popxboldls\fontsize{7}{7}\selectfont\color{#3}\MakeUppercase{#1}};%
}
\newcommand{\popsquiggle}[1]{%
  \tikz[baseline=-0.4ex]{
    \draw[#1,line width=2.6pt,line cap=round]
      plot [smooth] coordinates {(0,0.09) (0.34,0.01) (0.68,0.21) (1.02,0.01) (1.36,0.10)};
  }%
}
% Mot-clé surligné (marqueur orange sous le texte)
\newcommand{\cle}[1]{{\popxbold\color{popdark}#1}}

% ============================================================
% En-tête / pied
% ============================================================
\pagestyle{fancy}
\fancyhf{}
\renewcommand{\headrulewidth}{0pt}
\renewcommand{\footrulewidth}{0pt}
\lhead{\raisebox{-0.15ex}{\poplogo\fontsize{13}{13}\selectfont\color{popink}OnBuch\color{poporange}+}}
\rhead{\poppill{\DOCMATIERE\ \textbullet\ \DOCNIVEAU\ \textbullet\ Leçon \DOCLECON}{white}{popink}}
\lfoot{\popxbold\fontsize{7.6}{7.6}\selectfont\color{popmuted}\MakeUppercase{Cours signé OnBuch+}\ \textbullet\ {\popsemi\DOCTITRE}}
\rfoot{\tikz[baseline=-0.6ex]\node[rounded corners=8.5pt,fill=popink,minimum height=17pt,minimum width=17pt,inner xsep=5pt,inner sep=0]{\popxbold\fontsize{8}{8}\selectfont\color{popcream}\thepage\,/\,\pageref*{LastPage}};}

% ============================================================
% Titres
% ============================================================
\newcommand{\chaptitre}[2]{%
  \begin{tcolorbox}[enhanced,colback=poporange,colframe=popink,boxrule=2.6pt,arc=11pt,
      drop shadow=popink,left=15pt,right=15pt,top=12pt,bottom=11pt,before skip=3pt,after skip=18pt]
    \poppill{#2}{popink}{popcream}\par\vspace{9pt}
    {\raggedright\hyphenpenalty=10000\exhyphenpenalty=10000\popdisplay\fontsize{22}{24}\selectfont\color{popcream}\MakeUppercase{#1}\par}
  \end{tcolorbox}
}
% Section numérotée : pastille orange avec le numéro + titre Archivo
\newcounter{popsec}
\newcommand{\coursec}[1]{%
  \stepcounter{popsec}\setcounter{popsub}{0}%
  \par\vspace{16pt}\noindent
  \begin{minipage}{\linewidth}
  \tikz[baseline=-0.7ex]\node[draw=popink,line width=1.6pt,fill=poporange,rounded corners=4pt,minimum size=22pt,inner sep=0]{\popdisplay\fontsize{12}{12}\selectfont\color{popcream}\thepopsec};%
  \hspace{8pt}{\popdisplay\fontsize{15}{17}\selectfont\color{popink}\MakeUppercase{#1}}\par
  \vspace{4pt}\noindent{\color{poporange}\rule{36pt}{3.6pt}}
  \end{minipage}\par\nopagebreak\vspace{8pt}
}
\newcounter{popsub}
\newcommand{\courssub}[1]{%
  \stepcounter{popsub}%
  \par\vspace{9pt}\noindent{\popxbold\fontsize{11.5}{13}\selectfont\color{popink}{\color{poporange}\thepopsec.\thepopsub}\hspace{6pt}#1}\par\nopagebreak\vspace{4pt}
}

% ============================================================
% Boîtes pédagogiques — même recette : fond blanc, bordure épaisse colorée,
% ombre dure encre, badge plein. Titre optionnel [#1].
% ============================================================
\tcbset{popbase/.style={breakable,enhanced,colback=white,boxrule=2.3pt,arc=9pt,
  shadow={3pt}{-3pt}{0pt}{popink},left=14pt,right=14pt,top=11pt,bottom=11pt,before skip=11pt,after skip=15pt,
  fontupper=\popsemi\fontsize{10.6}{14.5}\selectfont\color{popink},
  fontlower=\popsemi\fontsize{10.6}{14.5}\selectfont\color{popink}}}
% Coche dessinée (le glyphe ✓ n'existe pas dans les polices du document)
\newcommand{\popcheck}[1]{\tikz[baseline=-0.5ex]\draw[#1,line width=1.6pt,line cap=round,line join=round](-0.13,0.0)--(-0.04,-0.09)--(0.13,0.1);}
\providecommand{\coche}{\popcheck{popgreen}}
\providecommand{\resultat}[1]{\par\smallskip\noindent\fcolorbox{popgreen}{popgreenL}{\parbox{\dimexpr\linewidth-2\fboxsep-2\fboxrule\relax}{\textbf{Résultat.} #1}}\par\smallskip}
\newcommand{\popboxhead}[3]{\poppill{#1}{#2}{white}\ifx\relax#3\relax\else\hspace{6pt}{\popxbold\fontsize{10.6}{12}\selectfont\color{popink}#3}\fi\par\vspace{7pt}}

\newtcolorbox{aretenir}[1][]{popbase,colframe=popgreen,before upper={\popboxhead{À retenir}{popgreen}{#1}}}
% Texte en chinois (document de lecture, dialogue, extrait) : cadre
% d'étude. Un titre optionnel (source, auteur) s'affiche dans l'en-tête.
\newtcolorbox{textebox}[1][]{popbase,colframe=popblue,before upper={\popboxhead{文本 · Texte}{popblue}{#1}},after upper={}}
% Marques ✓ / ✗ (forme correcte / fautive) souvent utilisées par les modèles
\providecommand{\cross}{\textcolor{poppink}{\ensuremath{\times}}}
\providecommand{\xmark}{\cross}\providecommand{\crossmark}{\cross}\providecommand{\cmark}{\ensuremath{\checkmark}}
% Ligne de réponse / justification à compléter (inventée par les modèles)
\providecommand{\justification}[1]{\par\noindent\textit{Justification~:} \dotfill\par}
% \textipa (package tipa non chargé) : la police ne couvre pas l'API -> simple passage
\providecommand{\textipa}[1]{#1}
% Soulignement (ulem non chargé) : \uline, \ul
\providecommand{\uline}[1]{\underline{#1}}\providecommand{\ul}[1]{\underline{#1}}
% \glossaire{\item ...} inventée par les modèles : liste de glossaire
\providecommand{\glossaire}[1]{\begin{itemize}#1\end{itemize}}
% \alert (beamer) inventée par les modèles : texte mis en évidence
\providecommand{\alert}[1]{\textbf{\textcolor{poppink}{#1}}}
% variantes ASCII d'ordinaux écrites en macro
\providecommand{\iere}{\textsuperscript{re}}\providecommand{\ieme}{\textsuperscript{e}}\providecommand{\ere}{\textsuperscript{re}}\providecommand{\eme}{\textsuperscript{e}}\providecommand{\er}{\textsuperscript{er}}\providecommand{\ie}{\textsuperscript{e}}
% Ordinaux français écrits en macro par les modèles : 1\ière, 3\ième
\newcommand{\ière}{\textsuperscript{re}}\newcommand{\ième}{\textsuperscript{e}}
% \exemple (inventée par les modèles) : étiquette « Exemple : »
\providecommand{\exemple}{\textbf{Exemple~:}\ }
% \square utilisable aussi hors mode math (cases à cocher)
\AtBeginDocument{\let\squareMath\square\renewcommand{\square}{\ensuremath{\squareMath}}}
% Mot ou phrase en chinois à l'intérieur d'un paragraphe français
\newcommand{\zh}[1]{\ifmmode\text{#1}\else{#1}\fi}
\let\eng\zh \let\esp\zh \let\all\zh % alias tolérés
% flèches d'intonation inventées par les modèles
\providecommand{\textsearrow}{}\renewcommand{\textsearrow}{\ensuremath{\searrow}}\providecommand{\textnearrow}{}\renewcommand{\textnearrow}{\ensuremath{\nearrow}}
% \fr{..} : mot français (inventé par les modèles)
\providecommand{\fr}[1]{\textit{#1}}
% zhbox : encadré de texte chinois inventé par les modèles
\newenvironment{zhbox}{\begin{quote}}{\end{quote}}
% \courssubsub : titre de niveau 3 inventé par les modèles
\providecommand{\courssubsub}[1]{\par\medskip\noindent\textbf{#1}\par\smallskip}
% \vs : « versus » inventé par les modèles
\providecommand{\vs}{\textit{vs}\ }
% \blank : case à compléter inventée par les modèles
\providecommand{\blank}[1][]{\underline{\hspace{1.6em}}}
\newtcolorbox{exemplebox}[1][]{popbase,colframe=poporange,before upper={\popboxhead{Exemple}{poporange}{#1}}}
\newtcolorbox{definition}[1][]{popbase,colframe=popblue,before upper={\popboxhead{Définition}{popblue}{#1}}}
\newtcolorbox{propriete}[1][]{popbase,colframe=poppurple,before upper={\popboxhead{Propriété}{poppurple}{#1}}}
\newtcolorbox{methode}[1][]{popbase,colframe=popgold,before upper={\popboxhead{Méthode}{popgold}{#1}}}
\newtcolorbox{attention}[1][]{popbase,colframe=poppink,before upper={\popboxhead{Attention}{poppink}{#1}}}
\newtcolorbox{experience}[1][]{popbase,colframe=popblue,colback=popblueL,before upper={\popboxhead{Expérience}{popblue}{#1}}}
% « Le savais-tu ? » : carte violette pleine (contexte, culture, Cameroun)
\newtcolorbox{savaistu}[1][]{popbase,colback=poppurple,colframe=popink,
  fontupper=\popsemi\fontsize{10.4}{14.2}\selectfont\color{white},
  before upper={\poppill{Le savais-tu\,?}{white}{poppurple}\ifx\relax#1\relax\else\hspace{6pt}{\popxbold\color{white}#1}\fi\par\vspace{7pt}}}
% Exercice résolu : énoncé en haut, \tcblower puis la solution sur fond crème
\newcounter{popexo}\newcounter{popexr}
\newtcolorbox{exoresolu}[1][]{popbase,colframe=poporange,colbacklower=popcream,
  segmentation style={popink,line width=1.2pt,dashed},
  before upper={\stepcounter{popexr}\popboxhead{Exercice résolu \thepopexr}{poporange}{#1}},
  before lower={\poppill{Solution}{popink}{popcream}\par\vspace{6pt}}}
% Exercice (non résolu en place) — la correction est dans la partie Corrigés
\newtcolorbox{exercice}[1][]{popbase,colframe=popink,
  before upper={\stepcounter{popexo}\popboxhead{Exercice \thepopexo}{popink}{#1}}}
% Corrigé
\newtcolorbox{corrige}[1][]{popbase,colframe=popgreen,colback=popgreenL,
  before upper={\popboxhead{Corrigé}{popgreen}{#1}}}
% Objectifs / prérequis / bilan
\newtcolorbox{objectifs}{popbase,colframe=popink,colback=poporangeL,
  before upper={\popboxhead{Objectifs de la leçon}{popink}{}}}
\newtcolorbox{prerequis}{popbase,colframe=popink,colback=popblueL,
  before upper={\popboxhead{Prérequis}{popblue}{}}}
\newtcolorbox{bilan}{popbase,colframe=popink,colback=white,boxrule=2.8pt,
  before upper={\popboxhead{Fiche bilan}{poporange}{L'essentiel à connaître pour l'examen}}}
% Figure encadrée avec légende
\newtcolorbox{popfigure}[1][]{enhanced,breakable=false,colback=white,colframe=popline,boxrule=1.4pt,arc=8pt,
  left=10pt,right=10pt,top=10pt,bottom=8pt,before skip=10pt,after skip=12pt,halign=center,
  lower separated=false,halign lower=center,
  fontlower=\popsemi\fontsize{9}{11.5}\selectfont\color{popmuted},
  #1}
\newcommand{\legende}[1]{\par\vspace{6pt}{\popsemi\fontsize{9}{11.5}\selectfont\color{popmuted}#1}}

% ============================================================
% Signature OnBuch+
% ============================================================
\newcommand{\poplogoinline}[1]{{\poplogo\fontsize{#1}{#1}\selectfont\color{popink}OnBuch\color{poporange}+}}
\newcommand{\signatureonbuch}{%
  \begin{tcolorbox}[enhanced,colback=popink,colframe=popink,boxrule=2.2pt,arc=10pt,
      drop shadow=poporange,left=14pt,right=14pt,top=10pt,bottom=10pt,before skip=0pt,after skip=0pt]
    \tikz[baseline=-0.7ex]\node[circle,fill=popgreen,draw=popcream,line width=1.3pt,minimum size=22pt,inner sep=0]{\popcheck{white}};%
    \hspace{9pt}\begin{minipage}[c]{0.62\linewidth}
      {\popxbold\fontsize{12}{13}\selectfont\color{popcream}Signé\ \textbullet\ {\poplogo OnBuch\color{poporange}+}}\par\vspace{2pt}
      {\raggedright\popsemi\fontsize{8}{10.5}\selectfont\color{popline}\MakeUppercase{Équipe pédagogique OnBuch+}\\\MakeUppercase{Conforme au programme officiel MINESEC}\par}
    \end{minipage}\hfill
    {\poplogo\fontsize{22}{22}\selectfont\color{popcream}OnBuch\color{poporange}+}
  \end{tcolorbox}%
}

% ============================================================
% Page de garde
% #1 titre · #2 matière · #3 niveau · #4 module · #5 n° leçon · #6 durée ·
% #7 description / objectifs (texte ou itemize)
% ============================================================
\newcommand{\pagegarde}[7]{%
  \thispagestyle{empty}
  \newgeometry{a4paper,margin=1.7cm,top=1.6cm,bottom=1.4cm}%
  \begin{tikzpicture}[remember picture,overlay]
    \fill[poporange!18] ([xshift=-3.2cm,yshift=-3.6cm]current page.north east) circle (4.6cm);
    \fill[poppurple!14] ([xshift=2.4cm,yshift=7.6cm]current page.south west) circle (3.8cm);
  \end{tikzpicture}%
  {\poplogo\fontsize{30}{30}\selectfont\color{popink}OnBuch\color{poporange}+}\hfill
  \raisebox{0.6ex}{\poppill{Cours signé \textbullet\ Premium}{popink}{popcream}}\par
  \vspace{1pt}\popsquiggle{poppurple}\par
  \vspace{8pt}
  \begin{tcolorbox}[enhanced,colback=poporange,colframe=popink,boxrule=2.8pt,arc=13pt,
      drop shadow=popink,left=18pt,right=18pt,top=12pt,bottom=13pt,after skip=10pt]
    \poppill{#2 \textbullet\ #3}{popink}{popcream}\hspace{5pt}\poppill{#4}{white}{popink}\par\vspace{8pt}
    {\popdisplay\fontsize{14}{14}\selectfont\color{popink}LEÇON #5}\par\vspace{4pt}
    {\raggedright\hyphenpenalty=10000\exhyphenpenalty=10000\popdisplay\fontsize{25}{27}\selectfont\color{popcream}\MakeUppercase{#1}\par}
    \vspace{8pt}
    {\popxbold\fontsize{9.5}{9.5}\selectfont\color{popink}\MakeUppercase{Durée conseillée : #6}}
  \end{tcolorbox}
  \begin{tcolorbox}[enhanced,colback=white,colframe=popink,boxrule=2.2pt,arc=10pt,
      drop shadow=popink,left=16pt,right=16pt,top=10pt,bottom=10pt,after skip=8pt]
    \poppill{Dans cette leçon}{poppurple}{white}\par\vspace{5pt}
    \popsemi\fontsize{9.3}{12.6}\selectfont\color{popink}#7
  \end{tcolorbox}
  \noindent\begin{minipage}[t]{0.487\linewidth}
    \begin{tcolorbox}[enhanced,colback=popgreen,colframe=popink,boxrule=2.2pt,arc=10pt,
        drop shadow=popink,left=11pt,right=11pt,top=10pt,bottom=10pt,height=3.1cm,valign=center]
      \tcbox[colback=white,colframe=popink,boxrule=1.3pt,arc=5pt,boxsep=0pt,nobeforeafter,
        left=3pt,right=3pt,top=3pt,bottom=3pt,valign=center]{\qrcode[height=1.9cm]{https://play.google.com/store/apps/details?id=com.luvvix.onbuch&hl=fr}}%
      \hspace{8pt}\begin{minipage}[c]{0.5\linewidth}
        {\popdisplay\fontsize{11}{12.5}\selectfont\color{white}\MakeUppercase{Télécharge OnBuch}}\par\vspace{4pt}
        {\raggedright\popsemi\fontsize{8.4}{11}\selectfont\color{white}Gratuit \textbullet\ Google Play\\Tuteur IA, annales, concours, résultats\par}
      \end{minipage}
    \end{tcolorbox}
  \end{minipage}\hfill
  \begin{minipage}[t]{0.487\linewidth}
    \begin{tcolorbox}[enhanced,colback=poppurple,colframe=popink,boxrule=2.2pt,arc=10pt,
        drop shadow=popink,left=11pt,right=11pt,top=10pt,bottom=10pt,height=3.1cm,valign=center]
      \tikz[baseline=-0.7ex]\node[circle,fill=white,draw=popink,line width=1.3pt,minimum size=1.6cm,inner sep=0]{\popdisplay\fontsize{24}{24}\selectfont\color{poppurple}+};%
      \hspace{8pt}\begin{minipage}[c]{0.6\linewidth}
        {\popdisplay\fontsize{11}{12.5}\selectfont\color{white}\MakeUppercase{Passe à OnBuch+}}\par\vspace{4pt}
        {\raggedright\popsemi\fontsize{8.4}{11}\selectfont\color{white}Tous les cours signés, Tuteur IA illimité, fiches PDF — onglet OnBuch+ de l'app\par}
      \end{minipage}
    \end{tcolorbox}
  \end{minipage}
  \vspace{8pt}
  \popcontenu
  \vfill
  \signatureonbuch
  \restoregeometry
  \newpage
}

% Rangée « Ce cours contient » (page de garde)
\newcommand{\poptile}[3]{% #1 couleur #2 gros texte #3 légende
  \begin{tcolorbox}[enhanced,colback=#1,colframe=popink,boxrule=2pt,arc=9pt,shadow={2.5pt}{-2.5pt}{0pt}{popink},
      left=6pt,right=6pt,top=9pt,bottom=9pt,halign=center,valign=center,height=1.75cm,width=\linewidth,nobeforeafter]
    {\popdisplay\fontsize{12.5}{13}\selectfont\color{white}\MakeUppercase{#2}}\par\vspace{3pt}
    {\centering\popsemi\fontsize{7.8}{9.5}\selectfont\color{white}#3\par}
  \end{tcolorbox}}
\newcommand{\popcontenu}{%
  \par\noindent{\popxboldls\fontsize{7.5}{7.5}\selectfont\color{popmuted}CE COURS SIGNÉ CONTIENT}\par\vspace{7pt}
  \noindent\begin{minipage}[t]{0.235\linewidth}\poptile{popblue}{Cours}{complet, illustré, pas à pas}\end{minipage}\hfill
  \begin{minipage}[t]{0.235\linewidth}\poptile{poporange}{Méthodes}{et exercices résolus}\end{minipage}\hfill
  \begin{minipage}[t]{0.235\linewidth}\poptile{poppink}{Exercices}{type examen + corrigés}\end{minipage}\hfill
  \begin{minipage}[t]{0.235\linewidth}\poptile{popgreen}{Fiche bilan}{l'essentiel à retenir}\end{minipage}\par}

% Sommaire
\newtcolorbox{sommaire}{enhanced,colback=white,colframe=popink,boxrule=2.2pt,arc=10pt,
  drop shadow=popink,left=18pt,right=18pt,top=13pt,bottom=13pt,before skip=6pt,after skip=20pt,
  before upper={\poppill{Sommaire}{poppurple}{white}\par\vspace{9pt}\popsemi\fontsize{10.6}{14.5}\selectfont\color{popink}}}

% Signature de fin de cours — poussée tout en bas de la dernière page
\newcommand{\findecours}{%
  \par\vfill\nopagebreak
  \begin{center}\popsquiggle{poporange}\end{center}\vspace{4pt}
  \signatureonbuch
}
\AtBeginDocument{\def\llbracket{\mathopen{[\mkern-3.5mu[}}\def\rrbracket{\mathclose{]\mkern-3.5mu]}}} % intervalles d'entiers (glyphe ⟦ absent de la police)
% Glyphes absents de Fira Math : redéfinis à partir de glyphes disponibles
\AtBeginDocument{%
  \def\setminus{\mathbin{\backslash}}\let\smallsetminus\setminus
  \def\vdots{\mathord{\vbox{\baselineskip3.2pt\lineskiplimit0pt\kern4pt\hbox{.}\hbox{.}\hbox{.}}}}%
  \def\ddots{\mathinner{\mkern1mu\raise6.5pt\hbox{.}\mkern2mu\raise3.6pt\hbox{.}\mkern2mu\raise0.7pt\hbox{.}\mkern1mu}}%
  \def\triangle{\mathord{\scalebox{0.9}{$\symup{\Delta}$}}}%
  \def\nabla{\mathord{\raisebox{\depth}{\scalebox{1}[-1]{$\symup{\Delta}$}}}}%
  \def\checkmark{\popcheck{popdark}}%
  \def\star{\mathbin{\ast}}\def\bigstar{\mathord{\ast}}%
  \def\diamond{\mathbin{\scalebox{0.6}{$\Diamond$}}}%
  \def\bigcup{\mathop{\mathchoice{\raisebox{-0.3em}{\scalebox{1.7}{$\cup$}}}{\scalebox{1.25}{$\cup$}}{\cup}{\cup}}\displaylimits}%
  \def\bigcap{\mathop{\mathchoice{\raisebox{-0.3em}{\scalebox{1.7}{$\cap$}}}{\scalebox{1.25}{$\cap$}}{\cap}{\cap}}\displaylimits}%
  \def\lesseqqgtr{\mathrel{\lessgtr}}\def\bigoplus{\mathop{\oplus}\displaylimits}%
}
\providecommand{\ssi}{si et seulement si} % abréviation fréquente
\AtBeginDocument{\providecommand{\wideparen}{\overparen}} % arc de cercle

%% FIGFIX-BEGIN v4 — correctifs globaux des figures (docs/onbuch-generation/tools/figfix)
\usepackage{adjustbox}\usepackage{etoolbox}
% toute figure TikZ/pgfplots plus large que la ligne est réduite à la largeur disponible
\BeforeBeginEnvironment{tikzpicture}{\begin{adjustbox}{max width=\linewidth}}
\AfterEndEnvironment{tikzpicture}{\end{adjustbox}}
% légendes pgfplots lisibles par-dessus les courbes
\pgfplotsset{every axis/.append style={legend style={fill=white,fill opacity=0.92,text opacity=1,draw=black!15,inner sep=3pt}}}
% étiquettes d'arêtes (syntaxe edge["..."]) avec fond blanc
\tikzset{every edge quotes/.append style={fill=white,inner sep=1.2pt}}
%% FIGFIX-END
\begin{document}

\pagegarde{Leçon 22 — Expression écrite : culture chinoise et camerounaise (你喜欢喀麦隆还是中国艺术？为什么？)}{Chinois}{1ère A}{Module 3 : Citoyenneté et ouverture au monde}{ZHA22}{2 h}{Cette leçon d'expression écrite apprend à rédiger un texte simple en chinois comparant l'art camerounais et l'art chinois et justifiant sa préférence. Elle consolide l'écriture correcte des caractères du thème, l'usage des connecteurs logiques et la traduction (thème et version).
\vspace{4pt}
\begin{multicols}{2}\small
\begin{itemize}
\item À la fin de cette leçon, je sais lire et comprendre un texte modèle simple sur les arts du Cameroun et de la Chine.
\item Je sais exprimer par écrit une préférence avec 喜欢...还是... et la justifier avec 因为.
\item Je sais utiliser les connecteurs 首先、然后、还有、所以 pour structurer un texte.
\item Je sais écrire correctement les caractères du thème (艺、术、画、喜、欢、觉、得) en respectant radicaux et ordre des traits.
\item Je sais distinguer les caractères proches (画/面、术/木、欢/吹、觉/党).
\item Je sais traduire des phrases simples du français vers le chinois (thème).
\end{itemize}\end{multicols}}

\begin{sommaire}
\begin{multicols}{2}
\begin{enumerate}
\item Texte modèle et compréhension
\item Lexique du thème : arts et opinions
\item Structure du texte et connecteurs
\item Écriture des caractères du thème
\item Thème et version guidés
\item Production écrite et évaluation
\item Activité d'intégration
\item Méthodes et exercices résolus
\item Exercices
\item Corrigés des exercices
\item Fiche bilan
\end{enumerate}
\end{multicols}
\end{sommaire}

\begin{objectifs}
\begin{itemize}
\item À la fin de cette leçon, je sais lire et comprendre un texte modèle simple sur les arts du Cameroun et de la Chine.
\item Je sais exprimer par écrit une préférence avec 喜欢...还是... et la justifier avec 因为.
\item Je sais utiliser les connecteurs 首先、然后、还有、所以 pour structurer un texte.
\item Je sais écrire correctement les caractères du thème (艺、术、画、喜、欢、觉、得) en respectant radicaux et ordre des traits.
\item Je sais distinguer les caractères proches (画/面、术/木、欢/吹、觉/党).
\item Je sais traduire des phrases simples du français vers le chinois (thème).
\item Je sais traduire un court texte du chinois vers le français (version).
\item Je sais rédiger un texte de 6 à 10 phrases répondant à la question « 你喜欢喀麦隆还是中国艺术？为什么？ ».
\item Je sais m'auto-évaluer à l'aide d'une grille d'évaluation.
\end{itemize}
\end{objectifs}

\begin{prerequis}
\begin{itemize}
\item La structure de préférence 喜欢 et la question alternative 还是 (leçons de Seconde).
\item La justification avec 因为...所以...
\item Le vocabulaire de base des pays et nationalités : 中国、喀麦隆、人.
\item L'ordre général des traits (de haut en bas, de gauche à droite, horizontal avant vertical).
\item Les verbes d'opinion simples : 觉得、认为.
\end{itemize}
\end{prerequis}

\newpage

\coursec{Texte modèle et compréhension}

\courssub{Situation d'accroche et problématique}

Au lycée de Yaoundé, le club d'art organise une grande exposition : des masques bamiléké aux couleurs éclatantes, des sculptures fang élancées, des peintures à l'encre chinoise et de superbes calligraphies. Un correspondant chinois, venu visiter l'exposition, demande à Amina :

\zh{你喜欢喀麦隆还是中国艺术？为什么？}

\emph{Nǐ xǐhuan Kāmàilóng háishi Zhōngguó yìshù? Wèishénme?}

« Tu préfères l'art camerounais ou l'art chinois ? Pourquoi ? »

Amina sait répondre à l'oral, mais saura-t-elle rédiger un petit texte clair, structuré et correct en chinois ? C'est exactement la compétence que cette leçon va construire : \cle{exprimer par écrit une préférence culturelle argumentée}. Pour cela, nous allons d'abord lire et comprendre un texte modèle, puis en dégager la structure et les connecteurs, avant de les réutiliser dans nos propres productions.

\begin{attention}[还是 ou 或者 ?]
Dans la question \zh{你喜欢喀麦隆还是中国艺术？}, le mot \zh{还是} \emph{háishi} signifie « ou » mais \textbf{uniquement dans les questions} : il propose un choix entre deux options. Dans une phrase affirmative, « ou » se dit \zh{或者} \emph{huòzhě}. Compare :\\
\zh{你喜欢茶还是咖啡？} \emph{Nǐ xǐhuan chá háishi kāfēi?} « Tu aimes le thé ou le café ? » (question)\\
\zh{我喜欢茶或者咖啡。} \emph{Wǒ xǐhuan chá huòzhě kāfēi.} « J'aime le thé ou le café. » (affirmation)
\end{attention}

\courssub{Lecture du texte modèle}

\begin{textebox}[我喜欢中国艺术 — Wǒ xǐhuan Zhōngguó yìshù]
喀麦隆和中国都有很有意思的艺术。\\
\emph{Kāmàilóng hé Zhōngguó dōu yǒu hěn yǒu yìsi de yìshù.}\\
« Le Cameroun et la Chine ont tous deux un art très intéressant. »\\[4pt]
喀麦隆有很好看的面具和雕塑，中国有很美的画和书法。\\
\emph{Kāmàilóng yǒu hěn hǎokàn de miànjù hé diāosù, Zhōngguó yǒu hěn měi de huà hé shūfǎ.}\\
« Le Cameroun a de très beaux masques et de belles sculptures, la Chine a de très belles peintures et la calligraphie. »\\[4pt]
我更喜欢中国艺术。\\
\emph{Wǒ gèng xǐhuan Zhōngguó yìshù.}\\
« Je préfère l'art chinois. »\\[4pt]
为什么呢？\\
\emph{Wèishénme ne?}\\
« Pourquoi donc ? »\\[4pt]
首先，中国画的颜色很好看；\\
\emph{Shǒuxiān, Zhōngguó huà de yánsè hěn hǎokàn;}\\
« D'abord, les couleurs de la peinture chinoise sont très belles ; »\\[4pt]
然后，我觉得书法很有意思；\\
\emph{Ránhòu, wǒ juéde shūfǎ hěn yǒu yìsi;}\\
« ensuite, je trouve la calligraphie très intéressante ; »\\[4pt]
还有，中国艺术有几千年的历史。\\
\emph{Hái yǒu, Zhōngguó yìshù yǒu jǐ qiān nián de lìshǐ.}\\
« de plus, l'art chinois a plusieurs milliers d'années d'histoire. »\\[4pt]
所以，我喜欢中国艺术。你呢？\\
\emph{Suǒyǐ, wǒ xǐhuan Zhōngguó yìshù. Nǐ ne?}\\
« C'est pourquoi j'aime l'art chinois. Et toi ? »
\end{textebox}

\begin{definition}[Glossaire du texte]
\begin{center}
\begin{tabularx}{\linewidth}{|X|X|X|X|}
\hline
\rowcolor{popblueL} Caractères & Pinyin & Nature & Traduction \\
\hline
艺术 & yìshù & nom & art \\
\hline
面具 & miànjù & nom & masque \\
\hline
雕塑 & diāosù & nom & sculpture \\
\hline
画 & huà & nom & peinture, tableau \\
\hline
书法 & shūfǎ & nom & calligraphie \\
\hline
更 & gèng & adverbe & davantage, encore plus \\
\hline
颜色 & yánsè & nom & couleur \\
\hline
觉得 & juéde & verbe & trouver, penser, estimer \\
\hline
历史 & lìshǐ & nom & histoire \\
\hline
有意思 & yǒu yìsi & expression & intéressant \\
\hline
好看 & hǎokàn & adjectif & beau (à regarder) \\
\hline
还是 & háishi & conjonction & ou (dans une question) \\
\hline
\end{tabularx}
\end{center}
\end{definition}

\begin{methode}[Comment lire un texte modèle]
\begin{enumerate}
\item Lire une première fois le texte entier pour saisir le \cle{sens global}, sans s'arrêter sur chaque caractère inconnu.
\item Relire en \cle{soulignant les connecteurs} (首先, 然后, 还有, 所以) : ce sont les piliers de la construction.
\item Relever la \cle{structure} du texte : introduction, préférence, arguments, conclusion.
\item \cle{Réutiliser la structure} avec ses propres idées : remplacer le sujet, la préférence et les arguments, en gardant le squelette.
\end{enumerate}
\end{methode}

\courssub{Repérage de la structure du texte}

Observons le texte : il ne se contente pas de dire « j'aime l'art chinois ». Il suit un plan en quatre temps, très facile à imiter.

\begin{itemize}
\item \textbf{Phrase d'introduction} (présentation du sujet, sans opinion) : \zh{喀麦隆和中国都有很有意思的艺术。} \emph{Kāmàilóng hé Zhōngguó dōu yǒu hěn yǒu yìsi de yìshù.} « Le Cameroun et la Chine ont tous deux un art très intéressant. »
\item \textbf{Phrase de préférence} (l'opinion, avec \cle{更} « davantage ») : \zh{我更喜欢中国艺术。} \emph{Wǒ gèng xǐhuan Zhōngguó yìshù.} « Je préfère l'art chinois. »
\item \textbf{Arguments} (annoncés par la question \zh{为什么呢？} puis enchaînés par les connecteurs 首先, 然后, 还有) : les couleurs, la calligraphie, l'histoire millénaire.
\item \textbf{Conclusion} (introduite par \cle{所以} « donc, c'est pourquoi ») : \zh{所以，我喜欢中国艺术。你呢？} « C'est pourquoi j'aime l'art chinois. Et toi ? »
\end{itemize}

\begin{popfigure}
\begin{tikzpicture}[
  bloc/.style={draw=popblue, fill=popblueL, rounded corners=4pt, minimum width=3.4cm, minimum height=1.1cm, align=center, font=\small},
  fleche/.style={-{Stealth[length=3mm]}, thick, popdark}
]
\node[bloc, align=center] (intro) at (0,0){\textbf{Introduction}\\喀麦隆和中国都有…};
\node[bloc, fill=poporangeL, draw=poporange, align=center] (pref) at (4.6,0){\textbf{Préférence}\\我更喜欢…};
\node[bloc, fill=popgreenL, draw=popgreen, align=center] (arg) at (9.2,0){\textbf{Arguments}\\为什么？首先…然后…还有…};
\node[bloc, fill=poppinkL, draw=poppink, align=center] (conc) at (13.8,0){\textbf{Conclusion}\\所以，我喜欢…};
\draw[fleche] (intro) -- (pref);
\draw[fleche] (pref) -- (arg);
\draw[fleche] (arg) -- (conc);
\end{tikzpicture}
\legende{Plan du texte modèle en quatre blocs : introduction, préférence, arguments, conclusion.}
\end{popfigure}

\begin{propriete}[Les connecteurs du texte argumentatif]
Pour enchaîner les arguments en chinois, on utilise des \cle{connecteurs} placés en début de phrase, suivis d'une virgule chinoise « ， » :
\begin{itemize}
\item \zh{首先} \emph{shǒuxiān} : « d'abord, premièrement » — premier argument.
\item \zh{然后} \emph{ránhòu} : « ensuite » — deuxième argument.
\item \zh{还有} \emph{hái yǒu} : « de plus, il y a aussi » — argument supplémentaire.
\item \zh{所以} \emph{suǒyǐ} : « donc, c'est pourquoi » — conclusion logique.
\end{itemize}
Structure générale : \textbf{首先，argument 1；然后，argument 2；还有，argument 3。所以，conclusion。}\\
Remarque : entre deux arguments, on peut utiliser le point-virgule chinois « ； » ; la dernière phrase se termine par le point chinois « 。 ».
\end{propriete}

\begin{exemplebox}[Connecteurs en action]
\zh{首先，中国画的颜色很好看；然后，我觉得书法很有意思。所以，我喜欢中国艺术。}\\
\emph{Shǒuxiān, Zhōngguó huà de yánsè hěn hǎokàn; ránhòu, wǒ juéde shūfǎ hěn yǒu yìsi. Suǒyǐ, wǒ xǐhuan Zhōngguó yìshù.}\\
« D'abord, les couleurs de la peinture chinoise sont très belles ; ensuite, je trouve la calligraphie très intéressante. Donc, j'aime l'art chinois. »\\[4pt]
\zh{首先，喀麦隆的面具很好看；还有，喀麦隆的音乐很有意思。所以，我喜欢喀麦隆艺术。}\\
\emph{Shǒuxiān, Kāmàilóng de miànjù hěn hǎokàn; hái yǒu, Kāmàilóng de yīnyuè hěn yǒu yìsi. Suǒyǐ, wǒ xǐhuan Kāmàilóng yìshù.}\\
« D'abord, les masques camerounais sont très beaux ; de plus, la musique camerounaise est très intéressante. Donc, j'aime l'art camerounais. »
\end{exemplebox}

\begin{attention}[Pièges fréquents des francophones]
\begin{itemize}
\item Ne traduis pas « d'abord » par \zh{第一} dans ce contexte : \zh{第一} \emph{dì yī} numérote (premier, deuxième), tandis que \zh{首先} introduit le premier argument d'un raisonnement.
\item \zh{更} (gèng) signifie « davantage » : \zh{我更喜欢中国艺术} = « je préfère l'art chinois ». Sans 更, la phrase dit simplement « j'aime l'art chinois », sans comparaison.
\item En chinois, le connecteur \zh{所以} peut suivre \zh{因为} dans le même texte, contrairement au français où « parce que... donc » est incorrect. Ici, la question \zh{为什么呢？} joue le rôle du « parce que ».
\item Attention à la ponctuation : après 首先, 然后, 还有, 所以, on met toujours la virgule chinoise « ， ».
\item L'adjectif seul ne suffit pas comme prédicat : on dit \zh{颜色很好看} « les couleurs sont belles », avec l'adverbe \zh{很} devant l'adjectif, et jamais le verbe 是 devant un adjectif.
\end{itemize}
\end{attention}

\courssub{Questions de compréhension}

\begin{exoresolu}[Comprendre le texte modèle]
Énoncé : réponds en chinois aux quatre questions suivantes, d'après le texte modèle.\\
1) 作者更喜欢什么艺术？\emph{(Zuòzhě gèng xǐhuan shénme yìshù? — Quel art l'auteur préfère-t-il ?)}\\
2) 他为什么喜欢中国艺术？\emph{(Tā wèishénme xǐhuan Zhōngguó yìshù? — Pourquoi aime-t-il l'art chinois ?)}\\
3) 喀麦隆有什么艺术？\emph{(Kāmàilóng yǒu shénme yìshù? — Quel art le Cameroun possède-t-il ?)}\\
4) 中国艺术有多少年的历史？\emph{(Zhōngguó yìshù yǒu duōshao nián de lìshǐ? — Combien d'années d'histoire l'art chinois a-t-il ?)}
\tcblower
Solution détaillée :
\begin{enumerate}
\item \zh{作者更喜欢中国艺术。} \emph{Zuòzhě gèng xǐhuan Zhōngguó yìshù.} « L'auteur préfère l'art chinois. » (Repère : phrase de préférence avec 更.)
\item \zh{因为中国画的颜色很好看，书法很有意思，而且中国艺术有几千年的历史。} \emph{Yīnwèi Zhōngguó huà de yánsè hěn hǎokàn, shūfǎ hěn yǒu yìsi, érqiě Zhōngguó yìshù yǒu jǐ qiān nián de lìshǐ.} « Parce que les couleurs de la peinture chinoise sont très belles, la calligraphie est très intéressante, et l'art chinois a plusieurs milliers d'années d'histoire. » (Repère : les trois arguments après 首先, 然后, 还有.)
\item \zh{喀麦隆有很好看的面具和雕塑。} \emph{Kāmàilóng yǒu hěn hǎokàn de miànjù hé diāosù.} « Le Cameroun a de très beaux masques et de belles sculptures. »
\item \zh{中国艺术有几千年的历史。} \emph{Zhōngguó yìshù yǒu jǐ qiān nián de lìshǐ.} « L'art chinois a plusieurs milliers d'années d'histoire. » (\zh{几千} \emph{jǐ qiān} = « plusieurs milliers ».)
\end{enumerate}
\end{exoresolu}

\begin{exercice}[À toi de repérer]
Dans le texte modèle, recopie : a) la phrase d'introduction ; b) la phrase de préférence ; c) les trois arguments avec leur connecteur ; d) la conclusion. Souligne dans chaque cas le mot-clé qui signale la fonction de la phrase (都, 更, 首先, 然后, 还有, 所以).
\end{exercice}

\begin{corrige}[Exercice « À toi de repérer »]
a) Introduction : \zh{喀麦隆和中国都有很有意思的艺术。} (mot-clé : \zh{都}, « tous les deux »).\\
b) Préférence : \zh{我更喜欢中国艺术。} (mot-clé : \zh{更}, « davantage »).\\
c) Arguments : \zh{首先，中国画的颜色很好看；} / \zh{然后，我觉得书法很有意思；} / \zh{还有，中国艺术有几千年的历史。}\\
d) Conclusion : \zh{所以，我喜欢中国艺术。你呢？} (mot-clé : \zh{所以}, « donc »).
\end{corrige}

\begin{savaistu}[Deux arts admirés dans le monde]
La calligraphie (\zh{书法} \emph{shūfǎ}) est considérée en Chine comme le plus noble des arts : bien écrire les caractères au pinceau est un signe de culture et de maîtrise de soi. Au Cameroun, les masques bamiléké et les statues fang sont célèbres dans le monde entier : on les admire dans les plus grands musées. Dire \zh{我更喜欢…} n'est donc pas rejeter un art, mais exprimer un goût personnel argumenté — exactement ce que fait le texte modèle.
\end{savaistu}

\begin{aretenir}[L'essentiel de la section]
Un texte de préférence argumentée en chinois suit quatre étapes : \textbf{introduction} (présenter les deux sujets avec \zh{都有}), \textbf{préférence} (\zh{我更} + verbe), \textbf{arguments} enchaînés par \zh{首先}, \zh{然后}, \zh{还有}, et \textbf{conclusion} avec \zh{所以}. Ce squelette est réutilisable pour n'importe quel sujet : art, musique, cuisine, sport.
\end{aretenir}

\coursec{Lexique du thème : arts et opinions}

Dans la section précédente, nous avons lu et compris le texte modèle \zh{你喜欢喀麦隆还是中国艺术？} (\emph{Nǐ xǐhuan Kāmàilóng háishi Zhōngguó yìshù ?} — « Tu aimes l'art camerounais ou l'art chinois ? »). Nous avons repéré les grandes idées du texte. Pour pouvoir, à votre tour, écrire un texte sur ce thème, il faut maintenant vous construire une \cle{boîte à outils lexicale} : les noms des arts, les verbes pour donner son avis, et les petits mots qui permettent de comparer et de classer. C'est l'objet de cette deuxième section. Observez d'abord, mémorisez ensuite, réutilisez enfin.

\courssub{Le vocabulaire des arts : observation et glossaire}

\textbf{Observation.} Relisez ces phrases tirées du texte modèle et soulignez mentalement les mots qui désignent des arts :

\zh{中国的书法有几千年的历史。} \emph{Zhōngguó de shūfǎ yǒu jǐ qiān nián de lìshǐ.} « La calligraphie chinoise a plusieurs millénaires d'histoire. »

\zh{喀麦隆的面具和雕塑很有名。} \emph{Kāmàilóng de miànjù hé diāosù hěn yǒumíng.} « Les masques et les sculptures du Cameroun sont très célèbres. »

\zh{我喜欢听音乐，也喜欢看画。} \emph{Wǒ xǐhuan tīng yīnyuè, yě xǐhuan kàn huà.} « J'aime écouter de la musique, et j'aime aussi regarder des peintures. »

Vous remarquez que les mots désignant les arts sont souvent composés de \cle{deux caractères} : 艺术 (art), 书法 (calligraphie), 面具 (masque), 雕塑 (sculpture), 音乐 (musique), 舞蹈 (danse). C'est une caractéristique du chinois moderne : la plupart des noms sont \emph{disyllabiques}. Certains mots, en revanche, sont monosyllabiques et servent souvent de radicaux ou de racines dans des mots composés : 画 (peinture / peindre), 舞 (danse / danser — littéraire ou dans des composés).

\begin{definition}[Le mot-clé : 艺术 yìshù]
\zh{艺术} (\emph{yìshù}, n.) signifie « l'art » en général. C'est le mot parapluie de toute cette leçon. On dit : \zh{中国艺术} (l'art chinois), \zh{喀麦隆艺术} (l'art camerounais), \zh{传统艺术} (l'art traditionnel). Le caractère 艺 contient la clé de l'herbe 艹 ; 术 seul signifie « technique, savoir-faire » : l'art est donc étymologiquement un « savoir-faire cultivé ».
\end{definition}

Voici le glossaire complet du thème, à apprendre par cœur :

\begin{center}
\begin{tabularx}{\linewidth}{|X|X|X|X|}
\hline
\rowcolor{popblueL} Caractères & Pinyin & Nature & Traduction \\
\hline
艺术 & yìshù & n. & art \\
\hline
画 & huà & n. / v. & peinture ; peindre, dessiner \\
\hline
书法 & shūfǎ & n. & calligraphie \\
\hline
面具 & miànjù & n. & masque \\
\hline
雕塑 & diāosù & n. & sculpture \\
\hline
音乐 & yīnyuè & n. & musique \\
\hline
舞蹈 & wǔdǎo & n. & danse (nom) \\
\hline
跳舞 & tiàowǔ & v. & danser (verbe : sauter + danse) \\
\hline
颜色 & yánsè & n. & couleur \\
\hline
历史 & lìshǐ & n. & histoire \\
\hline
文化 & wénhuà & n. & culture \\
\hline
\end{tabularx}
\end{center}

\begin{attention}[Pièges de prononciation et de sens]
\begin{itemize}
\item \zh{音乐} se lit \emph{yīnyuè} (et non \emph{yīnlè}) : le caractère 乐 a deux lectures, \emph{yuè} dans 音乐 (musique) et \emph{lè} dans 快乐 (joyeux). C'est un caractère \cle{polyphonique} classique.
\item \zh{画} porte le quatrième ton (\emph{huà}) ; ne le confondez pas avec \zh{话} (\emph{huà}, parole) : même pinyin, même ton, mais sens différent ! 看画 = regarder des peintures ; 说话 = parler.
\item \zh{面具} : citation \emph{miànjù} (deux quatrièmes tons). En chaîne parlée, le premier quatrième ton devient un deuxième ton par \cle{changement de ton} (sandhi) : on prononce \emph{miánjù}. Entraînez-vous à bien descendre le second ton.
\item \zh{舞} (\emph{wǔ}) seul est surtout un morphème lié (littéraire) ; le nom courant est \zh{舞蹈} (\emph{wǔdǎo}) et le verbe est \zh{跳舞} (\emph{tiàowǔ}, litt. « sauter la danse »).
\end{itemize}
\end{attention}

\courssub{Exprimer son opinion : verbes et adjectifs}

Pour répondre à la question « Tu aimes l'art camerounais ou chinois ? Pourquoi ? », il ne suffit pas de nommer les arts : il faut \cle{donner son avis}. Voici les outils indispensables.

\begin{center}
\begin{tabularx}{\linewidth}{|X|X|X|X|}
\hline
\rowcolor{popblueL} Caractères & Pinyin & Nature & Traduction \\
\hline
喜欢 & xǐhuan & v. & aimer \\
\hline
觉得 & juéde & v. & trouver, penser (avis) \\
\hline
有意思 & yǒu yìsi & adj. & intéressant, amusant \\
\hline
好看 & hǎokàn & adj. & beau à voir, joli \\
\hline
美 & měi & adj. & beau (esthétique, soutenu) \\
\hline
有名 & yǒumíng & adj. & célèbre \\
\hline
\end{tabularx}
\end{center}

\textbf{Observation.} Comparez les deux phrases suivantes :

\zh{我喜欢面具。} \emph{Wǒ xǐhuan miànjù.} « J'aime les masques. » (verbe + nom)

\zh{我觉得面具很有意思。} \emph{Wǒ juéde miànjù hěn yǒu yìsi.} « Je trouve les masques très intéressants. » (verbe + phrase complète)

La différence est importante : \zh{喜欢} est suivi d'un \emph{nom} ou d'un \emph{verbe} (j'aime quelque chose, j'aime faire quelque chose), tandis que \zh{觉得} est suivi d'une \emph{proposition entière} (sujet + adjectif). En français, « trouver » fonctionne pareil : « je trouve \emph{que} les masques \emph{sont} intéressants ».

\begin{exemplebox}[Donner son avis : modèles à réutiliser]
\zh{我觉得面具很有意思。} \emph{Wǒ juéde miànjù hěn yǒu yìsi.} « Je trouve les masques très intéressants. »\\
\zh{我觉得中国画的颜色很美。} \emph{Wǒ juéde Zhōngguó huà de yánsè hěn měi.} « Je trouve les couleurs de la peinture chinoise très belles. »\\
\zh{中国的书法最有名。} \emph{Zhōngguó de shūfǎ zuì yǒumíng.} « La calligraphie chinoise est la plus célèbre. »\\
\zh{喀麦隆的舞蹈很好看。} \emph{Kāmàilóng de wǔdǎo hěn hǎokàn.} « Les danses camerounaises sont très belles à voir. »
\end{exemplebox}

\begin{attention}[L'adjectif seul ne suffit pas : le rôle de 很]
En français on dit « les masques sont intéressants » avec le verbe « être ». En chinois, un adjectif seul après le sujet sonne \emph{comparatif} ou incomplet : \zh{面具很有意思} utilise \zh{很} (\emph{hěn}) non pas pour dire « très », mais comme \cle{liaison obligatoire} entre le sujet et l'adjectif (sauf si l'adjectif est déjà modifié par 更, 最, 不, etc.). Retenez la structure : \textbf{Sujet + 很 + adjectif}. Ne traduisez jamais « être » par 是 devant un adjectif : \zh{面具是很有意思} est incorrect dans ce sens.
\end{attention}

\courssub{La question alternative avec 还是}

\begin{definition}[还是 háishi : le « ou » des questions]
\zh{还是} (\emph{háishi}) est une conjonction qui signifie « ou », mais elle s'emploie \cle{uniquement dans les questions alternatives}, où l'on propose un choix. Structure : \textbf{A + 还是 + B ?}\\
\zh{你喜欢喀麦隆还是中国艺术？} \emph{Nǐ xǐhuan Kāmàilóng háishi Zhōngguó yìshù ?} « Tu aimes l'art camerounais ou l'art chinois ? »\\
\zh{你喜欢书法还是音乐？} \emph{Nǐ xǐhuan shūfǎ háishi yīnyuè ?} « Tu aimes la calligraphie ou la musique ? »
\end{definition}

\begin{attention}[Erreur fréquente des francophones : 还是 vs 或者]
En français, le même mot « ou » sert dans les questions \emph{et} dans les affirmations. En chinois, c'est impossible :
\begin{itemize}
\item \textbf{Question} : 还是. \zh{你喜欢画还是雕塑？} « Tu aimes la peinture ou la sculpture ? »
\item \textbf{Affirmation} : 或者 (\emph{huòzhě}). \zh{我喜欢画或者雕塑。} « J'aime la peinture ou la sculpture. »
\end{itemize}
Dans la \textbf{réponse} à une question en 还是, on choisit une option \emph{sans} répéter 还是 : \zh{我喜欢中国艺术。} (\emph{Wǒ xǐhuan Zhōngguó yìshù.}) « J'aime l'art chinois. » Dire \zh{我喜欢还是中国艺术} est une faute typique du francophone.
\end{attention}

\courssub{Comparer et classer : 更 et 最}

\textbf{Observation.} Lisez ce mini-dialogue entre deux élèves de Yaoundé :

\zh{——你喜欢喀麦隆艺术还是中国艺术？} \emph{——Nǐ xǐhuan Kāmàilóng yìshù háishi Zhōngguó yìshù ?} « — Tu aimes l'art camerounais ou l'art chinois ? »

\zh{——我更喜欢中国艺术。我最喜欢书法。} \emph{——Wǒ gèng xǐhuan Zhōngguó yìshù. Wǒ zuì xǐhuan shūfǎ.} « — Je préfère l'art chinois. Ce que j'aime le plus, c'est la calligraphie. »

Repérez les deux petits mots placés \emph{avant} le verbe : \zh{更} (\emph{gèng}, encore plus) et \zh{最} (\emph{zuì}, le plus). Ils se placent toujours \textbf{juste devant l'adjectif ou le verbe} qu'ils modifient.

\begin{propriete}[更 + verbe/adjectif : la préférence comparative]
\zh{更} (\emph{gèng}) exprime le degré supérieur : « encore plus », et devant un verbe d'opinion il traduit « préférer ». Structure : \textbf{Sujet + 更 + verbe/adjectif}.\\
\zh{我更喜欢中国画。} \emph{Wǒ gèng xǐhuan Zhōngguó huà.} « Je préfère la peinture chinoise. »\\
\zh{我觉得雕塑更有意思。} \emph{Wǒ juéde diāosù gèng yǒu yìsi.} « Je trouve la sculpture plus intéressante. »\\
Remarque : avec 更, on n'utilise pas 很 devant l'adjectif (on dit \zh{更美}, jamais \zh{更很美}).
\end{propriete}

\begin{propriete}[最 + verbe/adjectif : le superlatif]
\zh{最} (\emph{zuì}) exprime le degré maximal : « le plus ». Structure : \textbf{Sujet + 最 + verbe/adjectif}.\\
\zh{我最喜欢书法。} \emph{Wǒ zuì xǐhuan shūfǎ.} « Ce que j'aime le plus, c'est la calligraphie. »\\
\zh{中国的书法最有名。} \emph{Zhōngguó de shūfǎ zuì yǒumíng.} « La calligraphie chinoise est la plus célèbre. »
\end{propriete}

\begin{center}
\begin{tabularx}{\linewidth}{|X|X|X|}
\hline
\rowcolor{popblueL} Structure & Exemple (caractères + pinyin) & Traduction \\
\hline
A + 还是 + B ? & \zh{你喜欢音乐还是舞蹈？} Nǐ xǐhuan yīnyuè háishi wǔdǎo ? & Tu aimes la musique ou la danse ? \\
\hline
Sujet + 更 + verbe & \zh{我更喜欢音乐。} Wǒ gèng xǐhuan yīnyuè. & Je préfère la musique. \\
\hline
Sujet + 最 + verbe & \zh{我最喜欢跳舞。} Wǒ zuì xǐhuan tiàowǔ. & Ce que j'aime le plus, c'est danser. \\
\hline
Sujet + 最 + adjectif & \zh{面具最有意思。} Miànjù zuì yǒu yìsi. & Les masques sont les plus intéressants. \\
\hline
\end{tabularx}
\end{center}

\courssub{Carte mentale du lexique}

\begin{popfigure}
\begin{center}
\begin{tikzpicture}[
  mindmap,
  concept color=popblue,
  every node/.style={concept, align=center, font=\small},
  level 1 concept/.append style={level distance=3.4cm, sibling angle=90},
  level 2 concept/.append style={level distance=2.4cm},
  grow cyclic
]
\node[concept color=popblue, text=white, align=center]{\zh{艺术}\\ yìshù}
  child[concept color=popgreen] { node {arts visuels}
    child { node {\zh{画} huà} }
    child { node {\zh{书法} shūfǎ} }
    child { node {\zh{雕塑} diāosù} }
    child { node {\zh{面具} miànjù} }
  }
  child[concept color=poporange] { node[align=center]{musique\\ et danse}
    child { node {\zh{音乐} yīnyuè} }
    child { node {\zh{舞蹈} wǔdǎo} }
    child { node {\zh{跳舞} tiàowǔ} }
  }
  child[concept color=poppurple] { node {opinions}
    child { node {\zh{喜欢} xǐhuan} }
    child { node {\zh{觉得} juéde} }
    child { node {\zh{有意思} yǒu yìsi} }
  }
  child[concept color=popgold] { node {qualités}
    child { node {\zh{好看} hǎokàn} }
    child { node {\zh{美} měi} }
    child { node {\zh{有名} yǒumíng} }
  };
\end{tikzpicture}
\end{center}
\legende{Carte mentale du lexique de la leçon : quatre familles de mots autour de 艺术.}
\end{popfigure}

\courssub{Les collocations utiles : les mots qui vont ensemble}

En chinois comme en français, certains verbes vont naturellement avec certains noms : ce sont les \cle{collocations}. Les apprendre par paires évite les traductions mot à mot incorrectes.

\begin{center}
\begin{tabularx}{\linewidth}{|X|X|X|}
\hline
\rowcolor{popblueL} Collocation & Pinyin & Traduction \\
\hline
\zh{看画} & kàn huà & regarder des peintures (visiter une expo) \\
\hline
\zh{画画} & huà huà & peindre, dessiner (activité) \\
\hline
\zh{写字 / 写书法 / 练书法} & xiě zì / xiě shūfǎ / liàn shūfǎ & écrire des caractères / faire de la calligraphie / s'entraîner à la calligraphie \\
\hline
\zh{听音乐} & tīng yīnyuè & écouter de la musique \\
\hline
\zh{跳舞} & tiào wǔ & danser \\
\hline
\zh{做雕塑 / 雕刻} & zuò diāosù / diāokè & faire de la sculpture / sculpter \\
\hline
\zh{有几千年的历史} & yǒu jǐ qiān nián de lìshǐ & avoir plusieurs millénaires d'histoire \\
\hline
\end{tabularx}
\end{center}

\begin{exemplebox}[Chaque collocation dans une phrase]
\zh{星期天我喜欢去美术馆看画。} \emph{Xīngqītiān wǒ xǐhuan qù měishùguǎn kàn huà.} « Le dimanche, j'aime aller au musée regarder des peintures. »\\
\zh{我的老师每天练书法。} \emph{Wǒ de lǎoshī měitiān liàn shūfǎ.} « Mon professeur s'entraîne à la calligraphie tous les jours. »\\
\zh{她喜欢听音乐，不喜欢跳舞。} \emph{Tā xǐhuan tīng yīnyuè, bù xǐhuan tiào wǔ.} « Elle aime écouter de la musique, mais n'aime pas danser. »\\
\zh{中国的书法有几千年的历史。} \emph{Zhōngguó de shūfǎ yǒu jǐ qiān nián de lìshǐ.} « La calligraphie chinoise a plusieurs millénaires d'histoire. »
\end{exemplebox}

\begin{attention}[Le bon verbe pour le bon art : pas de calque sur « faire »]
Ne traduisez pas « faire » partout ! Chaque art a son verbe spécifique :
\begin{itemize}
\item Peinture : \zh{看画} (regarder), \zh{画画} (pratiquer).
\item Calligraphie : \zh{写字}, \zh{写书法}, \zh{练书法}.
\item Musique : \zh{听音乐} (écouter).
\item Danse : \zh{跳舞} (danser — litt. « sauter »).
\item Sculpture : \zh{做雕塑} (faire de la sculpture), \zh{雕刻} (sculpter, tailler).
\end{itemize}
Dire \zh{做音乐} pour « faire de la musique » ou \zh{做舞} pour « danser » est un calque du français à proscrire.
\end{attention}

\courssub{Texte de réinvestissement}

\begin{textebox}[Les arts chez moi et en Chine]
\shortstack[l]{\zh{我叫玛丽，我住在雅温得。我非常喜欢艺术。喀麦隆的面具和雕塑很有名，我觉得它们很有意思。我们的舞蹈也很好看，颜色很美。}\\ \zh{中国的文化不一样：中国的书法有几千年的历史，中国画也很有名。}\\ \zh{你问我喜欢喀麦隆艺术还是中国艺术？我更喜欢喀麦隆的舞蹈，但是我最喜欢中国的书法。我觉得写书法很有意思，所以我每天写字。艺术没有国界，它让我们了解别的文化。}}\\
\emph{Wǒ jiào Mǎlì, wǒ zhù zài Yǎwēndé. Wǒ fēicháng xǐhuan yìshù. Kāmàilóng de miànjù hé diāosù hěn yǒumíng, wǒ juéde tāmen hěn yǒu yìsi. Wǒmen de wǔdǎo yě hěn hǎokàn, yánsè hěn měi. Zhōngguó de wénhuà bù yīyàng : Zhōngguó de shūfǎ yǒu jǐ qiān nián de lìshǐ, Zhōngguó huà yě hěn yǒumíng. Nǐ wèn wǒ xǐhuan Kāmàilóng yìshù háishi Zhōngguó yìshù ? Wǒ gèng xǐhuan Kāmàilóng de wǔdǎo, dànshì wǒ zuì xǐhuan Zhōngguó de shūfǎ. Wǒ juéde xiě shūfǎ hěn yǒu yìsi, suǒyǐ wǒ měitiān xiě zì. Yìshù méiyǒu guójiè, tā ràng wǒmen liǎojiě bié de wénhuà.}\\
« Je m'appelle Marie et j'habite à Yaoundé. J'aime beaucoup l'art. Les masques et les sculptures du Cameroun sont très célèbres, je les trouve très intéressants. Nos danses aussi sont belles à voir, les couleurs sont magnifiques. La culture chinoise est différente : la calligraphie chinoise a plusieurs millénaires d'histoire, et la peinture chinoise est aussi très célèbre. Tu me demandes si j'aime l'art camerounais ou l'art chinois ? Je préfère les danses camerounaises, mais ce que j'aime le plus, c'est la calligraphie chinoise. Je trouve que faire de la calligraphie est très intéressant, alors j'écris des caractères tous les jours. L'art n'a pas de frontières, il nous fait connaître d'autres cultures. »
\end{textebox}

\textbf{Glossaire du texte.} 住 \emph{zhù} (habiter) ; 非常 \emph{fēicháng} (extrêmement) ; 不一样 \emph{bù yīyàng} (différent) ; 但是 \emph{dànshì} (mais) ; 所以 \emph{suǒyǐ} (donc) ; 国界 \emph{guójiè} (frontière) ; 了解 \emph{liǎojiě} (comprendre, connaître) ; 美术馆 \emph{měishùguǎn} (musée des beaux-arts) ; 练 \emph{liàn} (s'entraîner, pratiquer).

\textbf{Questions de compréhension.}
\begin{enumerate}
\item \zh{玛丽觉得喀麦隆的面具怎么样？} (Que pense Marie des masques camerounais ?)
\item \zh{她最喜欢什么艺术？} (Quel art aime-t-elle le plus ?)
\item \zh{她为什么每天写字？} (Pourquoi écrit-elle des caractères tous les jours ?)
\end{enumerate}

\courssub{Exercice résolu et entraînement}

\begin{exoresolu}[Compléter avec 更, 最 ou 还是]
Énoncé. Complétez les phrases avec \zh{更}, \zh{最} ou \zh{还是}, puis traduisez-les en français.\\
1. 你喜欢音乐\_\_\_\_舞蹈？\\
2. 我\_\_\_\_喜欢书法。\\
3. 我觉得面具\_\_\_\_有意思。
\tcblower
Solution détaillée.\\
1. \zh{还是} : c'est une \emph{question alternative} entre deux options (musique / danse). \zh{你喜欢音乐还是舞蹈？} \emph{Nǐ xǐhuan yīnyuè háishi wǔdǎo ?} « Tu aimes la musique ou la danse ? »\\
2. \zh{最} : phrase affirmative qui exprime le degré maximal. \zh{我最喜欢书法。} \emph{Wǒ zuì xǐhuan shūfǎ.} « Ce que j'aime le plus, c'est la calligraphie. » (\zh{更} serait possible avec un contexte de comparaison explicite, mais seul \zh{最} donne un sens complet isolé.)\\
3. \zh{最} : devant l'adjectif \zh{有意思}, on exprime le superlatif. \zh{我觉得面具最有意思。} \emph{Wǒ juéde miànjù zuì yǒu yìsi.} « Je trouve que les masques sont les plus intéressants. »
\end{exoresolu}

\begin{exercice}[À vous de jouer]
1. Traduisez en chinois : « Je préfère la musique camerounaise. » ; « La sculpture est la plus célèbre. » ; « Tu aimes la peinture ou la calligraphie ? »\\
2. Construisez deux phrases avec \zh{觉得} pour donner votre avis sur les masques et sur la danse.\\
3. Répondez en une phrase complète (sans \zh{还是}) à la question : \zh{你喜欢喀麦隆艺术还是中国艺术？}
\end{exercice}

\begin{corrige}[Exercice « À vous de jouer »]
1. \zh{我更喜欢喀麦隆音乐。} \emph{Wǒ gèng xǐhuan Kāmàilóng yīnyuè.} ; \zh{雕塑最有名。} \emph{Diāosù zuì yǒumíng.} ; \zh{你喜欢画还是书法？} \emph{Nǐ xǐhuan huà háishi shūfǎ ?}\\
2. Exemples possibles : \zh{我觉得面具很有意思。} ; \zh{我觉得喀麦隆的舞蹈很好看。}\\
3. Exemple : \zh{我更喜欢中国艺术。} ou \zh{我更喜欢喀麦隆艺术。} — une seule option, sans \zh{还是}.
\end{corrige}

Vous disposez maintenant de tout le lexique nécessaire : les noms des arts, les verbes d'opinion, la question alternative avec \zh{还是}, et les degrés de préférence avec \zh{更} et \zh{最}. Dans la prochaine section, nous apprendrons à \cle{organiser} ces mots dans un texte structuré grâce aux connecteurs.

\coursec{Texte modèle et compréhension}

Cette leçon vous guide pour rédiger un petit texte argumentatif en chinois sur le thème : \zh{你喜欢喀麦隆还是中国艺术？为什么？} (Nǐ xǐhuan Kāmàilóng háishi Zhōngguó yìshù? Wèishénme?). Nous commençons par l'étude d'un texte modèle authentique, simple et bien structuré, qui servira de référence tout au long de la leçon.

\begin{textebox}[Texte modèle : 我喜欢的艺术]
喀麦隆和中国都有很美的艺术。喀麦隆的木雕很有名，中国的画和书法也很好看。\\
Kāmàilóng hé Zhōngguó dōu yǒu hěn měi de yìshù. Kāmàilóng de mùdiāo hěn yǒumíng, Zhōngguó de huà hé shūfǎ yě hěn hǎokàn.\\
« Le Cameroun et la Chine ont tous deux un art très beau. La sculpture sur bois du Cameroun est très célèbre, la peinture et la calligraphie chinoises sont aussi très belles. »\\[0.3em]
对我来说，我更喜欢中国艺术。为什么呢？\\
Duì wǒ láishuō, wǒ gèng xǐhuan Zhōngguó yìshù. Wèishénme ne?\\
« Pour moi, je préfère l'art chinois. Pourquoi donc ? »\\[0.3em]
首先，中国画的颜色很好看。然后，中国书法很有意思，每一个字都像一幅画。还有，中国艺术有几千年的历史。\\
Shǒuxiān, Zhōngguó huà de yánsè hěn hǎokàn. Ránhòu, Zhōngguó shūfǎ hěn yǒuyìsi, měi yí ge zì dōu xiàng yì fú huà. Háiyǒu, Zhōngguó yìshù yǒu jǐ qiān nián de lìshǐ.\\
« D'abord, les couleurs de la peinture chinoise sont très belles. Ensuite, la calligraphie chinoise est très intéressante, chaque caractère ressemble à un tableau. De plus, l'art chinois a plusieurs millénaires d'histoire. »\\[0.3em]
因为中国艺术又美又有历史，所以我很喜欢它。所以，我喜欢中国艺术。你呢？\\
Yīnwèi Zhōngguó yìshù yòu měi yòu yǒu lìshǐ, suǒyǐ wǒ hěn xǐhuan tā. Suǒyǐ, wǒ xǐhuan Zhōngguó yìshù. Nǐ ne?\\
« Parce que l'art chinois est à la fois beau et chargé d'histoire, je l'aime beaucoup. Donc, j'aime l'art chinois. Et toi ? »
\end{textebox}

\textbf{Questions de compréhension et d'observation :}
\begin{enumerate}
\item Quelle est l'opinion de l'auteur ? Avec quelle phrase l'annonce-t-il explicitement ?
\item Combien d'arguments donne-t-il pour justifier son choix ? Souligne dans le texte les mots qui les introduisent.
\item Comment le texte se termine-t-il ? À qui s'adresse la dernière question et quel est son effet ?
\item Repère la structure \zh{因为...所以...} dans le texte. Quelle en est la traduction naturelle en français ?
\end{enumerate}

\begin{savaistu}[Arts traditionnels : Cameroun et Chine]
Le \textbf{Cameroun} est réputé pour sa \cle{sculpture sur bois} (\zh{木雕 mùdiāo}), notamment chez les Bamileke (masques, statues de chefs, pirogues sculptées) et dans la région de l'Ouest. Les \cle{tissus traditionnels} comme le \emph{ndop} (tissu bleu à motifs blancs) ou le \emph{toghu} (velours brodé du Nord-Ouest) sont aussi des arts majeurs.
La \textbf{Chine} est le berceau de la \cle{calligraphie} (\zh{书法 shūfǎ}, « méthode de l'écriture ») et de la \cle{peinture à l'encre} (\zh{中国画 zhōngguó huà}), qui partagent les mêmes outils (pinceau, encre, papier, pierre à encre) et la même esthétique du trait (\zh{笔触 bǐchù}). La calligraphie est considérée comme l'art suprême, révélant le caractère de l'artiste.
\end{savaistu}

\begin{center}
\begin{tabularx}{\linewidth}{|X|X|X|X|}\hline
\rowcolor{popblueL} Caractères & Pinyin & Nature & Traduction \\ \hline
艺术 & yìshù & nom & art \\ \hline
木雕 & mùdiāo & nom & sculpture sur bois \\ \hline
有名 & yǒumíng & adjectif & célèbre, connu \\ \hline
画 & huà & nom/verbe & peinture, tableau / peindre \\ \hline
书法 & shūfǎ & nom & calligraphie \\ \hline
好看 & hǎokàn & adjectif & beau (visuellement) \\ \hline
颜色 & yánsè & nom & couleur \\ \hline
有意思 & yǒuyìsi & adjectif & intéressant, amusant \\ \hline
像 & xiàng & verbe & ressembler à \\ \hline
历史 & lìshǐ & nom & histoire \\ \hline
几千年 & jǐ qiān nián & locution & plusieurs millénaires \\ \hline
又...又... & yòu... yòu... & structure & à la fois... et... \\ \hline
\end{tabularx}
\end{center}

\coursec{Lexique du thème : arts, opinions et comparaisons}

Pour exprimer ses goûts artistiques et comparer deux cultures, il faut maîtriser un lexique précis : noms d'arts, verbes d'appréciation, adjectifs de description, et structures de comparaison.

\courssub{Les noms d'arts et disciplines}

\begin{center}
\begin{tabularx}{\linewidth}{|X|X|X|X|}\hline
\rowcolor{popblueL} Caractères & Pinyin & Nature & Traduction \\ \hline
音乐 & yīnyuè & nom & musique \\ \hline
舞蹈 & wǔdǎo & nom & danse \\ \hline
戏剧 & xìjù & nom & théâtre, opéra \\ \hline
雕塑 & diāosù & nom & sculpture (général) \\ \hline
陶艺 & táoyì & nom & céramique, poterie d'art \\ \hline
剪纸 & jiǎnzhǐ & nom & découpage de papier \\ \hline
建筑 & jiànzhù & nom & architecture \\ \hline
传统艺术 & chuántǒng yìshù & nom & art traditionnel \\ \hline
现代艺术 & xiàndài yìshù & nom & art moderne \\ \hline
\end{tabularx}
\end{center}

\courssub{Verbes et adjectifs pour exprimer son opinion}

\begin{center}
\begin{tabularx}{\linewidth}{|X|X|X|X|}\hline
\rowcolor{popblueL} Caractères & Pinyin & Nature & Traduction \\ \hline
喜欢 & xǐhuan & verbe & aimer, apprécier \\ \hline
爱 & ài & verbe & aimer (plus fort que \zh{喜欢}) \\ \hline
觉得 & juéde & verbe & penser, trouver (avis subjectif) \\ \hline
认为 & rènwéi & verbe & considérer que, estimer que (plus formel) \\ \hline
欣赏 & xīnshǎng & verbe & apprécier (esthétiquement) \\ \hline
漂亮 & piàoliang & adjectif & joli, beau (visuel) \\ \hline
美 & měi & adjectif & beau (esthétique, littéraire) \\ \hline
有名 & yǒumíng & adjectif & célèbre \\ \hline
独特 & dútè & adjectif & unique, particulier \\ \hline
精彩 & jīngcǎi & adjectif & splendide, remarquable \\ \hline
\end{tabularx}
\end{center}

\courssub{Structures de comparaison : 比 et 更}

\begin{propriete}[Comparatif d'inégalité avec \zh{比} (bǐ)]
Structure : \textbf{Sujet A + \zh{比} + Sujet B + Adjectif / Verbe d'état}. \emph{Pas de « de » après l'adjectif.}
\zh{中国艺术比喀麦隆艺术古老。} (Zhōngguó yìshù bǐ Kāmàilóng yìshù gǔlǎo.) — « L'art chinois est plus ancien que l'art camerounais. »
\zh{我觉得中国书法比画难。} (Wǒ juéde Zhōngguó shūfǎ bǐ huà nán.) — « Je trouve la calligraphie chinoise plus difficile que la peinture. »
\end{propriete}

\begin{propriete}[Comparatif avec \zh{更} (gèng, « encore plus »)]
\zh{更} se place \textbf{devant} l'adjectif ou le verbe d'état pour intensifier la comparaison, souvent après \zh{比} ou seul si le terme de comparaison est connu.
\zh{我更喜欢中国艺术。} (Wǒ gèng xǐhuan Zhōngguó yìshù.) — « Je préfère l'art chinois. » (Litt. : « J'aime \emph{encore plus} l'art chinois. »)
\zh{喀麦隆的木雕很有名，中国的书法更有名。} (Kāmàilóng de mùdiāo hěn yǒumíng, Zhōngguó de shūfǎ gèng yǒumíng.) — « La sculpture camerounaise est célèbre, la calligraphie chinoise l'est encore plus. »
\end{propriete}

\begin{attention}[Pièges classiques des francophones sur la comparaison]
\begin{itemize}
\item \textbf{N'ajoute pas \zh{的} après l'adjectif dans la structure \zh{比}} : \zh{他比我高} (Tā bǐ wǒ gāo), \textbf{pas} \zh{他比我高的}.
\item \textbf{\zh{更} vs \zh{比}} : \zh{比} introduit le terme de comparaison (B) ; \zh{more} intensifie la qualité. On peut les combiner : \zh{A 比 B 更...}.
\item \textbf{\zh{喜欢} vs \zh{爱}} : \zh{爱} est plus fort, plus affectif (famille, patrie, passion). Pour l'art, \zh{喜欢} ou \zh{欣赏} sont plus naturels.
\item \textbf{Adjectifs à un seul caractère} : \zh{美} (beau), \zh{旧} (vieux), \zh{新} (neuf), \zh{贵} (cher). En chinois, un adjectif monosyllabique seul après le sujet sonne incomplet. Ajoute \zh{很} (hěn, « très », souvent vide de sens purement intensif) : \zh{中国艺术很美}. Dans une phrase \zh{比}, \zh{很} disparaît : \zh{中国艺术比喀麦隆艺术美}.
\end{itemize}
\end{attention}

\coursec{Structure du texte et connecteurs}

Dans la section précédente, nous avons acquis le vocabulaire des arts et des opinions. Pour répondre par écrit à la question \zh{你喜欢喀麦隆还是中国艺术？为什么？}, il faut savoir \cle{organiser} ces mots dans un texte clair et logique. C'est l'objet de cette section : apprendre le plan type d'un petit texte argumentatif chinois et maîtriser les connecteurs qui le font tenir debout.

\courssub{Le plan type en quatre parties}

Le texte modèle de la section 1 suit un plan en quatre parties, très facile à réutiliser dans n'importe quelle production écrite d'opinion :

\begin{enumerate}
\item \textbf{Introduction générale} (\zh{引言 yǐnyán}) : on présente le sujet sans donner encore son avis. Modèle : \zh{喀麦隆和中国都有...} (Kāmàilóng hé Zhōngguó dōu yǒu... — « Le Cameroun et la Chine ont tous deux... »). Le mot \zh{都} (dōu, « tous les deux ») est idéal pour comparer deux pays ou deux arts.
\item \textbf{Annonce de la préférence} (\zh{观点 guāndiǎn}) : on donne clairement son opinion. Modèles : \zh{我更喜欢...} (wǒ gèng xǐhuan... — « je préfère... ») ou \zh{对我来说，...} (duì wǒ láishuō... — « pour moi... »).
\item \textbf{Les arguments} (\zh{理由 lǐyóu}) : chaque argument est introduit par un connecteur : \zh{首先} (d'abord), \zh{然后} (ensuite), \zh{还有} (de plus). La question \zh{为什么呢？} sert de transition entre l'opinion et les arguments.
\item \textbf{Conclusion} (\zh{结论 jiélùn}) : on reprend l'opinion avec \zh{所以} (donc), puis on pose une question au lecteur : \zh{你呢？} (Nǐ ne? — « Et toi ? »).
\end{enumerate}

\begin{popfigure}
\begin{tikzpicture}[
  bloc/.style={draw=popblue, thick, fill=popblueL, rounded corners=3pt, align=center, minimum width=3.5cm, minimum height=1.2cm, font=\small},
  fleche/.style={-{Stealth[length=2.5mm]}, thick, popdark}]
\node[bloc, align=center] (intro) at (0,0){引言 yǐnyán\\Introduction\\\zh{喀麦隆和中国都有...}};
\node[bloc, fill=poporangeL, draw=poporange, align=center] (opinion) at (5.5,0){观点 guāndiǎn\\Opinion\\\zh{我更喜欢...}};
\node[bloc, fill=popgreenL, draw=popgreen, align=center] (raisons) at (11,0){理由 lǐyóu\\Arguments\\\zh{首先 $\rightarrow$ 然后 $\rightarrow$ 还有}};
\node[bloc, fill=poppinkL, draw=poppink, align=center] (concl) at (11,-2.8){结论 jiélùn\\Conclusion\\\zh{所以... 你呢？}};
\draw[fleche] (intro) -- (opinion);
\draw[fleche] (opinion) -- (raisons);
\draw[fleche] (raisons) -- (concl);
\draw[fleche, dashed, poppurple] (opinion.south) to[bend right=15] node[below, font=\footnotesize, align=center, text=poppurple]{transition :\\\zh{为什么呢？}} (raisons.south west);
\end{tikzpicture}
\legende{Organigramme du texte argumentatif simple : introduction $\rightarrow$ opinion $\rightarrow$ arguments $\rightarrow$ conclusion.}
\end{popfigure}

\begin{methode}[Structurer un texte argumentatif simple en quatre étapes]
\begin{enumerate}
\item \textbf{Introduction} : une phrase générale sur le sujet, sans opinion. Utilise \zh{A和B都有...} ou \zh{A和B都很有名}.
\item \textbf{Opinion} : une phrase claire avec \zh{我觉得...} ou \zh{对我来说，我更喜欢...}. Ajoute la transition \zh{为什么呢？} pour annoncer tes arguments.
\item \textbf{Arguments} : une idée = une phrase. Chaque argument est introduit par un connecteur : \zh{首先} pour le premier, \zh{然后} pour le suivant, \zh{还有} pour ajouter une idée ou un exemple.
\item \textbf{Conclusion} : reprends ton opinion avec \zh{所以} (\zh{所以，我喜欢...}), puis termine par la question au lecteur \zh{你呢？}.
\end{enumerate}
Règle d'or : en chinois comme en français, un bon petit texte n'est pas un texte compliqué, c'est un texte \cle{bien ordonné}.
\end{methode}

\courssub{Les connecteurs un par un}

\begin{definition}[首先 shǒuxiān — « d'abord, premièrement »]
\zh{首先} est un connecteur qui \cle{ouvre le premier argument}. Il se place en tête de phrase, suivi d'une virgule chinoise \zh{，}. Il annonce au lecteur : « voici ma première raison ».
\end{definition}

\begin{exemplebox}[首先 en tête de phrase]
\zh{首先，中国画的颜色很好看。}\\
\emph{Shǒuxiān, Zhōngguó huà de yánsè hěn hǎokàn.}\\
« D'abord, les couleurs de la peinture chinoise sont très belles. »\\[0.3em]
\zh{首先，喀麦隆的音乐很有名。}\\
\emph{Shǒuxiān, Kāmàilóng de yīnyuè hěn yǒumíng.}\\
« D'abord, la musique camerounaise est très célèbre. »
\end{exemplebox}

\begin{definition}[然后 ránhòu — « ensuite, puis »]
\zh{然后} \cle{enchaîne les arguments} (ou les actions). Il se place aussi en tête de phrase, suivi d'une virgule. Attention : \zh{然后} indique simplement l'ordre (chronologique ou logique), il ne marque pas la cause.
\end{definition}

\begin{exemplebox}[然后 pour enchaîner]
\zh{然后，中国书法很有意思。}\\
\emph{Ránhòu, Zhōngguó shūfǎ hěn yǒuyìsi.}\\
« Ensuite, la calligraphie chinoise est très intéressante. »\\[0.3em]
\zh{我们先去了博物馆，然后去了艺术市场。}\\
\emph{Wǒmen xiān qùle bówùguǎn, ránhòu qùle yìshù shìchǎng.}\\
« Nous sommes d'abord allés au musée, puis au marché d'art. »
\end{exemplebox}

\begin{definition}[还有 háiyǒu — « de plus, il y a aussi »]
\zh{还有} \cle{ajoute} un argument ou un exemple supplémentaire. C'est le connecteur de l'addition : il signifie littéralement « il y a encore ». Il se place en tête de phrase, suivi d'une virgule.
\end{definition}

\begin{exemplebox}[还有 pour ajouter]
\zh{还有，中国艺术有几千年的历史。}\\
\emph{Háiyǒu, Zhōngguó yìshù yǒu jǐ qiān nián de lìshǐ.}\\
« De plus, l'art chinois a plusieurs millénaires d'histoire. »\\[0.3em]
\zh{还有，我的中国朋友教我写汉字。}\\
\emph{Háiyǒu, wǒ de Zhōngguó péngyou jiāo wǒ xiě Hànzì.}\\
« De plus, mon ami chinois m'apprend à écrire les caractères chinois. »
\end{exemplebox}

\begin{propriete}[所以 suǒyǐ — « donc » : conclure, et la paire 因为...所以...]
\zh{所以} (suǒyǐ, « donc, c'est pourquoi ») introduit la conséquence ou la conclusion. Deux emplois :
\begin{itemize}
\item \textbf{La paire} \zh{因为...所以...} : \zh{因为} (yīnwèi, « parce que ») introduit la cause, \zh{所以} introduit la conséquence, \emph{dans la même phrase}. Structure : \zh{因为} + cause，\zh{所以} + conséquence。
\item \zh{所以} \textbf{seul}, en tête de phrase de conclusion : il reprend l'opinion annoncée au début.
\end{itemize}
Exemple de la paire : \zh{因为中国艺术又美又有历史，所以我很喜欢它。} (Yīnwèi Zhōngguó yìshù yòu měi yòu yǒu lìshǐ, suǒyǐ wǒ hěn xǐhuan tā. — « Parce que l'art chinois est à la fois beau et chargé d'histoire, je l'aime beaucoup. »)
\end{propriete}

\begin{attention}[因为...所以... : une paire correcte en chinois, fautive en français]
En français, on ne peut pas dire « \emph{Parce qu'il est beau, donc je l'aime} » : « parce que » et « donc » ne cohabitent pas. En chinois, c'est l'inverse : la paire \zh{因为...所以...} est tout à fait correcte et même très naturelle. Deux pièges pour les francophones :
\begin{itemize}
\item Ne traduis pas mot à mot : \zh{因为...所以...} se traduit par « parce que... » seul, ou par « comme..., ... ».
\item Tu peux aussi utiliser \zh{因为} seul ou \zh{所以} seul : la paire complète n'est pas obligatoire, mais elle est toujours correcte.
\end{itemize}
\end{attention}

\courssub{Formules d'ouverture et de clôture}

\begin{definition}[Formules d'ouverture : 我觉得... / 对我来说，...]
Pour donner son opinion poliment et clairement :
\begin{itemize}
\item \zh{我觉得...} (wǒ juéde... — « je pense que..., je trouve que... ») : la formule d'opinion la plus courante.
\item \zh{对我来说，...} (duì wǒ láishuō... — « pour moi..., en ce qui me concerne... ») : introduit un point de vue personnel, souvent suivi de \zh{我更喜欢...}.
\end{itemize}
\end{definition}

\begin{propriete}[La question 那些？ : une excellente transition]
\zh{那些？} (Wèishénme ne? — « Pourquoi donc ? ») placée juste après l'opinion annonce les arguments qui vont suivre. Le mot-outil final \zh{呢} (ne) adoucit la question et la rend plus naturelle. Séquence type : \zh{我更喜欢中国艺术。那些？首先...} (« Je préfère l'art chinois. Pourquoi donc ? D'abord... »).
\end{propriete}

\begin{definition}[Formule de clôture : 所以，我喜欢... 你呢？]
La conclusion reprend l'opinion avec \zh{所以}, puis s'adresse au lecteur : \zh{你呢？} (Nǐ ne? — « Et toi ? »). Cette petite question finale est très appréciée en chinois : elle ouvre le dialogue et donne un texte vivant. Modèle complet : \zh{所以，我喜欢中国艺术。你呢？} (Suǒyǐ, wǒ xǐhuan Zhōngguó yìshù. Nǐ ne? — « Donc, j'aime l'art chinois. Et toi ? »).
\end{definition}

\begin{center}
\begin{tabularx}{\linewidth}{|X|X|X|}\hline
\rowcolor{popblueL} Fonction & Connecteur chinois + pinyin & Équivalent français \\ \hline
Ouvrir le 1\ier{} argument & \zh{首先} shǒuxiān & d'abord, premièrement \\ \hline
Enchaîner les arguments & \zh{然后} ránhòu & ensuite, puis \\ \hline
Ajouter un argument & \zh{还有} háiyǒu & de plus, en outre \\ \hline
Donner son opinion & \zh{我觉得} wǒ juéde & je pense que \\ \hline
Introduire son point de vue & \zh{对我来说} duì wǒ láishuō & pour moi \\ \hline
Transition vers les raisons & \zh{那些？} wèishénme ne & pourquoi donc ? \\ \hline
Conclure (conséquence) & \zh{所以} suǒyǐ & donc, c'est pourquoi \\ \hline
Questionner le lecteur & \zh{你呢？} nǐ ne & et toi ? \\ \hline
Cause (dans paire) & \zh{因为} yīnwèi & parce que \\ \hline
\end{tabularx}
\end{center}

\begin{exoresolu}[Remettre un texte dans l'ordre]
Énoncé : voici les quatre phrases d'un petit texte, dans le désordre. Remets-les dans l'ordre logique et indique la fonction de chacune (introduction, opinion, argument, conclusion).\\
(a) \zh{首先，喀麦隆的木雕很漂亮。} (Shǒuxiān, Kāmàilóng de mùdiāo hěn piàoliang.)\\
(b) \zh{所以，我喜欢喀麦隆艺术。你呢？} (Suǒyǐ, wǒ xǐhuan Kāmàilóng yìshù. Nǐ ne?)\\
(c) \zh{喀麦隆和中国都有很好的艺术。} (Kāmàilóng hé Zhōngguó dōu yǒu hěn hǎo de yìshù.)\\
(d) \zh{对我来说，我更喜欢喀麦隆艺术。} (Duì wǒ láishuō, wǒ gèng xǐhuan Kāmàilóng yìshù.)
\tcblower
Solution détaillée : l'ordre correct est \textbf{(c) $\rightarrow$ (d) $\rightarrow$ (a) $\rightarrow$ (b)}.\\
1. (c) est l'\textbf{introduction} : elle présente le sujet avec la structure \zh{A和B都有...}, sans donner d'opinion.\\
2. (d) est l'\textbf{opinion} : la formule \zh{对我来说} + \zh{我更喜欢} annonce clairement la préférence.\\
3. (a) est le premier \textbf{argument}, signalé par le connecteur \zh{首先}.\\
4. (b) est la \textbf{conclusion} : \zh{所以} reprend l'opinion et \zh{你呢？} s'adresse au lecteur.\\
Traduction du texte reconstitué : « Le Cameroun et la Chine ont tous deux un très bon art. Pour moi, je préfère l'art camerounais. D'abord, la sculpture sur bois du Cameroun est très belle. Donc, j'aime l'art camerounais. Et toi ? »
\end{exoresolu}

\begin{exercice}[À toi de jouer]
1. Traduis en chinois, en utilisant les connecteurs étudiés : « D'abord, la musique camerounaise est très belle. Ensuite, j'aime beaucoup la sculpture sur bois. De plus, l'art camerounais a une longue histoire. Donc, j'aime l'art camerounais. Et toi ? »\\
2. Complète avec le connecteur qui convient (\zh{首先}, \zh{然后}, \zh{还有}, \zh{所以}) : \zh{我更喜欢中国画。\_\_\_\_，颜色很好看。\_\_\_\_，书法很有意思。\_\_\_\_，中国画有几千年的历史。\_\_\_\_，我喜欢中国艺术。}
\end{exercice}

\begin{corrige}[Exercice « À toi de jouer »]
1. \zh{首先，喀麦隆的音乐很好听。然后，我很喜欢木雕。还有，喀麦隆艺术有很长的历史。所以，我喜欢喀麦隆艺术。你呢？} (Shǒuxiān, Kāmàilóng de yīnyuè hěn hǎotīng. Ránhòu, wǒ hěn xǐhuan mùdiāo. Háiyǒu, Kāmàilóng yìshù yǒu hěn cháng de lìshǐ. Suǒyǐ, wǒ xǐhuan Kāmàilóng yìshù. Nǐ ne?) — Remarque : pour la musique, on dit \zh{好听} (hǎotīng, « beau à écouter ») et non \zh{好看}.\\
2. \zh{首先} (premier argument) ; \zh{然后} (enchaînement) ; \zh{还有} (ajout) ; \zh{所以} (conclusion).
\end{corrige}

\begin{attention}[Les trois erreurs fréquentes des francophones]
\begin{itemize}
\item \textbf{Oublier la virgule après le connecteur} : écris \zh{首先，...} et non \zh{首先...} collé à la phrase.
\item \textbf{Empiler les connecteurs} : ne dis pas \zh{首先然后...} ; un seul connecteur par phrase suffit.
\item \textbf{Traduire « et toi ? » par} \zh{和你？} : la formule correcte est \zh{你呢？}, avec le mot-outil \zh{呢}.
\end{itemize}
\end{attention}

\begin{aretenir}[L'essentiel de la section]
Un texte d'opinion simple se construit en quatre parties : introduction (\zh{A和B都有...}), opinion (\zh{对我来说，我更喜欢...}), arguments (\zh{首先...然后...还有...}) et conclusion (\zh{所以，我喜欢...你呢？}). La paire \zh{因为...所以...} est correcte en chinois, et la question \zh{那些？} est la transition idéale entre l'opinion et les arguments. Retiens : une idée = une phrase = un connecteur.
\end{aretenir}

\coursec{Écriture des caractères du thème}

Maîtriser l'écriture des caractères clés du thème (arts, opinions, connecteurs) est indispensable pour l'expression écrite. Nous travaillons ici l'ordre des traits, les radicaux (clés) et les confusions fréquentes.

\courssub{Analyse de caractères clés : radicaux et ordre des traits}

\begin{center}
\begin{tabularx}{\linewidth}{|X|X|X|X|X|}\hline
\rowcolor{popblueL} Caractère & Pinyin & Radical (Clé) & Nb traits & Ordre des traits (description) \\ \hline
艺 & yì & \zh{艹} (cǎo, herbe) & 7 & Horizontal (herbe), vertical, horizontal (herbe) ; point, horizontal, crochet vertical, horizontal. \\ \hline
术 & shù & \zh{术} (shù, art/méthode) & 6 & Horizontal, vertical, horizontal replié, point, horizontal, crochet vertical. \\ \hline
书 & shū & \zh{乙} (yǐ, second) & 4 & Horizontal, crochet vertical, horizontal, crochet final (comme \zh{乙} + trait). \\ \hline
法 & fǎ & \zh{氵} (sān diǎn shuǐ, eau) & 8 & Trois points d'eau ; horizontal, vertical, horizontal replié, point, horizontal. \\ \hline
雕 & diāo & \zh{木} (mù, bois) & 16 & Bois (horizontal, vertical, trait gauche, trait droit) ; haut de \zh{周} (土 + 口 + 土), puis bas. \\ \hline
历 & lì & \zh{止} (zhǐ, arrêter) & 5 & \zh{厂} (han, falaise) : horizontal, vertical replié ; \zh{止} : horizontal, vertical, horizontal, crochet. \\ \hline
史 & shǐ & \zh{口} (kǒu, bouche) & 5 & Bouche (carré) ; vertical, horizontal replié, horizontal, horizontal. \\ \hline
论 & lùn & \zh{讠} (yán, parole) & 7 & Parole (point, courbe) ; \zh{仑} (lún) : personne, vertical, horizontal, vertical, horizontal. \\ \hline
比 & bǐ & \zh{比} (bǐ, comparer) & 4 & Trait gauche, trait droit, horizontal, trait final montant. \\ \hline
\end{tabularx}
\end{center}

\courssub{Caractères proches et pièges visuels (辨析 biànxī)}

\begin{exemplebox}[Paires à ne pas confondre]
\begin{itemize}
\item \zh{艺} (yì, art, technique) \textbf{vs} \zh{草} (cǎo, herbe) : \zh{艺} a \zh{未} (wèi, pas encore) en bas ; \zh{草} a \zh{早} (zǎo, matin) en bas.
\item \zh{书} (shū, livre/écrire) \textbf{vs} \zh{画} (huà, peinture/dessiner) : \zh{书} = \zh{乙} + trait ; \zh{画} = \zh{聿} (yù, pinceau) + \zh{田} (champ) + \zh{一} + \zh{凵}.
\item \zh{法} (fǎ, loi/méthode) \textbf{vs} \zh{去} (qù, aller) : \zh{法} a l'eau \zh{氵} à gauche ; \zh{去} a \zh{土} (terre) en haut à droite.
\item \zh{历} (lì, calendrier/histoire) \textbf{vs} \zh{经} (jīng, traverser/classique) : \zh{历} a \zh{止} en bas ; \zh{经} a \zh{巾} (tissu) à gauche et \zh{工} + \zh{土} à droite.
\item \zh{论} (lùn, discuter/théorie) \textbf{vs} \zh{伦} (lún, ordre/éthique) : \zh{论} a la parole \zh{讠} ; \zh{伦} a la personne \zh{亻}.
\end{itemize}
\end{exemplebox}

\begin{experience}[Atelier d'écriture guidée]
\textbf{Matériel} : Papier à carreaux chinois (米字格 mǐzìgé) ou feuilles quadrillées.\\
\textbf{Consigne} :
\begin{enumerate}
\item Écris 5 fois chaque caractère du tableau ci-dessus en respectant l'ordre des traits et la proportion dans le carré.
\item Pour chaque caractère, identifie son radical et donne un autre caractère partageant ce radical (ex: \zh{艺} $\rightarrow$ \zh{草、花、菜}).
\item Échange ta feuille avec un camarade : il doit corriger l'ordre des traits (numéroter les traits 1, 2, 3... sur un exemplaire) et entourer les erreurs de proportion.
\end{enumerate}
\end{experience}

\coursec{Thème et version guidés}

La traduction (thème : français $\rightarrow$ chinois ; version : chinois $\rightarrow$ français) permet de fixer les structures et le vocabulaire. La méthode est rigoureuse : analyser la phrase française, identifier la structure chinoise cible, choisir le vocabulaire, écrire les caractères, vérifier l'ordre des mots et la ponctuation.

\begin{methode}[Méthode de traduction Thème / Version]
\begin{enumerate}
\item \textbf{Analyse} : Repère le sujet, le verbe, les compléments. Identifie la structure grammaticale chinoise correspondante (Sujet + \zh{比} + Objet + Adj, \zh{因为}...\zh{所以}..., connecteurs en tête de phrase).
\item \textbf{Transposition} : Ne traduis pas mot à mot. Pense en « blocs chinois » (ex: « pour moi » $\rightarrow$ \zh{对我来说} ; « d'abord » $\rightarrow$ \zh{首先}).
\item \textbf{Rédaction} : Écris en caractères simplifiés. Ajoute le pinyin au brouillon si nécessaire.
\item \textbf{Vérification} : Ordre des mots (Sujet - Temps - Lieu - Verbe - Complément) ? Ponctuation chinoise ? Tons ? Caractères corrects (pas de traditionnels) ?
\end{enumerate}
\end{methode}

\courssub{Thème guidé (Français $\rightarrow$ Chinois)}

\begin{exoresolu}[Thème n°1]
\textbf{Énoncé} : « Le Cameroun et la Chine ont tous deux un art très célèbre. Pour moi, je préfère l'art camerounais. D'abord, la sculpture sur bois est très unique. Ensuite, la musique est très belle à écouter. Donc, j'aime l'art camerounais. Et toi ? »
\tcblower
\textbf{Analyse structurelle} :
\begin{itemize}
\item Phrase 1 : Introduction, comparaison deux pays $\rightarrow$ \zh{A和B都有...}
\item Phrase 2 : Opinion personnelle $\rightarrow$ \zh{对我来说，我更喜欢...}
\item Phrase 3 : Argument 1 $\rightarrow$ \zh{首先，...很独特}
\item Phrase 4 : Argument 2 $\rightarrow$ \zh{然后，...很好听}
\item Phrase 5 : Conclusion $\rightarrow$ \zh{所以，我喜欢... 你呢？}
\end{itemize}
\textbf{Proposition} :
\zh{喀麦隆和中国都有很有名的艺术。对我来说，我更喜欢喀麦隆艺术。首先，木雕很独特。然后，音乐很好听。所以，我喜欢喀麦隆艺术。你呢？}\\
\emph{Kāmàilóng hé Zhōngguó dōu yǒu hěn yǒumíng de yìshù. Duì wǒ láishuō, wǒ gèng xǐhuan Kāmàilóng yìshù. Shǒuxiān, mùdiāo hěn dútè. Ránhòu, yīnyuè hěn hǎotīng. Suǒyǐ, wǒ xǐhuan Kāmàilóng yìshù. Nǐ ne?}
\end{exoresolu}

\begin{exercice}[Thème : à vous de jouer]
Traduis les phrases suivantes en chinois (caractères + pinyin).
1. Je pense que la calligraphie chinoise est plus difficile que la peinture.
2. Parce que l'art camerounais a une longue histoire, il est très intéressant.
3. D'abord, j'aime la danse camerounaise. Ensuite, j'aime la musique. Donc, j'aime la culture camerounaise.
4. Les couleurs de cette peinture sont très belles. De plus, elle a plusieurs centaines d'années d'histoire.
\end{exercice}

\begin{corrige}[Exercice Thème]
1. \zh{我觉得中国书法比画难。} (Wǒ juéde Zhōngguó shūfǎ bǐ huà nán.)\\
2. \zh{因为喀麦隆艺术有很长的历史，所以很有意思。} (Yīnwèi Kāmàilóng yìshù yǒu hěn cháng de lìshǐ, suǒyǐ hěn yǒuyìsi.)\\
3. \zh{首先，我喜欢喀麦隆的舞蹈。然后，我喜欢音乐。所以，我喜欢喀麦隆文化。} (Shǒuxiān, wǒ xǐhuan Kāmàilóng de wǔdǎo. Ránhòu, wǒ xǐhuan yīnyuè. Suǒyǐ, wǒ xǐhuan Kāmàilóng wénhuà.)\\
4. \zh{这幅画的颜色很好看。还有，它有几百年的历史。} (Zhè fú huà de yánsè hěn hǎokàn. Háiyǒu, tā yǒu jǐ bǎi nián de lìshǐ.) — Note : classifieur \zh{幅} (fú) pour tableau/peinture.
\end{corrige}

\courssub{Version guidée (Chinois $\rightarrow$ Français)}

\begin{exoresolu}[Version n°1]
\textbf{Énoncé} : \zh{喀麦隆的雕塑和中国的陶艺都很有名。我觉得中国陶艺更精彩。那些？首先，颜色很美。然后，有几千年的历史。所以，我欣赏中国艺术。你呢？}
\tcblower
\textbf{Analyse} :
\begin{itemize}
\item \zh{喀麦隆的雕塑和中国的陶艺} : Sujets coordonnés (sculpture CM / céramique Chine).
\item \zh{都很有名} : \zh{都} + adjectif.
\item \zh{我觉得...更精彩} : Opinion + comparatif \zh{更}.
\item \zh{那些？} : Transition.
\item \zh{首先/然后/所以} : Connecteurs standards.
\item \zh{欣赏} : Verbe « apprécier » (niveau HSK 3/4).
\end{itemize}
\textbf{Traduction} : « La sculpture du Cameroun et la céramique de Chine sont toutes deux très célèbres. Je trouve la céramique chinoise encore plus remarquable. Pourquoi donc ? D'abord, les couleurs sont très belles. Ensuite, elle a plusieurs millénaires d'histoire. Donc, j'apprécie l'art chinois. Et toi ? »
\end{exoresolu}

\begin{exercice}[Version : à vous de jouer]
Traduis en français les phrases suivantes.
1. \zh{对我来说，喀麦隆的音乐比中国音乐更有名。}
2. \zh{因为木雕又美又独特，所以我喜欢。}
3. \zh{首先，书法像画。然后，颜色很漂亮。还有，历史很长。}
4. \zh{所以，我认为喀麦隆艺术很精彩。你呢？}
\end{exercice}

\begin{corrige}[Exercice Version]
1. « Pour moi, la musique camerounaise est plus célèbre que la musique chinoise. »
2. « Parce que la sculpture sur bois est à la fois belle et unique, je l'aime. »
3. « D'abord, la calligraphie ressemble à une peinture. Ensuite, les couleurs sont très jolies. De plus, l'histoire est très longue. »
4. « Donc, je pense que l'art camerounais est remarquable. Et toi ? »
\end{corrige}

\coursec{Production écrite et évaluation}

C'est l'étape finale : tu vas produire ton propre texte en t'appuyant sur tout ce qui a été vu.

\begin{experience}[Production écrite : « Mon art préféré »]
\textbf{Consigne} : Rédige un court texte (environ 60--80 caractères chinois) pour répondre à la question : \zh{你喜欢喀麦隆还是中国艺术？为什么？}
\textbf{Contraintes} :
\begin{itemize}
\item Respecte le plan en 4 parties (Introduction, Opinion + \zh{那些？}, 2--3 arguments avec \zh{首先/然后/还有}, Conclusion avec \zh{所以} + \zh{你呢？}).
\item Utilise au moins 5 mots du lexique de la leçon (arts, adjectifs, connecteurs).
\item Écris les caractères correctement (soigne l'ordre des traits).
\item Ajoute le pinyin au-dessus ou à côté si tu n'es pas sûr de ton écriture.
\end{itemize}
\textbf{Grille d'auto-évaluation (coche les cases)} :
\begin{itemize}
\item[$\square$] Mon introduction présente les deux pays/arts sans opinion (\zh{都有...}).
\item[$\square$] Mon opinion est claire (\zh{我更喜欢...} ou \zh{我觉得...}).
\item[$\square$] J'ai utilisé \zh{那些？} comme transition.
\item[$\square$] J'ai au moins 2 arguments distincts, chacun introduit par un connecteur (\zh{首先}, \zh{然后}, \zh{还有}).
\item[$\square$] Ma conclusion utilise \zh{所以} et finit par \zh{你呢？}.
\item[$\square$] Pas de faute de ponctuation chinoise (，。？).
\item[$\square$] Les caractères sont lisibles, traits dans l'ordre.
\end{itemize}
\end{experience}

\begin{exercice}[Évaluation formative : Expression écrite notée]
\textbf{Sujet} : « Compare la musique et la danse au Cameroun et en Chine. Laquelle préfères-tu ? Donne deux raisons. »
\textbf{Format} : Texte suivi (pas de liste), 5--6 phrases minimum. Caractères + pinyin + traduction française.
\textbf{Durée} : 20 minutes.
\textbf{Barème indicatif (10 pts)} :
\begin{itemize}
\item Respect du plan et connecteurs (3 pts)
\item Vocabulaire précis et varié (arts, adjectifs) (2 pts)
\item Grammaire correcte (\zh{比}, \zh{更}, \zh{因为...所以...}, ordre des mots) (3 pts)
\item Écriture des caractères (propreté, ordre traits) + Pinyin correct (2 pts)
\end{itemize}
\end{exercice}

\begin{aretenir}[Synthèse de la leçon ZHA22]
\begin{itemize}
\item \textbf{Vocabulaire} : Arts (\zh{木雕, 书法, 画, 音乐, 舞蹈, 雕塑, 陶艺}), Opinions (\zh{喜欢, 觉得, 欣赏, 更喜欢}), Adjectifs (\zh{美, 独特, 精彩, 好听, 有意思}).
\item \textbf{Grammaire} : Comparatif \zh{比} / \zh{plus} ; Cause-conséquence \zh{因为...所以...} ; Structure \zh{又...又...}.
\item \textbf{Textuelle} : Plan argumentatif en 4 étapes (Intro $\rightarrow$ Opinion $\rightarrow$ Arguments $\rightarrow$ Conclusion). Connecteurs : \zh{首先, 然后, 那些？, 所以, 你呢？}.
\item \textbf{Écriture} : Radicaux clés \zh{艹, 木, 氵, 讠, 止, 口} ; Ordre des traits rigoureux ; Distinction caractères proches (\zh{艺/草, 书/画, 法/去, 历/经, 论/伦}).
\item \textbf{Culture} : Sculpture bois / Tissus (Cameroun) vs Calligraphie / Peinture encre (Chine). Respect mutuel et ouverture.
\end{itemize}
\end{aretenir}

\coursec{Écriture des caractères du thème}

Après avoir analysé la structure du texte argumentatif et les connecteurs logiques qui permettent d'articuler une opinion nuancée sur les arts chinois et camerounais, nous abordons maintenant une compétence fondamentale : l'écriture correcte des caractères du thème. En chinois, l'ordre des traits (\cle{笔顺 bǐshùn}) n'est pas une contrainte calligraphique arbitraire ; il conditionne la structure interne du caractère, sa lisibilité en écriture cursive, sa reconnaissance par les logiciels de saisie (IME) et, in fine, sa mémorisation à long terme. Un trait omis, inversé ou mal proportionné peut transformer un mot en un non-sens ou en un autre mot existant. Nous allons disséquer les sept caractères clés de cette leçon, identifier leurs radicaux (\cle{部首 bùshǒu}), maîtriser leur tracé trait par trait, et désamorcer les confusions visuelles les plus fréquentes chez les apprenants francophones.

\courssub{Vocabulaire cible : analyse radicale et structurelle}

Avant de tracer, observons les « briques » qui composent nos caractères. Le tableau ci-dessous récapitule pour chaque graphème son radical (clé de dictionnaire), son nombre total de traits, sa structure (composition spatiale) et son sens principal dans le contexte de notre thème « Arts et opinions ».

\begin{center}
\begin{tabularx}{\linewidth}{|X|X|X|X|X|}
\hline
\rowcolor{popblueL} \textbf{Caractère} & \textbf{Pinyin} & \textbf{Radical (Clé)} & \textbf{Nb traits} & \textbf{Structure / Sens clé} \\
\hline
艺 & yì & 艹 (cǎo, herbe) & 4 & Haut-bas : \zh{艹} + \zh{乙} — Art, talent, compétence (forme simplifiée de \zh{藝}) \\
\hline
术 & shù & 木 (mù, bois) & 5 & Haut-bas : forme simplifiée de \zh{術} — Technique, art, méthode \\
\hline
画 & huà & 田 (tián, champ) & 8 & Enclos : \zh{田} + traits intérieurs — Peinture, dessin, tableau \\
\hline
喜 & xǐ & 口 (kǒu, bouche) & 12 & Haut-bas : \zh{士} + \zh{口} + \zh{口} — Joie, aimer, apprécier \\
\hline
欢 & huān & 欠 (qiàn, souffle) & 22 & Gauche-droite : \zh{又} (main) + partie droite (contenant \zh{欠}) — Joie, réjouissance, \zh{喜欢} \\
\hline
觉 & jué & 见 (jiàn, voir) & 9 & Haut-bas : \zh{⺍} + \zh{冖} + \zh{见} — Sentir, percevoir, penser (\zh{觉得}) \\
\hline
得 & de / dé / děi & 彳 (chì, pas) & 11 & Gauche-droite : \zh{彳} + partie droite — Particule de complément de degré / obtenir / devoir \\
\hline
\end{tabularx}
\end{center}

\courssub{Décomposition trait par trait et ordre d'écriture}

Nous appliquons ici les règles universelles de l'ordre des traits : \textbf{de haut en bas} (\cle{从上到下}), \textbf{de gauche à droite} (\cle{从左到右}), \textbf{l'horizontale avant la verticale qui la coupe} (\cle{横竖先横}), \textbf{l'enveloppe avant le contenu} (\cle{先外后内}), \textbf{fermer l'enveloppe en dernier} (\cle{后封口}), \textbf{le trait central avant les ailes} (\cle{中间先两边}).

\begin{propriete}[Ordre des traits : 艺 \zh{yì} (4 traits)]
\textbf{Radical :} \zh{艹} (cǎozìtóu, herbe sur la tête). \textbf{Structure :} Haut-bas.
\textbf{Décomposition :}
1. \zh{一} (Héng, horizontale) — trait supérieur du radical \zh{艹}.
2. \zh{丨} (Shù, verticale) — trait central du radical \zh{艹}, traverse l'horizontale.
3. \zh{丨} (Shù, verticale) — trait droit descendant, amorce la partie basse (\zh{乙}).
4. \zh{乙} (Yǐ, crochet final) — trait courbe vers la gauche avec un petit crochet montant à la fin. \textit{Attention : ce n'est pas un trait droit, c'est une courbe unique.}
\end{propriete}

\begin{popfigure}
\legende{Grille 米字格 (mǐzìgé) pour 艺  — Ordre des traits 1 à 4}
\begin{tikzpicture}[scale=1.2, every node/.style={font=\small}]
% Grille 米字格
\draw[popline] (-1.5,-1.5) rectangle (1.5,1.5);
\draw[popline] (-1.5,0) -- (1.5,0);
\draw[popline] (0,-1.5) -- (0,1.5);
\draw[popline] (-1.5,1.5) -- (1.5,-1.5);
\draw[popline] (-1.5,-1.5) -- (1.5,1.5);
% Traits du caractère 艺
% Trait 1
\draw[popblue, line width=2pt, ->] (-0.8,0.8) -- (0.8,0.8) node[midway, above, popdark] {1};
% Trait 2
\draw[popblue, line width=2pt, ->] (0,0.8) -- (0,0.1) node[midway, right, popdark] {2};
% Trait 3
\draw[popblue, line width=2pt, ->] (0,0) -- (0,-0.9) node[midway, right, popdark] {3};
% Trait 4
\draw[popblue, line width=2pt, ->] (0,-0.9) .. controls (-0.6,-1.2) and (-0.8,-1.0) .. (-0.9,-0.7) node[near end, left, popdark] {4};
% Numéros repères
\node[popblue, circle, fill=popblue!10, draw, minimum size=5mm] at (-1.0, 1.0) {1};
\node[popblue, circle, fill=popblue!10, draw, minimum size=5mm] at (0.2, 0.4) {2};
\node[popblue, circle, fill=popblue!10, draw, minimum size=5mm] at (0.2, -0.4) {3};
\node[popblue, circle, fill=popblue!10, draw, minimum size=5mm] at (-1.0, -0.8) {4};
\end{tikzpicture}
\end{popfigure}

\begin{propriete}[Ordre des traits : 术 \zh{shù} (5 traits)]
\textbf{Radical :} \zh{木} (mù, bois) — \textit{la forme simplifiée 术 dérive de 術, le radical bois reste la clé de dictionnaire}.
\textbf{Structure :} Haut-bas.
\textbf{Décomposition :}
1. \zh{一} (Horizontale) — en haut.
2. \zh{丨} (Verticale) — descend du milieu de l'horizontale.
3. \zh{丿} (Piě, oblique gauche) — part du bas de la verticale vers le bas-gauche.
4. \zh{丶} (Diǎn, point) — en bas à gauche, après le piě.
5. \zh{丶} (Diǎn, point) — en bas à droite, \textbf{ce point final distingue 术 de 木}.
\end{propriete}

\begin{popfigure}
\legende{Grille 米字格 pour 术  — Ordre des traits 1 à 5}
\begin{tikzpicture}[scale=1.2, every node/.style={font=\small}]
\draw[popline] (-1.5,-1.5) rectangle (1.5,1.5);
\draw[popline] (-1.5,0) -- (1.5,0);
\draw[popline] (0,-1.5) -- (0,1.5);
\draw[popline] (-1.5,1.5) -- (1.5,-1.5);
\draw[popline] (-1.5,-1.5) -- (1.5,1.5);
% 1
\draw[popblue, line width=2pt, ->] (-0.8,1.0) -- (0.8,1.0) node[midway, above, popdark] {1};
% 2
\draw[popblue, line width=2pt, ->] (0,1.0) -- (0,-0.2) node[midway, right, popdark] {2};
% 3
\draw[popblue, line width=2pt, ->] (0,-0.2) -- (-0.8,-1.0) node[midway, left, popdark] {3};
% 4
\fill[popblue] (-0.6,-1.1) circle (2pt); \node[popblue] at (-1.0,-1.1) {4};
% 5
\fill[popblue] (0.6,-1.1) circle (2pt); \node[popblue] at (1.0,-1.1) {5};
% Repères
\node[popblue, circle, fill=popblue!10, draw, minimum size=5mm] at (-1.0, 1.2) {1};
\node[popblue, circle, fill=popblue!10, draw, minimum size=5mm] at (0.2, 0.4) {2};
\node[popblue, circle, fill=popblue!10, draw, minimum size=5mm] at (-1.0, -0.6) {3};
\node[popblue, circle, fill=popblue!10, draw, minimum size=5mm] at (-1.0, -1.2) {4};
\node[popblue, circle, fill=popblue!10, draw, minimum size=5mm] at (1.0, -1.2) {5};
\end{tikzpicture}
\end{popfigure}

\begin{attention}[Piège majeur : \zh{术} (shù) \textbf{≠} \zh{木} (mù)]
\zh{木} (bois, arbre) s'écrit en 4 traits : \zh{一} → \zh{丨} → \zh{丿} → \zh{丶} (le dernier point est à gauche).
\zh{术} (technique, art) s'écrit en 5 traits : \zh{一} → \zh{丨} → \zh{丿} → \zh{丶} (gauche) → \zh{丶} (\textbf{droite}).
Oublier le 5\textsuperscript{e} trait (le point en bas à droite) transforme « technique/art » en « bois/arbre ». \zh{艺术} (yìshù, art) deviendrait \zh{艺木} (sans sens).
\zh{中医} (médecine chinoise) utilise \zh{术} dans \zh{针灸术} (technique d'acupuncture). Ne l'écrivez jamais avec \zh{木}.
\end{attention}

\begin{propriete}[Ordre des traits : 画 \zh{huà} (8 traits)]
\textbf{Radical :} \zh{田} (tián, champ). \textbf{Structure :} Enclos (le champ enferme les traits intérieurs).
\textbf{Décomposition (Norme GB/T 13000.1) :}
1. \zh{一} (Horizontale haute du \zh{田}).
2. \zh{丨} (Verticale gauche du \zh{田}).
3. \zh{一} (Horizontale médiane du \zh{田} — traverse l'enclos).
4. \zh{丨} (Verticale centrale — traverse l'horizontale médiane, s'arrête avant le bas).
5. \zh{(trait brisé)} (Zhé : horizontale basse + verticale droite du \zh{田} en un trait continu — \textbf{ferme l'enclos}).
6. \zh{一} (Horizontale intérieure haute — « manche du pinceau »).
7. \zh{丨} (Verticale intérieure basse — « poils du pinceau »).
8. \zh{丶} (Point final — bout du pinceau).
\end{propriete}

\begin{popfigure}
\legende{Grille 米字格 pour 画  — Ordre des traits 1 à 8}
\begin{tikzpicture}[scale=1.1, every node/.style={font=\tiny}]
\draw[popline] (-1.8,-1.8) rectangle (1.8,1.8);
\draw[popline] (-1.8,0) -- (1.8,0);
\draw[popline] (0,-1.8) -- (0,1.8);
\draw[popline] (-1.8,1.8) -- (1.8,-1.8);
\draw[popline] (-1.8,-1.8) -- (1.8,1.8);
% 1 Top horiz
\draw[popblue, line width=1.5pt, ->] (-1.2,1.2) -- (1.2,1.2) node[midway, above, popdark] {1};
% 2 Left vert
\draw[popblue, line width=1.5pt, ->] (-1.2,1.2) -- (-1.2,-0.6) node[midway, left, popdark] {2};
% 3 Middle horiz
\draw[popblue, line width=1.5pt, ->] (-1.2,0) -- (1.2,0) node[midway, above, popdark] {3};
% 4 Middle vert
\draw[popblue, line width=1.5pt, ->] (0,0.6) -- (0,-0.6) node[midway, right, popdark] {4};
% 5 Bottom + Right (Hook)
\draw[popblue, line width=1.5pt, ->] (-1.2,-1.2) -- (1.2,-1.2) -- (1.2,-0.6) node[midway, right, popdark] {5};
% 6 Inner top horiz
\draw[popblue, line width=1.5pt, ->] (-0.6,0.6) -- (0.6,0.6) node[midway, above, popdark] {6};
% 7 Inner vert
\draw[popblue, line width=1.5pt, ->] (0,0.3) -- (0,-0.9) node[midway, right, popdark] {7};
% 8 Dot
\fill[popblue] (0.1,-0.9) circle (1.5pt); \node[popblue] at (0.5, -1.1) {8};
\end{tikzpicture}
\end{popfigure}

\begin{attention}[Confusion \zh{画} (huà) vs \zh{面} (miàn)]
\begin{center}
\begin{tabularx}{\linewidth}{|X|X|X|}
\hline
\rowcolor{poporangeL} \textbf{Caractère} & \textbf{Structure clé} & \textbf{Exemple discriminant} \\
\hline
\zh{画} (huà) & Contient \zh{田} \textbf{FERMÉ} en bas (trait 5 : (trait brisé)). À l'intérieur : 2 traits (横 + 竖) + 1 point. & \zh{画画} (huàhuà, peindre), \zh{画报} (huàbào, magazine illustré). \\
\hline
\zh{面} (miàn) & Contient \zh{田} \textbf{OUVERT} en bas (pas de trait de fermeture). Au-dessus : 3 traits (\zh{一} + \zh{丨} + \zh{一}). & \zh{面条} (miàntiáo, nouilles), \zh{面具} (miànjù, masque), \zh{见面} (jiànmiàn, se rencontrer). \\
\hline
\end{tabularx}
\end{center}
\textbf{Astuce mnémotechnique :} Dans \zh{面具} (masque), le visage (\zh{面}) est « ouvert » pour respirer. Dans \zh{画} (peinture), le cadre (\zh{田}) est « fermé » pour tenir la toile.
\end{attention}

\begin{propriete}[Ordre des traits : 喜 \zh{xǐ} (12 traits)]
\textbf{Radical :} \zh{口} (kǒu, bouche) — \textit{en bas}. \textbf{Structure :} Haut-bas (\zh{士} + \zh{口} + \zh{口}).
\textbf{Décomposition (Haut → Bas) :}
\textit{Partie haute \zh{士} (shì, lettré) :}
1. \zh{一} (Horizontale courte, haut).
2. \zh{一} (Horizontale longue, bas du \zh{士}).
3. \zh{丨} (Verticale centrale — ne traverse pas les bouches).
\textit{Première bouche \zh{口} (traits 4–6) :}
4. \zh{丨} (Verticale gauche).
5. \zh{(trait brisé)} (Horizontale haute + verticale droite en un trait).
6. \zh{一} (Horizontale basse).
\textit{Deuxième bouche \zh{口} (traits 7–9) :}
7. \zh{丨} (Verticale gauche).
8. \zh{(trait brisé)} (Horizontale haute + verticale droite).
9. \zh{一} (Horizontale basse).
\textit{Note :} Les deux \zh{口} sont alignés horizontalement, de même largeur. Le trait 3 s'arrête au-dessus du trait 4.
\end{propriete}

\begin{popfigure}
\legende{Grille 米字格 pour 喜  — Structure haut-bas (12 traits)}
\begin{tikzpicture}[scale=1.0, every node/.style={font=\tiny}]
\draw[popline] (-1.8,-1.8) rectangle (1.8,1.8);
\draw[popline] (-1.8,0) -- (1.8,0);
\draw[popline] (0,-1.8) -- (0,1.8);
\draw[popline] (-1.8,1.8) -- (1.8,-1.8);
\draw[popline] (-1.8,-1.8) -- (1.8,1.8);
% Shi (top)
\draw[popblue, line width=1.2pt] (-0.6,1.2) -- (0.6,1.2); % 1
\draw[popblue, line width=1.2pt] (-1.0,0.6) -- (1.0,0.6); % 2
\draw[popblue, line width=1.2pt] (0,1.2) -- (0,0.6); % 3
% Kou 1 (left bottom)
\draw[popblue, line width=1.2pt] (-1.0,0) -- (-1.0,-0.6); % 4
\draw[popblue, line width=1.2pt] (-1.0,0) -- (0,0) -- (0,-0.6); % 5
\draw[popblue, line width=1.2pt] (-1.0,-0.6) -- (0,-0.6); % 6
% Kou 2 (right bottom)
\draw[popblue, line width=1.2pt] (0,0) -- (0,-0.6); % 7 (shared left line? No, separate box) -> (1.0,0) -- (1.0,-0.6)
\draw[popblue, line width=1.2pt] (0,0) -- (1.0,0) -- (1.0,-0.6); % 8
\draw[popblue, line width=1.2pt] (0,-0.6) -- (1.0,-0.6); % 9
% Numbers approx
\node[popblue] at (-1.5, 1.5) {1,2,3};
\node[popblue] at (-1.5, -0.3) {4-6};
\node[popblue] at (1.5, -0.3) {7-9};
\end{tikzpicture}
\end{popfigure}

\begin{propriete}[Ordre des traits : 欢 \zh{huān} (22 traits — forme simplifiée)]
\textbf{Radical :} \zh{欠} (qiàn, souffle/manque). \textbf{Structure :} Gauche-Droite.
\textbf{Règle :} \textbf{Partie gauche avant partie droite} (\cle{先左后右}).
\textbf{Partie gauche (2 traits) :} \zh{又} (yòu, main droite).
1. \zh{丿} (Piě, oblique gauche).
2. \zh{㇏} (Nà, trait plat droit — \textbf{sans crochet}).
\textbf{Partie droite (20 traits, de haut en bas) :} Structure complexe contenant le radical \zh{欠} en bas (4 traits : \zh{一}, \zh{丨}, \zh{丿}, \zh{丶}).
\textit{Note pédagogique :} Au lycée (HSK 3-4), on mémorise \zh{欢} comme un bloc « joie » (\zh{喜欢}). On s'assure d'écrire \zh{又} correctement (piě puis nà, pas de crochet) et le radical \zh{欠} en bas à droite (horizontal, verticale, oblique, point).
\end{propriete}

\begin{popfigure}
\legende{Structure Gauche-Droite de 欢 : 又 (gauche) + Partie droite (droite, radical 欠 en bas)}
\begin{tikzpicture}[scale=1.2, every node/.style={font=\small}]
\draw[popline] (-3,-1.5) rectangle (3,1.5);
\draw[popline] (0,-1.5) -- (0,1.5); % Separation
% Left: You
\node[popdark, align=center] at (-1.5, 0) {\textbf{Gauche (2 traits)}\\\zh{又}\\(yòu)\\\zh{丿} + \zh{㇏}};
% Right: Complex
\node[popdark, align=center] at (1.5, 0) {\textbf{Droite (20 traits)}\\\textit{Partie phonétique haute}\\\zh{欠} (radical, bas, 4 traits)};
% Draw You inside left box roughly
\draw[popblue, thick, ->] (-2.0, 0.5) -- (-1.2, -0.2) node[midway, above left] {1};
\draw[popblue, thick, ->] (-1.8, -0.2) -- (-1.0, -0.8) node[midway, below right] {2};
% Draw Qian inside right box roughly
\draw[popblue, thick, ->] (1.0, -0.5) -- (2.0, -0.5) node[midway, above] {17};
\draw[popblue, thick, ->] (1.5, -0.5) -- (1.5, -1.0) node[midway, right] {18};
\draw[popblue, thick, ->] (1.5, -1.0) -- (1.0, -1.3) node[midway, left] {19};
\fill[popblue] (2.0, -1.3) circle (1.5pt); \node[popblue] at (2.3, -1.3) {20};
\end{tikzpicture}
\end{popfigure}

\begin{attention}[Confusion \zh{欢} (huān) vs \zh{吹} (chuī)]
\begin{center}
\begin{tabularx}{\linewidth}{|X|X|X|}
\hline
\rowcolor{poporangeL} \textbf{Caractère} & \textbf{Différence graphique} & \textbf{Contexte} \\
\hline
\zh{欢} (huān) & Partie gauche : \zh{又} (main droite, traits : \zh{丿} + \zh{㇏}). Radical \zh{欠} en bas à droite. & \zh{喜欢} (xǐhuan, aimer), \zh{欢迎} (huānyíng, bienvenue). \\
\hline
\zh{吹} (chuī) & Partie gauche : \zh{口} (bouche, carré fermé). Radical \zh{欠} en bas à droite. & \zh{吹风} (chuīfēng, vent / sèche-cheveux), \zh{吹牛} (chuīniú, se vanter). \\
\hline
\end{tabularx}
\end{center}
\textbf{Règle d'or :} \zh{口} (bouche) est un carré \zh{口} ; \zh{又} (main) est ouvert \zh{又}. « Souffler » (\zh{吹}) passe par la \textbf{bouche} (\zh{口}). « Se réjouir » (\zh{欢}) vient du \textbf{cœur/main} (\zh{又} + \zh{欠} souffle de joie).
\end{attention}

\begin{propriete}[Ordre des traits : 觉 \zh{jué} (9 traits — forme simplifiée)]
\textbf{Radical :} \zh{见} (jiàn, voir). \textbf{Structure :} Haut-bas.
\textbf{Partie haute (3 traits) :}
1. \zh{一} (Horizontale courte, haut gauche).
2. \zh{一} (Horizontale courte, milieu gauche).
3. \zh{冖} (Mì, couvercle : horizontale large avec deux petites pointes descendantes aux extrémités — \textit{un seul trait continu}).
\textbf{Partie basse \zh{见} (6 traits dans ce contexte) :}
4. \zh{一} (Horizontale de \zh{见}).
5. \zh{丨} (Verticale).
6. \zh{(trait brisé)} (Angle : horizontale + crochet bas).
7. \zh{一} (Horizontale).
8. \zh{丨} (Verticale centrale descendante).
9. \zh{一} (Horizontale finale).
\end{propriete}

\begin{popfigure}
\legende{Grille 米字格 pour 觉  — Ordre des traits 1 à 9}
\begin{tikzpicture}[scale=1.1, every node/.style={font=\tiny}]
\draw[popline] (-1.8,-1.8) rectangle (1.8,1.8);
\draw[popline] (-1.8,0) -- (1.8,0);
\draw[popline] (0,-1.8) -- (0,1.8);
\draw[popline] (-1.8,1.8) -- (1.8,-1.8);
\draw[popline] (-1.8,-1.8) -- (1.8,1.8);
% Top part (3 strokes)
\draw[popblue, line width=1.5pt, ->] (-0.8,1.2) -- (-0.2,1.2) node[midway, above, popdark] {1};
\draw[popblue, line width=1.5pt, ->] (-0.8,0.7) -- (-0.2,0.7) node[midway, above, popdark] {2};
\draw[popblue, line width=1.5pt, ->] (-1.2,0.2) -- (1.2,0.2) .. controls (1.4,0.1) and (1.4,-0.1) .. (1.2,-0.2) node[midway, above, popdark] {3}; % Flattened mì
% Bottom part Jian (6 strokes)
\draw[popblue, line width=1.5pt, ->] (-1.0,-0.2) -- (1.0,-0.2) node[midway, above, popdark] {4};
\draw[popblue, line width=1.5pt, ->] (0,-0.2) -- (0,-0.8) node[midway, right, popdark] {5};
\draw[popblue, line width=1.5pt, ->] (-1.0,-0.8) -- (1.0,-0.8) -- (1.0,-1.2) node[midway, right, popdark] {6};
\draw[popblue, line width=1.5pt, ->] (-1.0,-1.2) -- (1.0,-1.2) node[midway, above, popdark] {7};
\draw[popblue, line width=1.5pt, ->] (0,-1.2) -- (0,-1.6) node[midway, right, popdark] {8};
\draw[popblue, line width=1.5pt, ->] (-0.8,-1.6) -- (0.8,-1.6) node[midway, below, popdark] {9};
\end{tikzpicture}
\end{popfigure}

\begin{propriete}[Ordre des traits : 得 \zh{de/dé/děi} (11 traits)]
\textbf{Radical :} \zh{彳} (chì, double pas, gauche). \textbf{Structure :} Gauche-Droite.
\textbf{Règle :} \textbf{Radical gauche \zh{彳} entièrement avant la droite}.
\textbf{Partie gauche \zh{彳} (3 traits) :}
1. \zh{丿} (Piě).
2. \zh{丿} (Piě).
3. \zh{丨} (Shù).
\textbf{Partie droite (8 traits) :}
4. \zh{一} (Horizontale).
5. \zh{丨} (Verticale).
6. \zh{(trait brisé)} (Angle : horizontale + crochet bas).
7. \zh{一} (Horizontale).
8. \zh{一} (Horizontale).
9. \zh{丨} (Verticale).
10. \zh{丿} (Piě).
11. \zh{㇏} (Nà).
\end{propriete}

\begin{popfigure}
\legende{Grille 米字格 pour 得  — Gauche 彳 (3 traits) + Droite (8 traits)}
\begin{tikzpicture}[scale=1.1, every node/.style={font=\tiny}]
\draw[popline] (-1.8,-1.8) rectangle (1.8,1.8);
\draw[popline] (-1.8,0) -- (1.8,0);
\draw[popline] (0,-1.8) -- (0,1.8);
\draw[popline] (-1.8,1.8) -- (1.8,-1.8);
\draw[popline] (-1.8,-1.8) -- (1.8,1.8);
% Left: Chi (3 strokes)
\draw[popblue, line width=1.5pt, ->] (-1.2,1.0) -- (-1.5,0.5) node[midway, left, popdark] {1};
\draw[popblue, line width=1.5pt, ->] (-1.2,0.2) -- (-1.5,-0.3) node[midway, left, popdark] {2};
\draw[popblue, line width=1.5pt, ->] (-1.35,-0.3) -- (-1.35,-1.0) node[midway, left, popdark] {3};
% Right: (8 strokes) - simplified representation
\draw[popblue, line width=1.5pt, ->] (0.2,1.0) -- (1.4,1.0) node[midway, above, popdark] {4};
\draw[popblue, line width=1.5pt, ->] (0.8,1.0) -- (0.8,0.2) node[midway, right, popdark] {5};
\draw[popblue, line width=1.5pt, ->] (0.2,0.2) -- (1.4,0.2) -- (1.4,-0.2) node[midway, right, popdark] {6};
\draw[popblue, line width=1.5pt, ->] (0.2,-0.2) -- (1.4,-0.2) node[midway, above, popdark] {7};
\draw[popblue, line width=1.5pt, ->] (0.2,-0.6) -- (1.4,-0.6) node[midway, above, popdark] {8};
\draw[popblue, line width=1.5pt, ->] (0.8,-0.6) -- (0.8,-1.2) node[midway, right, popdark] {9};
\draw[popblue, line width=1.5pt, ->] (0.8,-1.2) -- (0.3,-1.5) node[midway, left, popdark] {10};
\draw[popblue, line width=1.5pt, ->] (0.8,-1.2) -- (1.3,-1.5) node[midway, right, popdark] {11};
\end{tikzpicture}
\end{popfigure}

\courssub{Caractères proches : exercice de discrimination visuelle}

Les erreurs les plus coûteuses aux examens (et dans la vie réelle) viennent de caractères qui se ressemblent comme des frères jumeaux mais dont le sens et la prononciation divergent radicalement. Voici les cinq paires critiques de cette leçon.

\begin{exemplebox}[Autres paires à surveiller : \zh{觉} vs \zh{党} ; \zh{美} vs \zh{关}]
\begin{itemize}
\item \zh{觉} (jué, sentir) vs \zh{党} (dǎng, parti) : \zh{觉} a pour radical \zh{见} (voir) en bas. \zh{党} a pour radical \zh{儿} (enfant/pied) en bas + \zh{尚} en haut. \zh{觉得} (juéde) vs \zh{党派} (dǎngpài).
\item \zh{美} (měi, beau) vs \zh{关} (guān, fermer/concerner) : \zh{美} = \zh{羊} (mouton, 6 traits) + \zh{大} (grand, 3 traits). \zh{关} = \zh{门} (porte, radical) + \zh{关} (phonétique). \zh{美术} (měishù) vs \zh{关心} (guānxīn).
\end{itemize}
\end{exemplebox}

\courssub{Méthodologie : comment ancrer l'écriture dans la mémoire}

\begin{methode}[Mémoriser un caractère de manière durable (5 étapes)]
\begin{enumerate}
\item \textbf{Identifier le radical (部首 bùshǒu)} : C'est la « clé » du dictionnaire et l'indice sémantique. Ex: \zh{艹} pour \zh{艺} (plante/culture), \zh{欠} pour \zh{欢} (souffle/émotion).
\item \textbf{Compter et nommer les traits (笔画 bǐhuà)} : Épeler l'ordre à voix haute : « 一, 丨, 丨, 乙 » pour \zh{艺}. Cela active la mémoire kinesthésique et auditive.
\item \textbf{Écrire 5 fois dans le 米字格 (mǐzìgé)} : Respecter scrupuleusement l'ordre, la proportion (centrer le caractère) et la direction des traits. Ne pas « dessiner », mais « tracer ».
\item \textbf{Écrire le caractère dans un mot (词语 cíyǔ)} : \zh{艺术}, \zh{书法}, \zh{喜欢}, \zh{觉得}. Le contexte stabilise la forme.
\item \textbf{Produire une phrase personnelle (造句 zàojù)} : \zh{我喜欢中国画。} (Wǒ xǐhuan Zhōngguó huà.) L'usage personnel scelle l'acquisition.
\end{enumerate}
\end{methode}

\begin{savaistu}[L'étymologie vivante de \zh{艺} (yì) : l'art comme culture]
Le caractère traditionnel est \zh{藝}. En haut, \zh{艹} (plantes) ; au milieu, \zh{云} (nuages/vapeur) ; en bas, \zh{埶} (force/momentum, avec \zh{土} terre). L'ensemble évoque \textbf{une personne agenouillée plantant des végétaux avec soin} : l'art originel, c'est l'agriculture, la culture de la terre (\zh{种艺}).
La forme simplifiée \zh{艺} conserve \zh{艹} (herbe) et \zh{乙} (yǐ, deuxième tige céleste, évoquant une pousse qui émerge). \textbf{L'art (艺) est donc, étymologiquement, le fait de faire pousser le beau comme on fait pousser le riz.} Au Cameroun, on dit « \emph{La terre ne ment pas} » ; en Chine, « \emph{艺不压身} » (yì bù yā shēn — un talent ne pèse jamais sur son porteur, il l'enrichit). Cette vision utilitaire et vitale de l'art rapproche nos deux cultures.
\end{savaistu}

\courssub{Exercice résolu : écriture guidée et correction d'erreurs}

\begin{exoresolu}[Rédaction d'une mini-biographie artistique]
\textbf{Consigne :} Complétez le texte suivant en écrivant les caractères manquants (en gras dans la version pinyin) dans les cases vides. Respectez l'ordre des traits. Traduisez ensuite la phrase en français.

\begin{center}
\begin{tabularx}{\linewidth}{|X|X|}
\hline
\zh{我 \underline{\hspace{2.5cm}} \underline{\hspace{2.5cm}} 中国 \underline{\hspace{2.5cm}} 。} & \zh{Wǒ \textbf{xǐhuan} Zhōngguó \textbf{huà} .} \\
\hline
\zh{这 \underline{\hspace{2.5cm}} \underline{\hspace{2.5cm}} 很 \underline{\hspace{2.5cm}} \underline{\hspace{2.5cm}} 。} & \zh{Zhè \textbf{shì} \textbf{shù} hěn \textbf{yǒu} \textbf{yìsi} .} \\
\hline
\zh{你 \underline{\hspace{2.5cm}} \underline{\hspace{2.5cm}} 喀麦隆 \underline{\hspace{2.5cm}} \underline{\hspace{2.5cm}} 吗？} & \zh{Nǐ \textbf{juéde} Kāmàilóng \textbf{méi} \textbf{shù} \textbf{hǎo} ma?} \\
\hline
\end{tabularx}
\end{center}

\tcblower
\textbf{Solution détaillée :}

\begin{enumerate}
\item \zh{我 \textbf{喜欢} 中国 \textbf{画} 。} (Wǒ \textbf{xǐhuan} Zhōngguó \textbf{huà}.) — J'aime la peinture chinoise.
\begin{itemize}
\item \zh{喜} : 12 traits. Vérifier les deux \zh{口} alignés, le \zh{士} au-dessus ne touche pas le haut des bouches.
\item \zh{欢} : Gauche \zh{又} (piě + nà), droite radical \zh{欠}. Ne pas mettre \zh{口} à gauche (ce serait \zh{吹}).
\item \zh{画} : 8 traits. Fermer le \zh{田} (trait 5 : (trait brisé)) AVANT d'écrire les traits intérieurs (traits 6, 7, 8). Ne pas confondre avec \zh{面} (pas de fermeture bas, 3 traits au-dessus).
\end{itemize}

\item \zh{这 \textbf{是} \textbf{术} 很 \textbf{有} \textbf{意思} 。} (Zhè \textbf{shì} \textbf{shù} hěn \textbf{yǒu} \textbf{yìsi}.) — C'est une technique/art très intéressant.
\begin{itemize}
\item \zh{术} : 5 traits. \textbf{Vérifier le 5\textsuperscript{e} trait : le point en bas à droite (丶).} Sans lui, c'est \zh{木} (bois).
\item \zh{意思} (yìsi) : sens, intérêt. \zh{意} (radical \zh{心}, cœur) + \zh{思} (penser).
\item \textit{Note :} \zh{术} seul signifie « technique/méthode ». Pour « art » au sens large, on préfère \zh{艺术} (yìshù).
\end{itemize}

\item \zh{你 \textbf{觉得} 喀麦隆 \textbf{美术} \textbf{好} 吗？} (Nǐ \textbf{juéde} Kāmàilóng \textbf{méishù} \textbf{hǎo} ma?) — Penses-tu que les beaux-arts du Cameroun sont beaux/bons ?
\begin{itemize}
\item \zh{觉} : 9 traits. Haut : \zh{⺍} + \zh{冖}. Bas : \zh{见}. Ne pas confondre avec \zh{党} (radical \zh{儿} en bas).
\item \zh{美} : 9 traits. \zh{羊} (6 traits : \zh{丶}, \zh{丶}, \zh{丶}, \zh{丨}, \zh{(trait brisé)}, \zh{一}) + \zh{大} (3 traits : \zh{一}, \zh{丿}, \zh{㇏}). Ordre \zh{羊} : trois points, verticale, angle, horizontale basse. \zh{大} : horizontale, piě, nà.
\item \zh{术} : Rappel du point final (5\textsuperscript{e} trait).
\end{itemize}
\end{enumerate}
\end{exoresolu}

\courssub{Entraînement autonome : du tracé à la phrase}

\begin{exercice}[Devoir d'écriture — À rendre au prochain cours]
\begin{enumerate}
\item \textbf{Exercice de grille (米字格) :} Recopiez 5 fois chacun des 7 caractères de la leçon (\zh{艺, 术, 画, 喜, 欢, 觉, 得}) dans votre cahier d'écriture (format 米字格). Annotez le numéro de chaque trait au crayon à papier pour les 3 premiers exemplaires.
\item \textbf{Discrimination :} Entourez le caractère correct dans chaque paire :
\begin{itemize}
\item ( \zh{画} / \zh{面} ) \hspace{1cm} ( \zh{术} / \zh{木} ) \hspace{1cm} ( \zh{欢} / \zh{吹} ) \hspace{1cm} ( \zh{觉} / \zh{党} ) \hspace{1cm} ( \zh{美} / \zh{关} )
\end{itemize}
\item \textbf{Production écrite (3 phrases minimum) :} Répondez en chinois (caractères + pinyin) à la question du thème : \zh{你喜欢喀麦隆还是中国艺术？为什么？} (Nǐ xǐhuan Kāmàilóng háishi Zhōngguó yìshù? Wèishénme?)
\begin{itemize}
\item Utilisez au moins 5 des 7 caractères étudiés.
\item Utilisez au moins 2 connecteurs vus en Section 3 (\zh{因为/所以}, \zh{虽然/但是}, \zh{而且}, \zh{另外}).
\item Soignez la calligraphie : alignement, proportion, ordre des traits.
\end{itemize}
\end{enumerate}
\end{exercice}

\begin{corrige}[Exercice d'écriture — Éléments de correction]
\begin{enumerate}
\item \textbf{Grilles :} Vérification par l'enseignant : respect de l'ordre (haut→bas, gauche→droite, horizontale→verticale, enveloppe→contenu, fermeture finale), centrage dans le 米字格, proportion équilibrée (pas de trait qui « déborde » excessivement).
\item \textbf{Discrimination :} \zh{画} (peinture), \zh{术} (technique), \zh{欢} (joie), \zh{觉} (sentir), \zh{美} (beau).
\item \textbf{Production type (modèle) :}
\zh{我喜欢喀麦隆的艺术，也喜欢中国的艺术。} (Wǒ xǐhuan Kāmàilóng de yìshù, yě xǐhuan Zhōngguó de yìshù.) — J'aime l'art camerounais, et j'aime aussi l'art chinois.
\zh{因为喀麦隆的雕塑很有力量，中国的画很有意境。} (Yīnwèi Kāmàilóng de diāosù hěn yǒu lìliàng, Zhōngguó de huà hěn yǒu yìjìng.) — Car la sculpture camerounaise a de la force, la peinture chinoise a de l'atmosphère artistique.
\zh{我觉得两种艺术都很美，所以我都想学习。} (Wǒ juéde liǎng zhǒng yìshù dōu hěn měi, suǒyǐ wǒ dōu xiǎng xuéxí.) — Je trouve les deux arts très beaux, donc je veux tous les deux apprendre.
\begin{itemize}
\item \textbf{Vocabulaire bonus :} \zh{雕塑} (diāosù, sculpture), \zh{力量} (lìliàng, force), \zh{意境} (yìjìng, conception artistique/atmosphère), \zh{两种} (liǎng zhǒng, deux sortes).
\end{itemize}
\end{enumerate}
\end{corrige}

\coursec{Thème et version guidés}

Dans la section précédente, nous avons appris à écrire correctement les caractères du thème : \zh{艺} (art), \zh{术} (technique), \zh{画} (peinture), \zh{书} (livre/écrire), \zh{法} (loi/méthode), \zh{面} (visage/surface), \zh{具} (outil), en respectant l'ordre des traits et les radicaux. Nous allons maintenant utiliser ces caractères dans leur fonction première : \cle{traduire}. Traduire, c'est passer d'une langue à une autre en respectant deux choses : le \cle{sens} et la \cle{structure} de la langue d'arrivée. En chinois, on distingue deux exercices complémentaires : le \cle{thème} (français $\rightarrow$ chinois) et la \cle{version} (chinois $\rightarrow$ français). Ces deux exercices sont au cœur de l'évaluation de Première : ils vérifient à la fois votre grammaire, votre vocabulaire et votre écriture des caractères.

\courssub{Le thème : du français vers le chinois, étape par étape}

\begin{methode}[Méthode du thème en quatre étapes]
Pour traduire une phrase française en chinois, procédez toujours ainsi :
\begin{enumerate}
\item \textbf{Repérer la structure française} : qui fait l'action ? (sujet) Quelle action ? (verbe) Quoi ? (complément)
\item \textbf{Convertir en structure chinoise SVO} : Sujet + Verbe + Complément, sans article, sans conjugaison.
\item \textbf{Choisir les caractères appris} : n'inventez pas de mots ; utilisez uniquement le lexique de la leçon.
\item \textbf{Relire en vérifiant radicaux et traits} : un caractère mal écrit change le sens (\zh{术} art / \zh{木} arbre ; \zh{法} méthode / \zh{去} aller).
\end{enumerate}
\end{methode}

\textbf{Phrase 1 : « J'aime l'art chinois. »}

Étape 1 : sujet = \emph{je} ; verbe = \emph{aime} ; complément = \emph{l'art chinois}. Étape 2 : le chinois ne connaît ni article (\emph{l'}) ni conjugaison : 我 + 喜欢 + 中国艺术. Étape 3 : caractères connus. Étape 4 : vérification de \zh{喜} (composants 士 et 口) et \zh{欢} (clé 欠).

\begin{center}
\zh{我喜欢中国艺术。}\par
\emph{Wǒ xǐhuan Zhōngguó yìshù.}\par
« J'aime l'art chinois. »
\end{center}

\textbf{Phrase 2 : « Le Cameroun a de beaux masques. »}

Attention au verbe \emph{avoir} : ici il exprime la \cle{possession}, donc \zh{有} est correct. « De beaux masques » = adjectif + nom ; en chinois, l'adjectif attribut se place avant le nom, précédé de \zh{很} (voir la boîte Attention ci-dessous) : 很好看的面具.

\begin{center}
\zh{喀麦隆有很好看的面具。}\par
\emph{Kāmàilóng yǒu hěn hǎokàn de miànjù.}\par
« Le Cameroun a de beaux masques. »
\end{center}

\textbf{Phrase 3 : « Je trouve la calligraphie intéressante. »}

« Je trouve que... » se rend par \zh{我觉得} (je pense, je trouve). L'adjectif « intéressante » est ici \cle{attribut} (il décrit la calligraphie) : il faut donc \zh{很} devant \zh{有意思}.

\begin{center}
\zh{我觉得书法很有意思。}\par
\emph{Wǒ juéde shūfǎ hěn yǒu yìsi.}\par
« Je trouve la calligraphie intéressante. »
\end{center}

\textbf{Phrase 4 : « Je préfère la peinture chinoise parce que les couleurs sont belles. »}

« Préférer » = \zh{更喜欢} (aimer davantage). « Parce que » = \zh{因为}, qui introduit la cause. En chinois, la proposition de cause se place en général \emph{avant} la proposition principale, souvent avec \zh{所以} en miroir, mais on peut omettre 所以.

\begin{center}
\zh{我更喜欢中国画，因为颜色很好看。}\par
\emph{Wǒ gèng xǐhuan Zhōngguó huà, yīnwèi yánsè hěn hǎokàn.}\par
« Je préfère la peinture chinoise parce que les couleurs sont belles. »
\end{center}

\begin{exemplebox}[Exemple traduit et commenté]
\zh{因为中国艺术有几千年的历史，所以我更喜欢中国艺术。}\par
\emph{Yīnwèi Zhōngguó yìshù yǒu jǐ qiān nián de lìshǐ, suǒyǐ wǒ gèng xǐhuan Zhōngguó yìshù.}\par
« Parce que l'art chinois a plusieurs millénaires d'histoire, je le préfère. »\par
Observez : la paire \zh{因为}... \zh{所以}... (parce que... donc...) encadre les deux propositions ; la virgule chinoise \zh{，} les sépare ; \zh{有} exprime ici la possession (« l'art \emph{possède} une longue histoire ») ; \zh{的} relie la durée \zh{几千年} au nom \zh{历史}.
\end{exemplebox}

\begin{attention}[L'adjectif attribut exige 很]
En chinois, un adjectif seul comme attribut sonne \emph{incomplet} ou crée une comparaison implicite. On ajoute donc \zh{很} devant l'adjectif, \textbf{sans valeur intensive forte} : \zh{颜色很好看} (les couleurs sont belles) et non \zh{颜色好} (phrase qui sonne inachevée ou qui suggère une opposition : « ce sont les couleurs qui sont belles, pas autre chose »). De même : \zh{面具很好看} (les masques sont beaux), \zh{书法很有意思} (la calligraphie est intéressante). N'écrivez jamais une phrase attributive affirmative avec l'adjectif nu. Exception : dans une question avec \zh{吗} ou dans une négation, \zh{很} peut disparaître : \zh{面具好看吗？} (Les masques sont-ils beaux ?), \zh{颜色不好看} (les couleurs ne sont pas belles).
\end{attention}

\begin{attention}[Piège : « il y a » ne se traduit pas toujours par 有]
Le français « il y a » a deux emplois différents. \textbf{Possession / existence} : « Le Cameroun a de beaux masques » / « il y a de beaux masques » $\rightarrow$ \zh{有} est correct : \zh{喀麦隆有很好看的面具。} ; « il y a plusieurs millénaires d'histoire » au sens de « l'art \emph{possède} une longue histoire » $\rightarrow$ \zh{有几千年的历史} est correct. \textbf{Mais} « il y a + durée » au sens de « cela fait... que » (il y a deux ans que j'apprends le chinois) ne se traduit pas par \zh{有} : on dit \zh{我学中文学了两年了。} (Wǒ xué Zhōngwén xué le liǎng nián le.) Avant de poser \zh{有}, demandez-vous : s'agit-il de \emph{posséder} ou d'une \emph{durée écoulée} ?
\end{attention}

\courssub{Schéma de la structure de phrase}

\begin{popfigure}
\begin{tikzpicture}[node distance=0.35cm,
 bloc/.style={draw=popblue, fill=popblueL, rounded corners=3pt, minimum height=0.9cm, align=center, font=\small},
 plus/.style={font=\large\bfseries, text=poporange}]
\node[bloc, align=center] (s){Sujet\\ \zh{我} wǒ};
\node[plus, right=of s] (p1) {+};
\node[bloc, right=of p1, fill=poppinkL, draw=poppink, align=center] (adv){(很 / 更 / 也 / 都)\\ adverbe};
\node[plus, right=of adv] (p2) {+};
\node[bloc, right=of p2, fill=popgreenL, draw=popgreen, align=center] (v){Verbe ou\\ adjectif\\ \zh{更喜欢}};
\node[plus, right=of v] (p3) {+};
\node[bloc, right=of p3, fill=popgoldL, draw=popgold, align=center] (c){Complément\\ \zh{中国艺术}};
\node[below=0.5cm of v, align=center, font=\small\itshape, text=popmuted] {Exemple : \zh{我更喜欢中国艺术。} — Je préfère l'art chinois.};
\end{tikzpicture}
\legende{Structure de base : Sujet + (adverbe) + Verbe/Adjectif + Complément. Les adverbes 很, 更, 也, 都 se placent toujours AVANT le verbe ou l'adjectif.}
\end{popfigure}

\begin{propriete}[Ponctuation chinoise]
La virgule chinoise \zh{，} sépare les propositions à l'intérieur de la phrase ; le point \zh{。} termine la phrase ; le point d'interrogation \zh{？} suit les mêmes règles qu'en français (fin de question). Attention : la virgule chinoise est pleine chasse (plus large) et il n'y a \textbf{pas d'espace} avant la ponctuation, contrairement au français. Exemple : \zh{我觉得喀麦隆艺术很美，因为它有很长的历史。} (Wǒ juéde Kāmàilóng yìshù hěn měi, yīnwèi tā yǒu hěn cháng de lìshǐ. — Je trouve que l'art camerounais est très beau, parce qu'il a une très longue histoire.)
\end{propriete}

\begin{propriete}[Place de 也 et 都]
Les adverbes \zh{也} (aussi) et \zh{都} (tous) se placent \textbf{après le sujet et avant le verbe}, jamais en fin de phrase comme le français « aussi ». \zh{我也喜欢中国艺术。} (Wǒ yě xǐhuan Zhōngguó yìshù. — Moi aussi, j'aime l'art chinois.) \zh{我们都喜欢喀麦隆的音乐。} (Wǒmen dōu xǐhuan Kāmàilóng de yīnyuè. — Nous aimons tous la musique camerounaise.) Erreur fréquente du francophone : \zh{我喜欢中国艺术也} est incorrect.
\end{propriete}

\begin{propriete}[Le mot-outil 的 entre le déterminant et le nom]
Pour dire « l'art \emph{du Cameroun} », « la musique \emph{du Cameroun} », le chinois place le déterminant avant le nom, relié par \zh{的} : \zh{喀麦隆的艺术} (l'art du Cameroun), \zh{中国的音乐} (la musique de Chine), \zh{喀麦隆的舞} (la danse camerounaise). Schéma : Déterminant + \zh{的} + Nom. N'oubliez jamais \zh{的} entre un nom de pays et le nom qu'il détermine : \zh{喀麦隆艺术} (sans 的) est possible dans les titres, mais à l'écrit scolaire, préférez la forme complète avec \zh{的} quand le sens est « l'art du Cameroun ».
\end{propriete}

\courssub{La version : du chinois vers le français}

\begin{methode}[Stratégie de version]
\begin{enumerate}
\item \textbf{Identifier} dans chaque phrase chinoise le sujet, le verbe et le complément.
\item \textbf{Traduire par blocs de sens} (groupe sujet, groupe verbal, groupe complément), jamais caractère par caractère.
\item \textbf{Reconstituer} une phrase française correcte : articles, accords, conjugaison.
\item \textbf{Vérifier les tons et les caractères} : un ton mal lu change le mot (\zh{买} mǎi acheter / \zh{卖} mài vendre).
\end{enumerate}
\end{methode}

\begin{textebox}[Texte de version : L'art camerounais]
\zh{喀麦隆的音乐和舞很有意思。我觉得喀麦隆艺术很美，因为它有很长的历史。所以，我更喜欢喀麦隆艺术。你呢？}\\
\emph{Kāmàilóng de yīnyuè hé wǔ hěn yǒu yìsi. Wǒ juéde Kāmàilóng yìshù hěn měi, yīnwèi tā yǒu hěn cháng de lìshǐ. Suǒyǐ, wǒ gèng xǐhuan Kāmàilóng yìshù. Nǐ ne?}\\
« La musique et la danse du Cameroun sont très intéressantes. Je trouve que l'art camerounais est très beau, parce qu'il a une très longue histoire. C'est pourquoi je préfère l'art camerounais. Et toi ? »
\end{textebox}

Analysons la traduction phrase par phrase :

\begin{itemize}
\item \zh{喀麦隆的音乐和舞很有意思。} Bloc sujet : \zh{喀麦隆的音乐和舞} (la musique et la danse du Cameroun) ; bloc attribut : \zh{很有意思} (très intéressantes). En français, on accorde : « sont très intéressantes ». Notez \zh{和} (et), qui relie deux noms, jamais deux propositions.
\item \zh{我觉得喀麦隆艺术很美，因为它有很长的历史。} \zh{我觉得} = « je trouve que » ; \zh{它} (il) reprend « l'art » ; \zh{有很长的历史} = « possède une très longue histoire ».
\item \zh{所以，我更喜欢喀麦隆艺术。} \zh{所以} = « donc, c'est pourquoi » ; \zh{更喜欢} = « préférer ».
\item \zh{你呢？} Question elliptique typique : « Et toi ? » — le mot-outil \zh{呢} renvoie la question à l'interlocuteur.
\end{itemize}

\begin{center}
\begin{tabularx}{\linewidth}{|X|X|X|X|}
\hline
\rowcolor{popblueL} Caractères & Pinyin & Nature & Traduction \\
\hline
\zh{音乐} & yīnyuè & nom & musique \\
\hline
\zh{舞} & wǔ & nom & danse \\
\hline
\zh{有意思} & yǒu yìsi & loc. adj. & intéressant \\
\hline
\zh{觉得} & juéde & verbe & trouver, penser \\
\hline
\zh{美} & měi & adjectif & beau \\
\hline
\zh{历史} & lìshǐ & nom & histoire \\
\hline
\zh{所以} & suǒyǐ & connecteur & donc, c'est pourquoi \\
\hline
\zh{更} & gèng & adverbe & davantage, plus \\
\hline
\zh{呢} & ne & particule & (renvoie la question) \\
\hline
\end{tabularx}
\end{center}

\courssub{Corrections commentées : les trois erreurs types}

\begin{center}
\begin{tabularx}{\linewidth}{|X|X|X|}
\hline
\rowcolor{popblueL} Erreur du francophone & Correction & Règle \\
\hline
\zh{喜欢我中国艺术。} & \zh{我喜欢中国艺术。} & Ordre Sujet-Verbe-Complément, comme en français. \\
\hline
\zh{颜色好。} & \zh{颜色很好看。} & Adjectif attribut : ajouter \zh{很} (sans valeur intensive). \\
\hline
\zh{我喜欢中国艺术也。} & \zh{我也喜欢中国艺术。} & \zh{也} et \zh{都} avant le verbe, jamais en fin de phrase. \\
\hline
\end{tabularx}
\end{center}

\begin{exoresolu}[Thème guidé : deux phrases]
Énoncé : traduisez en chinois : a) « Moi aussi, je trouve la danse camerounaise intéressante. » b) « Parce que les masques sont beaux, j'aime l'art camerounais. »
\tcblower
Solution détaillée : a) Sujet : 我 ; « aussi » = \zh{也}, placé avant le verbe ; « je trouve » = \zh{觉得} ; « la danse camerounaise » = \zh{喀麦隆的舞} ; adjectif attribut : \zh{很有意思}. Phrase : \zh{我也觉得喀麦隆的舞很有意思。} (Wǒ yě juéde Kāmàilóng de wǔ hěn yǒu yìsi.) b) Cause d'abord : \zh{因为面具很好看} ; conséquence avec \zh{所以} : \zh{所以我喜欢喀麦隆艺术}. Phrase complète : \zh{因为面具很好看，所以我喜欢喀麦隆艺术。} (Yīnwèi miànjù hěn hǎokàn, suǒyǐ wǒ xǐhuan Kāmàilóng yìshù.) Vérification : virgule chinoise \zh{，} entre les propositions, point \zh{。} final, \zh{很} devant chaque adjectif attribut.
\end{exoresolu}

\begin{exoresolu}[Version guidée : une phrase à trous de sens]
Énoncé : traduisez en français : \zh{中国的书法有几千年的历史，我觉得很有意思。}
\tcblower
Solution détaillée : Étape 1, découpage en blocs : \zh{中国的书法} (la calligraphie de Chine) / \zh{有} (possède) / \zh{几千年的历史} (plusieurs millénaires d'histoire) ; puis \zh{我觉得} (je trouve) / \zh{很有意思} (très intéressante). Étape 2, reconstitution en français correct avec articles et accords : « La calligraphie chinoise a plusieurs millénaires d'histoire, je la trouve très intéressante. » Remarquez que \zh{它} n'apparaît pas dans la seconde proposition : le chinois omet souvent le pronom quand le sens est clair ; le français, lui, exige « je \emph{la} trouve ».
\end{exoresolu}

\begin{exercice}[À vous de traduire]
Thème : 1) « J'aime la musique chinoise. » 2) « La calligraphie a une longue histoire. » 3) « Nous préférons tous l'art camerounais. » Version : traduisez en français : \zh{中国画很有意思，我也很喜欢书法。你呢？}
\end{exercice}

\begin{corrige}[Thème et version]
Thème : 1) \zh{我喜欢中国音乐。} (Wǒ xǐhuan Zhōngguó yīnyuè.) 2) \zh{书法有很长的历史。} (Shūfǎ yǒu hěn cháng de lìshǐ.) — \zh{有} = posséder, donc correct. 3) \zh{我们都更喜欢喀麦隆艺术。} (Wǒmen dōu gèng xǐhuan Kāmàilóng yìshù.) — \zh{都} avant le verbe. Version : « La peinture chinoise est très intéressante, et j'aime aussi beaucoup la calligraphie. Et toi ? »
\end{corrige}

\begin{aretenir}[L'essentiel de la section]
\begin{itemize}
\item Thème : repérer la structure française, convertir en SVO chinois, choisir les caractères appris, relire traits et radicaux.
\item Version : identifier sujet / verbe / complément, traduire par blocs de sens, reconstituer un français correct.
\item Adjectif attribut : toujours \zh{很} + adjectif (\zh{很好看}, \zh{很有意思}) dans une phrase affirmative.
\item \zh{也} et \zh{都} se placent avant le verbe ; \zh{因为}... \zh{所以}... encadrent cause et conséquence.
\item \zh{有} traduit « avoir / il y a » seulement pour la possession ou l'existence, jamais pour « il y a + durée écoulée ».
\item \zh{的} relie le déterminant au nom : \zh{喀麦隆的艺术} (l'art du Cameroun).
\item Ponctuation : \zh{，} entre les propositions, \zh{。} en fin de phrase, \zh{？} pour les questions, sans espace avant.
\end{itemize}
\end{aretenir}

\coursec{Production écrite et évaluation}

Après avoir travaillé la traduction thème/version dans la section précédente, nous arrivons à l'aboutissement du module : la production écrite autonome. Vous disposez désormais du vocabulaire, des structures grammaticales (\zh{比较句}, \zh{连词}) et des modèles de caractères nécessaires. Cette section vous guide pas à pas vers la rédaction d'un texte cohérent, noté sur 20, conforme aux attentes du Baccalauréat.

\courssub{Tâche finale : Rédiger un texte d'opinion comparatif}

\cle{Consigne officielle} : \textbf{« 你喜欢喀麦隆还是中国艺术？为什么？ »} (Tu aimes l'art camerounais ou l'art chinois ? Pourquoi ?).

Il s'agit d'un \cle{texte argumentatif simple} : vous devez \textbf{choisir un camp} (ou nuancer), \textbf{donner au moins deux arguments} structurés par des connecteurs logiques, et \textbf{conclure}. La longueur visée est de \textbf{6 à 10 phrases} (environ 80--120 caractères chinois).

\begin{textebox}{Modèle de production attendue (niveau HSK 3--4)}
\zh{我更喜欢喀麦隆艺术。} \emph{Wǒ gèng xǐhuan Kāmàilóng yìshù.} « Je préfère l'art camerounais. »
\zh{首先，喀麦隆的面具雕刻非常精美，颜色鲜艳。} \emph{Shǒuxiān, Kāmàilóng de miànjù diāokè fēicháng jīngměi, yánsè xiānyàn.} « D'abord, la sculpture sur masques camerounaise est très raffinée, les couleurs sont vives. »
\zh{其次，我们的传统音乐和舞蹈充满活力，比如马康萨。} \emph{Qícì, wǒmen de chuántǒng yīnyuè hé wǔdǎo chōngmǎn huólì, bǐrú Mǎkāngsà.} « Ensuite, notre musique et notre danse traditionnelles sont pleines de vitalité, comme le Makossa. »
\zh{再者，喀麦隆艺术反映了我们多样的民族文化。} \emph{Zài zhě, Kāmàilóng yìshù fǎnyìng le wǒmen duōyàng de mínzú wénhuà.} « En outre, l'art camerounais reflète notre culture ethnique diverse. »
\zh{虽然中国书法和京剧也很有名，历史悠久，} \emph{Suīrán Zhōngguó shūfǎ hé Jīngjù yě hěn yǒumíng, lìshǐ yōujiǔ,} « Bien que la calligraphie chinoise et l'opéra de Pékin soient aussi célèbres et aient une longue histoire, »
\zh{但我更爱自己的文化。} \emph{dàn wǒ gèng ài zìjǐ de wénhuà.} « moi j'aime davantage ma propre culture. »
\zh{所以，我为喀麦隆艺术感到骄傲。} \emph{Suǒyǐ, wǒ wèi Kāmàilóng yìshù gǎndào jiāo'ào.} « Donc, je suis fier de l'art camerounais. »
\zh{你呢？} \emph{Nǐ ne?} « Et toi ? »
\end{textebox}

\begin{center}
\begin{tabularx}{\linewidth}{|X|X|X|X|}
\hline
\rowcolor{popblueL} \textbf{Caractères} & \textbf{Pinyin} & \textbf{Nature} & \textbf{Traduction} \\
\hline
面具 & miànjù & nom & masque \\
雕刻 & diāokè & verbe/nom & sculpter / sculpture \\
精美 & jīngměi & adj & raffiné, magnifique \\
鲜艳 & xiānyàn & adj & vif, éclatant (couleur) \\
传统 & chuántǒng & adj/nom & tradition(nel) \\
充满 & chōngmǎn & verbe & être plein de \\
活力 & huólì & nom & vitalité, énergie \\
比如 & bǐrú & conj & par exemple \\
反映 & fǎnyìng & verbe & refléter \\
多样 & duōyàng & adj & varié, divers \\
民族 & mínzú & nom & ethnie, groupe ethnique \\
书法 & shūfǎ & nom & calligraphie \\
京剧 & Jīngjù & nom & opéra de Pékin \\
历史悠久 & lìshǐ yōujiǔ & loc. adj & avoir une longue histoire \\
唱腔 & chàngqiāng & nom & chant, mélodie (opéra) \\
骄傲 & jiāo'ào & adj/verbe & fier / fierté \\
\hline
\end{tabularx}
\end{center}

\courssub{Étapes de la production écrite : de l'idée au texte final}

Pour ne pas faire face à la page blanche, suivez cette \cle{méthode en 4 étapes} validée par l'expérience des correcteurs du Bac.

\begin{methode}{Réussir l'expression écrite au Bac : la méthode en 4 étapes}
\begin{enumerate}
\item \textbf{Brainstorming (3 min) :} Sur votre brouillon, faites deux colonnes : \zh{喀麦隆艺术} et \zh{中国艺术}. Notez 3--4 mots-clés par colonne (masques, makossa, tamtam, pagne, danse traditionnelle / calligraphie, opéra de Pékin, découpe de papier, musique traditionnelle). \textbf{Choisissez votre camp} (ou « les deux » avec nuance).
\item \textbf{Plan en 4 parties (2 min) :} Structurez votre texte :
      \begin{itemize}
      \item \textbf{Phrase 1 : Opinion claire} (\zh{我更喜欢... / 我喜欢...和...}).
      \item \textbf{Phrases 2--4 : Arguments + connecteurs} (\zh{首先...；其次...；再者...}).
      \item \textbf{Phrase 5 (optionnelle) : Concession} (\zh{虽然...但是...}) pour montrer la maturité.
      \item \textbf{Dernière phrase : Conclusion avec \zh{所以}} (\zh{所以我喜欢...}).
      \end{itemize}
\item \textbf{Rédaction soignée (10 min) :} Écrivez au propre en caractères. \textbf{Respectez l'ordre des mots} (Sujet + Temps + Lieu + Verbe + Complément). Utilisez \zh{很} devant les adjectifs. Vérifiez les classificateurs (\zh{个, 种, 首, 幅}).
\item \textbf{Relecture active (5 min) :} Relisez \textbf{à voix basse} (bouche en mouvement). Cochez : \cle{Caractères} (traits, radicaux) ? \cle{Ordre des mots} ? \cle{Connecteurs} ? \cle{很/不} devant adj ? \cle{Ponctuation chinoise} (，。；！) ?
\end{enumerate}
\end{methode}

\begin{popfigure}
\legende{Organigramme de la structure du texte argumentatif (4 parties)}
\begin{tikzpicture}[
    node distance=1.0cm and 1.2cm,
    box/.style={draw=popblue, thick, rounded corners, fill=popblueL, text width=3.0cm, align=center, minimum height=1.8cm},
    arrow/.style={-Stealth, thick, popdark},
    title/.style={font=\bfseries\small, popdark}
]
\node[box, title, align=center] (opinion){1. Opinion claire\\ \zh{我更喜欢喀麦隆艺术。}};
\node[box, title, right=of opinion, align=center] (arg){2. Arguments (2--3)\\ \zh{首先... 其次... 再者...}};
\node[box, title, right=of arg, align=center] (concession){3. Concession (optionnel)\\ \zh{虽然中国书法很美，但是...}};
\node[box, title, right=of concession, align=center] (conclu){4. Conclusion\\ \zh{所以我为...骄傲。}};

\draw[arrow] (opinion) -- node[above, font=\tiny] {donc} (arg);
\draw[arrow] (arg) -- node[above, font=\tiny] {mais} (concession);
\draw[arrow] (concession) -- node[above, font=\tiny] {donc} (conclu);
\end{tikzpicture}
\end{popfigure}

\courssub{Connecteurs logiques indispensables (Tableau récapitulatif)}

Maîtriser ces connecteurs, c'est gagner 4 points faciles sur le critère « Structure ».

\begin{center}
\begin{tabularx}{\linewidth}{|X|X|X|X|}
\hline
\rowcolor{popblueL} \textbf{Fonction} & \textbf{Connecteur chinois} & \textbf{Pinyin} & \textbf{Français} \\
\hline
Énumérer (1er arg) & 首先 & shǒuxiān & D'abord, en premier lieu \\
Énumérer (2e arg) & 其次 / 然后 & qícì / ránhòu & Ensuite, puis \\
Énumérer (3e arg) & 再者 / 还有 & zài zhě / hái yǒu & De plus, en outre \\
Contraste / Concession & 虽然...，但是/不过... & suīrán..., dànshì/búguò... & Bien que..., mais... \\
Cause & 因为...，所以... & yīnwèi..., suǒyǐ... & Parce que..., donc... \\
Conclusion & 所以 / 因此 & suǒyǐ / yīncǐ & Donc, par conséquent \\
\hline
\end{tabularx}
\end{center}

\courssub{Banque de phrases de secours (Trames à compléter)}

Pour les élèves en difficulté ou en panne d'inspiration, ces \cle{trames à trous} garantissent un texte structuré et grammatical. Remplacez les parties soulignées.

\begin{exemplebox}{Trame A : Choix unique (niveau standard)}
\zh{我更喜欢 \underline{喀麦隆艺术}。} (Je préfère \underline{l'art camerounais}.)
\zh{首先，\underline{喀麦隆的面具很精美}。} (D'abord, \underline{les masques camerounais sont raffinés}.)
\zh{其次，\underline{我们的音乐很有活力}。} (Ensuite, \underline{notre musique est pleine de vie}.)
\zh{所以，\underline{我为喀麦隆艺术感到骄傲}。} (Donc, \underline{je suis fier de l'art camerounais}.)
\end{exemplebox}

\begin{exemplebox}{Trame B : Nuance / Les deux (niveau avancé)}
\zh{我喜欢 \underline{喀麦隆艺术}，也喜欢 \underline{中国艺术}。} (J'aime \underline{l'art camerounais}, et aussi \underline{l'art chinois}.)
\zh{首先，\underline{喀麦隆的舞蹈很有趣}。} (D'abord, \underline{la danse camerounaise est amusante}.)
\zh{然后，\underline{中国书法很有历史}。} (Ensuite, \underline{la calligraphie chinoise a de l'histoire}.)
\zh{虽然 \underline{它们很不一样}，但是 \underline{我都喜欢}。} (Bien que \underline{elles soient très différentes}, \underline{je les aime toutes les deux}.)
\zh{所以，\underline{学习两种艺术很好}。} (Donc, \underline{apprendre les deux arts est bien}.)
\end{exemplebox}

\courssub{Grille d'évaluation sur 20 points (Barème officiel type Bac)}

Cette grille est utilisée pour l'auto-évaluation et l'évaluation croisée. Chaque critère est noté sur \textbf{4 points} (4 = Excellent, 3 = Bien, 2 = Passable, 1 = Insuffisant, 0 = Absent).

\begin{popfigure}
\legende{Grille d'évaluation détaillée (5 critères × 4 points = 20 pts)}
\begin{tikzpicture}[
    cell/.style={draw=popline, minimum height=1cm, align=center, font=\small},
    header/.style={cell, fill=popblue, text=white, font=\bfseries\small},
    crit/.style={cell, fill=popblueL, font=\bfseries\small, text width=3.5cm},
    score/.style={cell, text width=2.2cm},
    desc/.style={cell, fill=popcream, text width=3.5cm, font=\tiny}
]
% En-têtes
\node[header, text width=3.5cm] (c0) at (0,0) {Critère};
\node[header, text width=2.2cm] (c4) at (3.5,0) {4 pts \\ Excellent};
\node[header, text width=2.2cm] (c3) at (5.7,0) {3 pts \\ Bien};
\node[header, text width=2.2cm] (c2) at (7.9,0) {2 pts \\ Passable};
\node[header, text width=2.2cm] (c1) at (10.1,0) {1 pt \\ Insuffisant};
\node[header, text width=2.2cm] (c0) at (12.3,0) {0 pt \\ Absent};

% Ligne 1 : Sujet
\node[crit, align=center] (r1c0) at (0,-1){1. Respect du sujet\\ (Réponse à « pourquoi ? »)};
\node[score] (r1c4) at (3.5,-1) {Opinion claire, arguments pertinents, conclusion};
\node[score] (r1c3) at (5.7,-1) {Opinion claire, arguments présents};
\node[score] (r1c2) at (7.9,-1) {Opinion floue ou arguments hors sujet};
\node[score] (r1c1) at (10.1,-1) {Hors sujet total};
\node[score] (r1c0) at (12.3,-1) {Page blanche};

% Ligne 2 : Structure
\node[crit, align=center] (r2c0) at (0,-2){2. Structure \& Connecteurs\\ (首 先, 所 以, 虽 然...)};
\node[score] (r2c4) at (3.5,-2) {4 parties, connecteurs variés et justes};
\node[score] (r2c3) at (5.7,-2) {3 parties, connecteurs corrects};
\node[score] (r2c2) at (7.9,-2) {Structure lâche, 1--2 connecteurs};
\node[score] (r2c1) at (10.1,-2) {Phrases disjointes, sans lien};
\node[score] (r2c0) at (12.3,-2) {Aucun connecteur};

% Ligne 3 : Caractères
\node[crit, align=center] (r3c0) at (0,-3){3. Exactitude des caractères\\ (Traits, radicaux, 术 vs 木)};
\node[score] (r3c4) at (3.5,-3) {Zéro erreur, écriture soignée};
\node[score] (r3c3) at (5.7,-3) {1--2 erreurs mineures (traits)};
\node[score] (r3c2) at (7.9,-3) {3--4 erreurs ou radicaux faux};
\node[score] (r3c1) at (10.1,-3) {Beaucoup d'erreurs, illisible};
\node[score] (r3c0) at (12.3,-3) {Pas de caractères (pinyin seul)};

% Ligne 4 : Grammaire
\node[crit, align=center] (r4c0) at (0,-4){4. Grammaire \& Ordre des mots\\ (SVO, 很+adj, 比较句)};
\node[score] (r4c4) at (3.5,-4) {Phrases correctes, structures variées};
\node[score] (r4c3) at (5.7,-4) {Quelques erreurs mineures (了/过)};
\node[score] (r4c2) at (7.9,-4) {Erreurs fréquentes (ordre, 很 absent)};
\node[score] (r4c1) at (10.1,-4) {Phrases incompréhensibles};
\node[score] (r4c0) at (12.3,-4) {Aucune phrase correcte};

% Ligne 5 : Vocabulaire
\node[crit, align=center] (r5c0) at (0,-5){5. Richesse du vocabulaire\\ (Arts, adj, classif.)};
\node[score] (r5c4) at (3.5,-5) {Vocabulaire précis, varié (HSK 3+), classif. justes};
\node[score] (r5c3) at (5.7,-5) {Vocabulaire correct, un peu répétitif};
\node[score] (r5c2) at (7.9,-5) {Vocabulaire basique (HSK 1), répétitif};
\node[score] (r5c1) at (10.1,-5) {Mots inventés, anglicismes, français};
\node[score] (r5c0) at (12.3,-5) {Aucun mot chinois};
\end{tikzpicture}
\end{popfigure}

\courssub{Auto-évaluation et évaluation croisée entre pairs}

L'évaluation n'est pas une sanction, c'est un \cle{outil d'apprentissage}. Voici le protocole à suivre en classe (ou chez vous) :

\begin{enumerate}
\item \textbf{Auto-évaluation (5 min) :} Prenez votre grille. Lisez votre texte \textbf{à voix haute}. Cochez la case correspondante pour chaque critère. Soyez honnête : « Ai-je vraiment mis \zh{很} devant \zh{精美} ? Ai-je écrit \zh{术} correctement ? ». Notez votre total /20.
\item \textbf{Échange de copies (binôme) :} Donnez votre texte + grille vierge à votre voisin. Lui fait de même.
\item \textbf{Évaluation par les pairs (5 min) :} Votre pair lit votre texte, coche la grille, et \textbf{écrit un commentaire constructif} en bas (ex: « Argument 2 super, mais attention à l'ordre des mots dans la phrase 3 : \zh{我喜欢很音乐} $\to$ \zh{我很喜欢音乐} »).
\item \textbf{Retour et révision (5 min) :} Récupérez votre copie. Lisez le commentaire. \textbf{Corrigez vos erreurs au stylo vert} sur votre brouillon ou en marge. C'est là que se joue la progression.
\item \textbf{Remise finale :} Le professeur collecte les versions corrigées pour la note officielle.
\end{enumerate}

\begin{exercice}{Entraînement à l'évaluation par les pairs}
Voici une production d'élève fictif. Utilisez la grille pour la noter /20 et écrivez un commentaire constructif (2 phrases max).

\zh{我喜欢中国艺术。因为中国书法很美。中国功夫也很有意思。喀麦隆艺术也不错。所以我喜欢中国艺术。}
\end{exercice}

\begin{corrige}{Exercice -- Évaluation par les pairs}
\textbf{Note suggérée : 13/20}
\begin{itemize}
\item \textbf{Sujet (3/4) :} Opinion claire, mais arguments un peu généraux (« très beau », « intéressant »).
\item \textbf{Structure (3/4) :} 5 phrases, connecteurs \zh{因为/所以} utilisés, mais pas de \zh{首先/其次}, pas de concession.
\item \textbf{Caractères (3/4) :} Supposés corrects (texte tapé).
\item \textbf{Grammaire (2/4) :} Phrase 1 : \zh{我喜欢中国艺术} (OK). Phrase 2 : \zh{因为中国书法很美} (OK). Phrase 4 : \zh{喀麦隆艺术也不错} (contraste mal géré, devrait être \zh{虽然喀麦隆艺术也不错，但是...}).
\item \textbf{Vocabulaire (2/4) :} Vocabulaire HSK 2 basique (\zh{美, 有意思, 不错}). Pas de vocabulaire spécifique arts (masques, opéra, calligraphie, sculpture).
\end{itemize}
\textbf{Commentaire type :} « Ton opinion est claire et la structure 因为/所以 fonctionne. Pour monter à 16+, ajoute \zh{首先/其次}, précise \zh{书法} ou \zh{京剧} au lieu de dire juste "art", et essaie une phrase avec \zh{虽然...但是} pour nuancer. »
\end{corrige}

\begin{attention}[Erreurs fréquentes des francophones (Pièges à éviter avant de rendre)]

\begin{enumerate}
\item \textbf{Oublier \zh{很} devant l'adjectif attributif.}
      \zh{喀麦隆艺术精美。} $\times$ (Phrase incomplète en chinois standard)
      \zh{喀麦隆艺术\cle{很}精美。} $\checkmark$ « L'art camerounais est magnifique. »
\item \textbf{Inverser \zh{喜欢} et son complément (structure française « J'aime la Chine » $\to$ *\zh{我中国喜欢} ).}
      \textbf{Règle immuable :} \zh{Sujet + 喜欢 + Objet}. \zh{我喜欢喀麦隆。}
\item \textbf{Écrire \zh{术} (art, technique) sans le point final (le trait « 丶 » en bas à droite).}
      \zh{术} (correct) vs \zh{朮} (incorrect, forme ancienne/erreur de trait). Le radical est \zh{殳} (shū, lance/arme), le trait final est un point.
\item \textbf{Mélanger \zh{还是} (choix dans la QUESTION) et \zh{或者} (choix dans l'AFFIRMATION).}
      \zh{你喜欢喀麦隆\cle{还是}中国？} (Question)
      \zh{我喜欢喀麦隆\cle{或者}中国。} (Affirmation : « j'aime le Cameroun ou la Chine » / indifférent).
      \textbf{Piège Bac :} La consigne utilise \zh{还是}, votre réponse n'a \textbf{pas} à la réutiliser sauf si vous citez la question.
\item \textbf{Ordre des mots : Complément de lieu/temps APRÈS le verbe.}
      \zh{*我在喀麦隆喜欢艺术。} $\times$
      \zh{我\cle{在喀麦隆}\cle{喜欢}艺术。} $\checkmark$ (Lieu avant verbe).
\item \textbf{Classificateurs oubliés ou wrong.}
      \zh{一个艺术} $\times$ (art non comptable ainsi).
      \zh{一种艺术} $\checkmark$ (un type d'art) / \zh{一幅画} (un tableau) / \zh{一个面具} (un masque).
\end{enumerate}

\end{attention}

\begin{savaistu}[Compétence citoyenne : comparer sans dénigrer]
Savoir écrire « \zh{我更喜欢喀麦隆艺术，但是中国书法也很棒} » (Je préfère l'art camerounais, mais la calligraphie chinoise est super aussi), c'est exercer la \cle{citoyenneté mondiale}. Le module « Citoyenneté et ouverture au monde » vise à former des ambassadeurs culturels. Un bon texte argumentatif ne dit pas « l'autre culture est nulle », il dit « ma culture me touche plus \textbf{pour telle raison précise}, et je respecte l'autre ». C'est la différence entre un stéréotype et une opinion argumentée. Au Bac, les correcteurs valorisent l'ouverture d'esprit (critère « Richesse vocabulaire / Maturité »).

\end{savaistu}

\begin{exoresolu}[Analyse d'une production d'élève niveau « Bien » (16/20)]

\textbf{Texte de l'élève (Léa, 1ère A) :}
\zh{我更喜欢喀麦隆艺术。首先，我们的面具雕刻很精美，颜色很鲜艳。其次，马康萨音乐充满活力，让我们想跳舞。虽然中国京剧很有名，唱腔很美，但我听不太懂。所以，我为喀麦隆文化感到骄傲。你喜欢哪个？}

\textbf{Analyse détaillée avec la grille :}

\begin{center}
\begin{tabularx}{\linewidth}{|X|X|X|}
\hline
\rowcolor{popblueL} \textbf{Critère} & \textbf{Note} & \textbf{Justification} \\
\hline
1. Respect du sujet & 4/4 & Opinion nette (\zh{更喜欢}), 2 arguments concrets (masques, makossa), concession (京剧), conclusion (\zh{所以}), question finale engageante. \\
\hline
2. Structure \& Connecteurs & 4/4 & Plan 4 parties respecté. Connecteurs variés : \zh{首先、其次、虽然...但/但是、所以}. Ponctuation chinoise parfaite. \\
\hline
3. Caractères & 3/4 & \zh{雕刻, 鲜艳, 唱腔} bien écrits. \textbf{Erreur mineure :} \zh{腔} (qiāng) écrit sans le radical \zh{月} (viande/corps) à gauche (écrit *\zh{空} au lieu de \zh{腔}). \\
\hline
4. Grammaire & 3/4 & Ordre parfait. \zh{让我们想跳舞} (structure causative \zh{让} bien utilisée). \textbf{Point faible :} \zh{听不太懂} (complément de degré potentiel) correct mais \zh{不太} placé après \zh{听} est familier ; standard : \zh{听不大懂} ou \zh{不太听得懂} est plus formel. \\
\hline
5. Vocabulaire & 2/4 & Bon vocabulaire spécifique (\zh{面具雕刻, 鲜艳, 马康萨, 活力, 京剧, 唱腔, 骄傲}). \textbf{Manque :} Pas de classificateur numérique (\zh{一种}, \zh{首}) pour montrer la maîtrise HSK 3. Répétition de \zh{很}. \\
\hline
\textbf{TOTAL} & \textbf{16/20} & \textbf{Excellent travail. Progression : varier les adverbes de degré (\zh{非常、特别、挺}), ajouter un classificateur, vérifier le radical de \zh{腔}.} \\
\hline
\end{tabularx}
\end{center}

\end{exoresolu}

\begin{aretenir}[Check-list finale avant l'examen (Bac Blanc / Bac)]
\begin{itemize}
\item \textbf{Temps imparti :} 20 minutes pour 8 phrases. \textbf{Chronométrez-vous} à l'entraînement.
\item \textbf{Brouillon OBLIGATOIRE :} Plan + vocabulaire clé + phrases pièges (\zh{虽然...但是}).
\item \textbf{Caractères :} Écrivez grand, lisible. Si vous oubliez un caractère : \textbf{le remplacez par un synonyme connu} (ex: \zh{漂亮} au lieu de \zh{精美}), ne laissez pas de blanc ni de pinyin.
\item \textbf{Relecture :} 3 minutes minimum. Chassez les \zh{很} manquants, les \zh{还是}/\zh{或者}, l'ordre [Lieu/Temps] + Verbe.
\item \textbf{Le « plus » qui fait la différence :} Une phrase de concession (\zh{虽然...但是...}) et une question finale (\zh{你呢？}) montrent l'interaction orale/écrite.
\end{itemize}
\end{aretenir}


\newpage

\coursec{Activité d'intégration}

\coursec{Leçon 22 — Expression écrite : culture chinoise et camerounaise (你喜欢喀麦隆还是中国艺术？为什么？)}

\courssub{Objectifs de l'activité}
\begin{aretenir}[Objectifs visés]
À l'issue de cette activité, l'élève sera capable de :
\begin{itemize}
    \item Mobiliser le lexique des arts (peinture, sculpture, musique, danse, architecture) et des opinions pour comparer deux cultures.
    \item Structurer un texte argumentatif court (6--10 phrases) en utilisant des connecteurs logiques (\zh{首先}、\zh{然后}、\zh{还有}、\zh{所以}).
    \item Écrire correctement les caractères clés de la leçon (ordre des traits, radicaux) dans un contexte de production autonome.
    \item S'auto-évaluer et évaluer un pair à l'aide d'une grille de critères (contenu, langue, écriture, cohérence).
    \item Présenter oralement son point de vue avec une prononciation soignée (tons, liaisons).
\end{itemize}
\end{aretenir}

\courssub{Situation}
\begin{textebox}[Contexte : Le magazine du Club de Chinois du Lycée]
Le Club de Chinois de votre lycée (Lycée de la Réunification, Yaoundé) prépare un numéro spécial de son magazine bilingue \emph{Écho du Dragon / 龙声} intitulé \zh{« 喀麦隆艺术和中国艺术 »} (L'art camerounais et l'art chinois). La rédactrice en chef, votre professeure, lance un appel à contributions :

\shortstack[l]{\zh{亲爱的同学们：}\\ \zh{大家好！}\\ \zh{我们要出版《龙声》特刊，主题是“喀麦隆艺术和中国艺术”。}\\ \zh{请你写一篇短文（6--10句话），回答：}\\ \zh{“你喜欢喀麦隆还是中国艺术？为什么？”}\\ \zh{要求：}\\ \zh{1. 结构清晰（四个部分）；}\\ \zh{2. 用上至少三个连词（首先、然后、还有、所以）；}\\ \zh{3. 字迹工整，生字无误。}\\ \zh{优秀作品将在“中喀文化墙”展出，并有机会在校园文化节上朗读。}\\ \zh{期待你的佳作！}\\ \zh{中文社编辑部}}
\\
\textbf{Traduction :} « Chers camarades, bonjour ! Nous allons publier un numéro spécial de \emph{L'Écho du Dragon} sur le thème "L'art camerounais et l'art chinois". Merci de rédiger un court texte (6 à 10 phrases) répondant à la question : "Préférez-vous l'art camerounais ou l'art chinois ? Pourquoi ?" Exigences : structure claire (4 parties), utilisation d'au moins 3 connecteurs (d'abord, ensuite, de plus, donc), écriture soignée, pas d'erreurs de caractères. Les meilleurs travaux seront affichés sur le "Mur culturel Sino-Camerounais" et pourront être lus lors de la fête culturelle de l'école. Nous attendons vos chefs-d'œuvre ! Comité de rédaction du Club de Chinois. »
\end{textebox}

\courssub{Consignes}
\begin{methode}[Déroulé de l'activité en 5 étapes]
\begin{enumerate}
    \item \textbf{Analyse du sujet et planification (10 min).} En binôme, lisez l'appel à contributions. Identifiez la question centrale (\zh{你喜欢喀麦隆还是中国艺术？为什么？}) et les contraintes (longueur, connecteurs, structure en 4 parties). Complétez le canevas ci-dessous dans votre cahier de brouillon :
    \begin{center}
    \begin{tabularx}{\linewidth}{|X|X|}\hline
    \rowcolor{popblueL} \textbf{Partie du texte} & \textbf{Mes idées / Mots-clés (en chinois si possible)} \\ \hline
    1. Introduction : Mon choix global (\zh{我喜欢……艺术。}) &  \\ \hline
    2. Argument 1 : \zh{首先}… (Description d'un art précis) &  \\ \hline
    3. Argument 2 : \zh{然后 / 还有}… (Comparaison / Autre aspect) &  \\ \hline
    4. Conclusion : \zh{所以}… (Justification finale / Sentiment) &  \\ \hline
    \end{tabularx}
    \end{center}
    \begin{popfigure}
    \begin{tikzpicture}[node distance=1.2cm, >=Stealth,
        box/.style={draw=popblue, thick, rounded corners, fill=popblueL, text width=2.8cm, align=center, minimum height=1.5cm, font=\small},
        arrow/.style={->, thick, popdark}
    ]
    \node[box, align=center] (intro){1. Introduction\\ \zh{我喜欢…艺术。}\\ (Thèse claire)};
    \node[box, right=of intro, align=center] (arg1){2. Argument 1\\ \zh{首先}…\\ (Description visuelle / statuaire)};
    \node[box, right=of arg1, align=center] (arg2){3. Argument 2 / Preuve\\ \zh{然后 / 还有}…\\ (Expérience vécue / Danse)};
    \node[box, right=of arg2, align=center] (concl){4. Conclusion\\ \zh{所以}…\\ (Fierté / Identité)};
    \draw[arrow] (intro) -- (arg1);
    \draw[arrow] (arg1) -- (arg2);
    \draw[arrow] (arg2) -- (concl);
    \end{tikzpicture}
    \legende{Structure type du texte argumentatif en 4 parties}
    \end{popfigure}
    \item \textbf{Rédaction de l'article (25 min).} Rédigez votre article \textbf{uniquement en caractères chinois} (pinyin interdit dans la version finale) sur la fiche fournie (format A5). Respectez le plan en 4 parties. Utilisez \textbf{au moins 3 connecteurs} parmi la liste imposée. Vérifiez l'écriture de chaque caractère (ordre des traits, radicaux vus en classe : \zh{艺、文、舞、乐、比、较}).
    \item \textbf{Correction croisée par les pairs (15 min).} Échangez votre fiche avec votre binôme. Utilisez la \textbf{Grille d'évaluation} ci-dessous pour annoter le texte de votre camarade (cochez, soulignez, notez une observation bienveillante).
    \begin{center}
    \begin{tabularx}{\linewidth}{|X|c|c|c|}\hline
    \rowcolor{popblueL} \textbf{Critère} & \textbf{✓ Acquis} & \textbf{~ En cours} & \textbf{Commentaire / Exemple d'erreur} \\ \hline
    Contenu : Réponse claire à la question, 4 parties respectées & & & \\ \hline
    Langue : Au moins 3 connecteurs utilisés correctement & & & \\ \hline
    Vocabulaire : Lexique arts / opinions précis (min. 5 mots du cours) & & & \\ \hline
    Grammaire : Structures \zh{比}、\zh{因为……所以……}、\zh{觉得} correctes & & & \\ \hline
    Écriture : Caractères lisibles, traits dans l'ordre, pas de confusion & & & \\ \hline
    \end{tabularx}
    \end{center}
    \item \textbf{Réécriture et finalisation (10 min).} Récupérez votre fiche, lisez les commentaires de votre pair. Corrigez vos erreurs (caractères mal formés, connecteur manquant, ordre des mots) à l'encre verte. Remettez la version finale à la professeure.
    \item \textbf{Partage oral et vote (10 min).} Un volontaire lit son article à voix haute (travail sur la prosodie : groupes rythmiques, tons). La classe vote à main levée pour l'argument le plus convaincant (\zh{最有说服力的论点}). Les 3 meilleurs articles sont affichés au tableau d'expression culturelle.
\end{enumerate}
\end{methode}

\courssub{Ressources à mobiliser}
\begin{definition}[Lexique essentiel : Arts et Opinions (Rappel Leçon 22)]
\begin{center}
\begin{tabularx}{\linewidth}{|X|X|X|X|}\hline
\rowcolor{popblueL} \textbf{Caractères} & \textbf{Pinyin} & \textbf{Nature} & \textbf{Traduction} \\ \hline
艺术 & yìshù & n. & art \\ \hline
绘画 & huìhuà & n. & peinture \\ \hline
雕塑 & diāosù & n. & sculpture \\ \hline
音乐 & yīnyuè & n. & musique \\ \hline
舞蹈 & wǔdǎo & n. & danse \\ \hline
建筑 & jiànzhù & n. & architecture \\ \hline
书法 & shūfǎ & n. & calligraphie \\ \hline
剪纸 & jiǎnzhǐ & n. & papier découpé \\ \hline
面具 & miànjù & n. & masque \\ \hline
木雕 & mùdiāo & n. & sculpture sur bois \\ \hline
传统 & chuántǒng & adj./n. & traditionnel / tradition \\ \hline
现代 & xiàndài & adj./n. & moderne / modernité \\ \hline
独特 & dútè & adj. & unique, distinctif \\ \hline
丰富 & fēngfù & adj. & riche, varié \\ \hline
有名 & yǒumíng & adj. & célèbre \\ \hline
有力 & yǒulì & adj. & puissant, vigoureux \\ \hline
节奏感 & jiézòugǎn & n. & sens du rythme \\ \hline
文化 & wénhuà & n. & culture \\ \hline
骄傲 & jiāo'ào & adj. & fier \\ \hline
喜欢 & xǐhuan & v. & aimer, apprécier \\ \hline
觉得 & juéde & v. & trouver, penser (sentiment) \\ \hline
认为 & rènwéi & v. & considérer, estimer (opinion) \\ \hline
比较 & bǐjiào & adv. & relativement, assez \\ \hline
因为 & yīnwèi & conj. & parce que \\ \hline
所以 & suǒyǐ & conj. & donc, par conséquent \\ \hline
首先 & shǒuxiān & adv. & d'abord, en premier lieu \\ \hline
然后 & ránhòu & adv. & ensuite, puis \\ \hline
还有 & hái yǒu & adv. & de plus, en outre \\ \hline
虽然 & suīrán & conj. & bien que \\ \hline
但是 & dànshì & conj. & mais \\ \hline
让 & ràng & v. & faire, laisser, rendre (causatif) \\ \hline
\end{tabularx}
\end{center}
\end{definition}

\begin{propriete}[Structures grammaticales clés pour l'argumentation]
\begin{itemize}
    \item \textbf{Comparaison avec \zh{比} (bǐ) :} \zh{主语 A + 比 + 主语 B + 形容词} (A est plus [adj] que B). \emph{Ex:} \zh{中国书法比喀麦隆绘画难。} (La calligraphie chinoise est plus difficile que la peinture camerounaise.) \textbf{Attention :} Pas de \zh{很} après l'adjectif. Négation : \zh{没(有)比} (A n'est pas plus [adj] que B). \emph{Ex:} \zh{喀麦隆舞蹈没比中国舞蹈难。}
    \item \textbf{Justification \zh{因为……所以……} (yīnwèi... suǒyǐ...) :} Lie la cause à la conséquence. La proposition \zh{因为} précède souvent le sujet principal. \emph{Ex:} \zh{因为非洲鼓很有节奏感，所以我喜欢喀麦隆音乐。}
    \item \textbf{Expression de l'opinion \zh{觉得} (juéde) / \zh{认为} (rènwéi) :} \zh{主语 + 觉得 / 认为 + 从句}. \zh{觉得} exprime un ressenti ; \zh{认为} un jugement plus réfléchi. \emph{Ex:} \zh{我觉得中国剪纸很独特。} / \zh{我认为喀麦隆雕塑更有力量。}
    \item \textbf{Énumération / Structuration :} \zh{首先……，然后……，还有……，所以……}. Ces connecteurs se placent généralement en début de proposition (après le sujet éventuel : \zh{我首先觉得…} ou \zh{首先，…}).
    \item \textbf{Aspect expérientiel \zh{过} (guo) :} \zh{动词 + 过} indique une expérience vécue au moins une fois dans le passé. \emph{Ex:} \zh{我去过雅温得博物馆。} (Je suis allé au musée de Yaoundé.)
    \item \textbf{Structure causative \zh{让} (ràng) :} \zh{主语 + 让 + 对象 + 形容词/动词} (Le sujet fait que l'objet devient...). \emph{Ex:} \zh{喀麦隆艺术让我骄傲。}
\end{itemize}
\end{propriete}

\begin{attention}[Pièges fréquents pour les francophones (HSK 1-3)]
\begin{itemize}
    \item \textbf{Ordre des mots (Cause / Conséquence) :} En français : "J'aime l'art chinois \emph{parce qu'il est} ancien." En chinois : \zh{我喜欢中国艺术，因为它很古老。} ou \zh{因为中国艺术很古老，所以我喜欢它。} Ne dites pas : \zh{*我因为它很古老喜欢中国艺术。} (La proposition \zh{因为} ne s'insère pas entre le sujet et le verbe principal).
    \item \textbf{Connecteurs vs Conjonctions :} \zh{所以} (donc) répond à \zh{因为}. \zh{但是} (mais) répond à \zh{虽然}. Ne les mélangez pas : \zh{*虽然我很喜欢，所以我不买。} $\rightarrow$ \zh{虽然我很喜欢，但是我不买。}
    \item \textbf{Caractères proches (Confusions visuelles) :}
    \begin{itemize}
        \item \zh{艺} (yì, art, radical \zh{艹} cǎo zì tóu) vs \zh{培} (péi, cultiver, radical \zh{土} tǔ zì páng).
        \item \zh{舞} (wǔ, danse, radical \zh{舛} chuǎn) vs \zh{舟} (zhōu, bateau, radical \zh{舟}) / \zh{免} (miǎn, lapin). \zh{舞} a \zh{舛} (deux pieds inversés) en bas.
        \item \zh{乐} (yuè/lè, musique/joie) vs \zh{木} (mù, bois). \zh{乐} : trait horizontal du haut, verticale, trait horizontal du milieu, crochet final (\zh{乙}) en bas. \zh{木} : trait horizontal du haut, verticale, trait horizontal du milieu, trait horizontal du bas (plus long).
        \item \zh{比} (bǐ, comparer, radical \zh{比} bǐ zì páng) vs \zh{北} (běi, nord, radical \zh{匕} bǐ). \zh{比} a un trait "huit" \zh{八} en bas (deux traits séparés) ; \zh{北} a \zh{匕} (trait horizontal + crochet).
    \end{itemize}
    \item \textbf{Pinyin / Tons (Sandhi) :}
    \begin{itemize}
        \item \zh{艺术} (yì shù, 4-4), pas *yí shu.
        \item \zh{音乐} (yīn yuè, 1-4).
        \item \zh{舞蹈} (wǔ dǎo, 3-3 $\rightarrow$ sandhi 2-3 à l'oral : \emph{wú dǎo}).
        \item \zh{喀麦隆} (Kā mài lóng, 1-4-2).
        \item \zh{雅温得} (Yǎ wēn dé, 3-1-2).
    \end{itemize}
    \item \textbf{Classificateurs (Mesure) :} Ne pas oublier le classifieur pour les objets d'art : \zh{一幅} (yī fú) pour tableaux/peintures/calligraphie, \zh{一个} (yī gè) ou \zh{一件} (yī jiàn) pour sculptures/masques, \zh{支} (zhī) pour pinceaux/bâtons de tambour.
\end{itemize}
\end{attention}

\courssub{Proposition de production et corrigé commenté}
\begin{exoresolu}[Article modèle pour le magazine « 龙声 »]
\textbf{Énoncé :} Rédiger l'article de 6 à 10 phrases suivant le plan en 4 parties, avec au moins 3 connecteurs imposés.

\tcblower

\textbf{Production modèle (Version élève visée - Niveau HSK 2/3) :}

\begin{textebox}[Mon article : 我喜欢喀麦隆艺术]
\zh{我喜欢喀麦隆艺术。} \\
\zh{首先，喀麦隆的雕塑和面具很独特，颜色丰富。} \\
\zh{然后，传统舞蹈动作有力，节奏感强，看了很开心。} \\
\zh{还有，我在雅温得博物馆看过木雕，觉得很有文化。} \\
\zh{所以，喀麦隆艺术让我骄傲，我爱我的国家文化。}
\end{textebox}

\textbf{Analyse linguistique et commentaire des choix (pour l'enseignant / auto-correction guidée) :}

\begin{enumerate}
    \item \textbf{Phrase 1 (Introduction) :} \zh{我喜欢喀麦隆艺术。} (Wǒ xǐhuan Kāmàilóng yìshù.) Structure simple SVO. Réponse directe à la question. \textit{Choix :} Position claire dès le départ.
    \item \textbf{Phrase 2 (Argument 1 - Visuel/Statuaire) :} \zh{首先，喀麦隆的雕塑和面具很独特，颜色丰富。} (Shǒuxiān, Kāmàilóng de diāosù hé miànjù hěn dútè, yánsè fēngfù.)
    \begin{itemize}
        \item Connecteur \zh{首先} en tête de phrase (thématisation).
        \item \zh{喀麦隆的} (attributif avec \zh{的}) qualifie \zh{雕塑和面具}.
        \item Adjectifs predicatifs : \zh{很独特} (très unique), \zh{颜色丰富} (couleurs riches -- structure Sujet \zh{颜色} + Adj \zh{丰富}).
        \item \textbf{Caractères cibles :} \zh{雕} (diāo, radical \zh{木}, sculpture sur bois), \zh{塑} (sù, radical \zh{月/肉}, modeler), \zh{面} (miàn, visage/masque), \zh{具} (jù, outil/objet), \zh{独} (dú, seul/unique), \zh{特} (tè, spécial), \zh{丰} (fēng, abondant), \zh{富} (fù, riche).
    \end{itemize}
    \item \textbf{Phrase 3 (Argument 2 - Danse/Performance) :} \zh{然后，传统舞蹈动作有力，节奏感强，看了很开心。} (Ránhòu, chuántǒng wǔdǎo dòngzuò yǒulì, jiézòugǎn qiáng, kàn le hěn kāixīn.)
    \begin{itemize}
        \item Connecteur \zh{然后} (suite temporelle/logique).
        \item \zh{传统舞蹈} (danse traditionnelle) -- composé nom-nom.
        \item \zh{动作有力} (mouvements puissants -- \zh{有力} adj. "qui a de la force").
        \item \zh{节奏感强} (sens du rythme fort -- \zh{感} gǎn "sentiment/sens", \zh{强} qiáng "fort").
        \item \zh{看了很开心} (regardé [et c'était] très joyeux -- structure V + \zh{了} + adj. d'état psychologique). \textbf{Point de grammaire :} \zh{了} (le) marque l'achèvement de l'action "regarder" menant au résultat "joie".
    \end{itemize}
    \item \textbf{Phrase 4 (Argument 3 - Expérience personnelle / Preuve) :} \zh{还有，我在雅温得博物馆看过木雕，觉得很有文化。} (Hái yǒu, wǒ zài Yǎwēndé bówùguǎn kàn guò mùdiāo, juéde hěn yǒu wénhuà.)
    \begin{itemize}
        \item Connecteur \zh{还有} (addition d'un argument fort : l'expérience vécue).
        \item \zh{在...看过} (expérience passée : \zh{过} guo aspect expérientiel). \textbf{Rappel :} \zh{过} se place après le verbe.
        \item \zh{雅温得} (Yǎwēndé) -- Transcription phonétique de Yaoundé (capitale politique).
        \item \zh{木雕} (mùdiāo, sculpture sur bois -- mot composé).
        \item \zh{觉得很有文化} (trouve que c'est très culturel / porteur de culture). \zh{有文化} (yǒu wénhuà) = cultivé, riche en culture.
    \end{itemize}
    \item \textbf{Phrase 5 (Conclusion) :} \zh{所以，喀麦隆艺术让我骄傲，我爱我的国家文化。} (Suǒyǐ, Kāmàilóng yìshù ràng wǒ jiāo'ào, wǒ ài wǒ de guójiā wénhuà.)
    \begin{itemize}
        \item Connecteur \zh{所以} (conséquence logique des arguments 1, 2, 3).
        \item \zh{让...骄傲} (ràng... jiāo'ào) structure causative "rendre fier". \zh{让} (ràng) = faire/laisser/rendre.
        \item \zh{我爱我的国家文化} (J'aime la culture de mon pays) -- réaffirmation identitaire.
    \end{itemize}
\end{enumerate}

\textbf{Bilan des connecteurs utilisés :} \zh{首先} (1), \zh{然后} (2), \zh{还有} (3), \zh{所以} (4) $\rightarrow$ \textbf{4 connecteurs} (critère "au moins 3" respecté avec marge).
\textbf{Nombre de phrases :} 5 phrases complexes (équivalent 8-10 propositions) $\rightarrow$ \textbf{Critère longueur respecté}.
\textbf{Caractères cibles écrits :} \zh{艺、雕、塑、面、具、舞、蹈、动、作、节、奏、感、博、物、馆、木、骄、傲、国、家}... (Réinvestissement massif du lexique de la leçon).
\end{exoresolu}

\courssub{Bilan de l'activité}
\begin{aretenir}[Ce qu'il faut retenir de cette intégration]
\begin{itemize}
    \item \textbf{Écrire, c'est choisir et structurer :} Le plan en 4 parties (Thèse -- Argument 1 -- Argument 2/Preuve -- Conclusion) guide la pensée et la langue.
    \item \textbf{Les connecteurs sont le squelette du texte :} \zh{首先/然后/还有} construisent la progression ; \zh{所以} en tire la quintessence. Leur place (souvent en position 1 ou 2 après le sujet) conditionne la fluidité.
    \item \textbf{L'écriture du caractère ancre la mémorisation :} Passer du pinyin au caractère \zh{手写} (shǒuxiě) force l'analyse morphologique (radical + phonétique) et prévient les confusions (\zh{艺/培}, \zh{舞/免}).
    \item \textbf{La correction par les pairs (pair work) :} Lire le texte d'un autre avec une grille critique développe la méta-linguistique (je repère l'erreur \zh{比较} vs \zh{比} chez l'autre, donc je ne la ferai plus).
    \item \textbf{Ancrage culturel camerounais :} Citer \zh{雅温得博物馆}, \zh{木雕}, \zh{面具}, \zh{非洲鼓} (même si non écrit ici, évoqué en classe) ancre la langue chinoise dans \textbf{notre} réalité, rendant l'expression authentique et motivante.
\end{itemize}
\end{aretenir}

\courssub{Éléments socioculturels : Regards croisés}
\begin{savaistu}[Arts traditionnels : Cameroun vs Chine]
\begin{itemize}
    \item \textbf{Masques et Sculpture :} Au Cameroun (Grassfields, Bamiléké, Bamoun), les masques (\zh{面具} miànjù) et statues royales (\zh{雕塑} diāosù) en bois (\zh{木雕} mùdiāo) ou perles symbolisent le pouvoir, les ancêtres, les sociétés secrètes. En Chine, la sculpture sur bois (\zh{木雕}), jade (\zh{玉雕} yùdiāo) ou pierre orne temples et jardins ; les masques d'opéra (\zh{脸谱} liǎnpǔ) codifient le caractère des rôles (rouge = loyal, blanc = traître, noir = intègre).
    \item \textbf{Danse et Rythme :} La danse camerounaise (\zh{舞蹈} wǔdǎo) est souvent collective, ancrée, polyrythmique (tambours \zh{鼓} gǔ, balafon), liée aux rites (récolte, funérailles, initiation). La danse classique chinoise (ex: \zh{敦煌舞} Dūnhuáng wǔ) ou l'opéra de Pékin (\zh{京剧} Jīngjù) privilégient la codification gestuelle, l'élégance (\zh{韵} yùn), la virtuosité acrobatique.
    \item \textbf{Calligraphie vs Écriture :} La calligraphie chinoise (\zh{书法} shūfǎ) est art suprême, méditation sur le trait (\zh{笔画} bǐhuà), l'encre (\zh{墨} mò), le papier (\zh{纸} zhǐ). Le Cameroun a une tradition orale forte ; l'écriture (alphabet latin, écriture Bamoun \zh{巴姆姆文字} Bāmǔmǔ wénzì) est support de mémoire, moins art graphique autonome (sauf calligraphie arabe au Nord).
    \item \textbf{Point commun :} Dans les deux cultures, l'art n'est pas seulement "beaux-arts" muséal : il est \zh{生活} (shēnghuó, vie), \zh{信仰} (xìnyǎng, croyance), \zh{身份} (shēnfèn, identité).
\end{itemize}
\end{savaistu}

\courssub{Entraînement : Thème et Version guidés}
\begin{exercice}[Entraînement à la traduction (Thème / Version)]
\textbf{A. Thème (Français $\rightarrow$ Chinois). Traduisez en caractères + pinyin.}
\begin{enumerate}
    \item Je préfère l'art chinois parce qu'il est très ancien.
    \item La sculpture camerounaise est plus unique que la peinture chinoise.
    \item D'abord, j'aime la musique africaine. Ensuite, la danse est très joyeuse.
    \item J'ai vu des masques au musée de Yaoundé, donc je connais cette culture.
    \item L'art camerounais me rend fier, j'aime mon pays.
\end{enumerate}

\textbf{B. Version (Chinois $\rightarrow$ Français). Traduisez en français.}
\begin{enumerate}
    \item 我觉得中国书法比喀麦隆绘画难。
    \item 首先，喀麦隆的木雕很有名。然后，面具的颜色很丰富。
    \item 虽然中国音乐很独特，但是我更喜欢非洲鼓。
    \item 因为我在博物馆看过雕塑，所以我觉得喀麦隆艺术很有力量。
    \item 所以，我爱喀麦隆艺术，也尊重中国文化。
\end{enumerate}
\end{exercice}

\begin{corrige}[Exercice Thème / Version]
\textbf{A. Thème (Propositions de correction) :}
\begin{enumerate}
    \item \zh{我喜欢中国艺术，因为它很古老。} (Wǒ xǐhuan Zhōngguó yìshù, yīnwèi tā hěn gǔlǎo.) \textit{Note : "préférer" $\rightarrow$ \zh{喜欢} + \zh{因为}... structure cause.}
    \item \zh{喀麦隆的雕塑比中国的绘画独特。} (Kāmàilóng de diāosù bǐ Zhōngguó de huìhuà dútè.) \textit{Note : Structure \zh{比}, pas de \zh{很}. \zh{中国的} elliptique pour \zh{中国的绘画}.}
    \item \zh{首先，我喜欢非洲音乐。然后，舞蹈很开心。} (Shǒuxiān, wǒ xǐhuan Fēizhōu yīnyuè. Ránhòu, wǔdǎo hěn kāixīn.) \textit{Note : "musique africaine" = \zh{非洲音乐}. "joyeuse" pour danse $\rightarrow$ \zh{开心} ou \zh{有活力} (yǒu huólì).}
    \item \zh{我在雅温得博物馆看过面具，所以我知道这文化。} (Wǒ zài Yǎwēndé bówùguǎn kàn guò miànjù, suǒyǐ wǒ zhīdào zhè wénhuà.) \textit{Note : \zh{看过} (expérience). \zh{这文化} (cette culture).}
    \item \zh{喀麦隆艺术让我骄傲，我爱我的国家。} (Kāmàilóng yìshù ràng wǒ jiāo'ào, wǒ ài wǒ de guójiā.) \textit{Note : \zh{让...骄傲} structure causative.}
\end{enumerate}

\textbf{B. Version (Traduction française) :}
\begin{enumerate}
    \item Je trouve que la calligraphie chinoise est plus difficile que la peinture camerounaise.
    \item D'abord, la sculpture sur bois du Cameroun est très célèbre. Ensuite, les couleurs des masques sont très riches.
    \item Bien que la musique chinoise soit unique, je préfère le tambour africain.
    \item Parce que j'ai vu des sculptures au musée, je trouve que l'art camerounais a beaucoup de force / est très puissant.
    \item Donc, j'aime l'art camerounais et je respecte aussi la culture chinoise.
\end{enumerate}
\end{corrige}

\courssub{Prolongement possible : Concours « Jeune Ambassadeur Culturel »}
\begin{experience}[Projet interdisciplinaire : Mur culturel numérique]
Les 3 articles affichés au tableau sont photographiés. La professeure crée un \emph{Padlet} ou un \emph{QR Code} vers un mur virtuel \zh{« 中喀文化墙 »}. Les élèves enregistrent leur lecture (audio) via l'application OnBuch+ ou WeChat/Whisper. La classe (et pourquoi pas une classe partenaire en Chine via jumelage) vote pour le \zh{最佳文化大使} (Meilleur Ambassadeur Culturel). Critères bonus : Prononciation (tons, \zh{erhua} si pertinent), Intensité expressive, Originalité de l'argument.
\end{experience}

\newpage

\coursec{Méthodes et exercices résolus}

\coursec{Leçon 22 — Expression écrite : culture chinoise et camerounaise (你喜欢喀麦隆还是中国艺术？为什么？)}

\courssub{Texte modèle et compréhension}

\begin{methode}[Analyser un texte modèle d'opinion comparée (版文分析)]
\textbf{Objectif :} Comprendre la structure d'un texte argumentatif simple comparant deux cultures (Cameroun / Chine) et identifier les connecteurs logiques.
\begin{enumerate}
    \item \textbf{Lire globalement :} Identifier le sujet (arts/culture), l'auteur (un élève), la thèse (préférence) et les arguments.
    \item \textbf{Repérer la structure en 3 parties :} 
    \begin{itemize}
        \item \textit{Introduction :} Présentation du sujet + thèse claire (\textbf{我觉得 / 我认为} + opinion).
        \item \textit{Développement :} 2--3 arguments reliés par des connecteurs (\textbf{首先}..., \textbf{其次}..., \textbf{而且}...). Chaque argument = \textbf{Fait culturel} + \textbf{Sentiment personnel}.
        \item \textit{Conclusion :} Résumé + ouverture (\textbf{总之}..., \textbf{希望}...).
    \end{itemize}
    \item \textbf{Surligner les mots-outils de liaison :} \cle{因为 / 所以} (cause-conséquence), \cle{虽然 / 但是} (concession), \cle{和 / 还是} (choix/alternative), \cle{比} (comparaison).
    \item \textbf{Noter le vocabulaire clé du thème :} Noms d'arts, adjectifs d'appréciation, verbes d'opinion.
\end{enumerate}
\textbf{Réflexe Bac :} Avant d'écrire, faites un plan en 3 lignes (Intro / Arg 1, Arg 2 / Conclu) avec 1 connecteur par flèche.
\end{methode}

\begin{exoresolu}[Analyse du texte modèle : « Mon cœur balance »]
\textbf{Énoncé :} Lisez le texte ci-dessous. Répondez aux questions en français.
\begin{textebox}[Texte modèle : 你喜欢喀麦隆还是中国艺术？]
我叫马里，是喀麦隆的高中生。我喜欢中国艺术，也喜欢喀麦隆艺术。\\
\textbf{首先}，中国书法很有意思。写汉字要一横一竖，很安静，让我心静。\\
\textbf{其次}，喀麦隆的雕刻很有特色。木头雕刻的人和动物，形神兼备，充满生命力。\\
\textbf{虽然}中国书法很难，\textbf{但是}喀麦隆雕刻也需要耐心。\\
\textbf{总之}，我最喜欢中国书法，\textbf{因为}它能培养我的专注力。将来我想去北京学书法。
\end{textebox}
\begin{enumerate}
    \item Quel est le sujet du texte ? Quelle est la thèse de l'auteur ?
    \item Relevez les 4 connecteurs structurants (en caractères) et donnez leur fonction.
    \item Trouvez 3 mots de vocabulaire « Arts » et 2 adjectifs d'appréciation.
    \item Traduisez la phrase : « 虽然中国书法很难，但是喀麦隆雕刻也需要耐心。 »
\end{enumerate}
\tcblower
\textbf{Solution détaillée :}
\begin{enumerate}
    \item \textbf{Sujet :} Comparaison entre l'art chinois (calligraphie) et l'art camerounais (sculpture sur bois). \textbf{Thèse :} L'auteur aime les deux, mais préfère la calligraphie chinoise (我最喜欢中国书法).
    \item \textbf{Connecteurs :}
    \begin{center}
    \begin{tabularx}{\linewidth}{|X|X|X|}
    \hline \rowcolor{popblueL} Connecteur & Pinyin & Fonction \\ \hline
    首先 & shǒuxiān & Énumération (Argument 1) \\ \hline
    其次 & qícì & Énumération (Argument 2) \\ \hline
    虽然...但是... & suīrán... dànshì... & Concession / Opposition \\ \hline
    总之 & zǒngzhī & Conclusion / Résumé \\ \hline
    \end{tabularx}
    \end{center}
    \item \textbf{Vocabulaire Arts :} \zh{书法} (shūfǎ, calligraphie), \zh{雕刻} (diāokè, sculpture/gravure), \zh{汉字} (hànzì, sinogramme). \textbf{Adjectifs :} \zh{有意思} (yǒu yìsi, intéressant), \zh{有特色} (yǒu tèsè, typique/avec caractère), \zh{充满生命力} (chōngmǎn shēngmìnglì, plein de vitalité).
    \item \textbf{Traduction :} « Bien que la calligraphie chinoise soit difficile, la sculpture camerounaise demande aussi de la patience. » \\
    \textit{Analyse grammaticale :} \zh{虽然} (bien que) introduit la subordonnée de concession (sujet \zh{中国书法} + adjectif \zh{难}). \zh{但是} (mais) introduit la principale (sujet \zh{喀麦隆雕刻} + \zh{也} aussi + \zh{需要} avoir besoin + \zh{耐心} patience). Attention : en chinois, \zh{虽然} et \zh{但是} s'utilisent \textbf{ensemble} (corrélatif), on ne les sépare pas.
\end{enumerate}
\textbf{Conclusion :} Ce texte respecte la structure standard : Thèse $\rightarrow$ Arguments (connecteurs ordinaux) $\rightarrow$ Nuance (concession) $\rightarrow$ Conclusion justifiée (\zh{因为}).
\end{exoresolu}

\courssub{Lexique du thème : arts et opinions}

\begin{methode}[Construire son lexique thématique : Arts, Opinions, Comparaison (主题词汇构建)]
\textbf{Objectif :} Maîtriser le vocabulaire de fréquence élevée (HSK 2-3) pour décrire des œuvres, exprimer une préférence et comparer.
\begin{enumerate}
    \item \textbf{Classer par catégories mentales (Mindmap) :}
    \begin{itemize}
        \item \textit{Arts (名词) :} 书法, 绘画, 雕刻, 音乐, 戏曲, 舞蹈, 京剧, 面具, 鼓, 琵琶, 古筝.
        \item \textit{Verbes d'action/perception :} 看, 听, 学, 写, 画, 雕, 表演, 欣赏.
        \item \textit{Adjectifs d'appréciation (褒义/贬义) :} \textbf{Positif :} 美, 漂亮, 有意思, 有特色, 优美, 精彩, 深刻, 传统, 现代. \textbf{Négatif :} 难, 无聊, 奇怪, 复杂.
        \item \textit{Structures d'opinion :} 我喜欢 / 爱 / 觉得...好 / 认为...好 / 对...感兴趣 / 偏爱.
        \item \textit{Comparaison :} 比, 跟...一样, 不如, 最喜欢, 更.
    \end{itemize}
    \item \textbf{Mémoriser les radicaux-clés (写字线索) :}
    \begin{itemize}
        \item \zh{书} (shū, livre/écrire) : clé \zh{乙} (yǐ) --\textgreater\ 书法, 书店, 书包.
        \item \zh{画} (huà, dessin/peinture) : clé \zh{田} (tián) --\textgreater\ 绘画, 画家, 画展.
        \item \zh{雕} (diāo, sculpter) : clé \zh{石} (shí, pierre) --\textgreater\ 雕刻, 石雕, 木雕.
        \item \zh{艺} (yì, art/talent) : clé \zh{艹} (cǎo) --\textgreater\ 艺术, 艺术家, 园艺.
    \end{itemize}
    \item \textbf{Attention aux « faux amis » / homophones :} \zh{艺} (yì, art) $\neq$ \zh{易} (yì, facile/changer) $\neq$ \zh{意} (yì, sens/idée). \zh{雕} (diāo, sculpter) $\neq$ \zh{吊} (diào, pendre).
\end{enumerate}
\textbf{Astuce :} Apprenez les mots par \textbf{binômes opposés} (传统 vs 现代, 简单 vs 复杂, 安静 vs 热闹) et par \textbf{familles de radicaux}.
\end{methode}

\begin{exoresolu}[Exercices de lexique : Complétion, Caractères proches, Traduction]
\textbf{Énoncé :}
\begin{enumerate}
    \item \textbf{Complétez avec le bon caractère (radical indice) :}
    \begin{itemize}
        \item \zh{我\_\_\_ (clé \zh{艹}) 术很感兴趣。}
        \item \zh{这幅\_\_\_ (clé \zh{田}) 很漂亮。}
        \item \zh{他学习\_\_\_ (clé \zh{石}) 已经三年了。}
    \end{itemize}
    \item \textbf{Choisissez l'adjectif approprié (有意思 / 有特色 / 优美 / 无聊) :}
    \begin{itemize}
        \item Le masque traditionnel est très \_\_\_\_\_ (représentatif de la culture).
        \item Ce film est \_\_\_\_\_ (je m'endors).
        \item La musique du \zh{古筝} est très \_\_\_\_\_ (agréable à l'oreille, élégante).
        \item Apprendre les caractères, c'est \_\_\_\_\_ (ça passionne).
    \end{itemize}
    \item \textbf{Thème (Français $\rightarrow$ Chinois) :}
    \begin{itemize}
        \item « Je pense que l'opéra de Pékin est plus intéressant que le football. »
        \item « La sculpture sur bois camerounaise a beaucoup de caractère. »
    \end{itemize}
\end{enumerate}
\tcblower
\textbf{Solution détaillée :}
\begin{enumerate}
    \item \textbf{Caractères (radicaux) :}
    \begin{itemize}
        \item \zh{艺} (yì) --\textgreater\ \zh{艺术} (yìshù). Clé \zh{艹} (herbe/plante, haut).
        \item \zh{画} (huà) --\textgreater\ \zh{画} (huà, tableau/peinture). Clé \zh{田} (champ, milieu).
        \item \zh{雕} (diāo) --\textgreater\ \zh{雕刻} (diāokè). Clé \zh{石} (pierre, gauche). \textbf{Piège :} Ne pas confondre avec \zh{吊} (diào, clé \zh{口} bouche).
    \end{itemize}
    \item \textbf{Adjectifs :}
    \begin{itemize}
        \item \zh{有特色} (yǒu tèsè) -- « typique, distinctif » (masque traditionnel).
        \item \zh{无聊} (wúliáo) -- « ennuyeux » (film).
        \item \zh{优美} (yōuměi) -- « gracieux, esthétique » (musique, danse, écriture).
        \item \zh{有意思} (yǒu yìsi) -- « intéressant, amusant » (apprentissage).
    \end{itemize}
    \item \textbf{Thème :}
    \begin{itemize}
        \item \zh{我认为京剧比足球有意思。} (Wǒ rènwéi Jīngjù bǐ zúqiú yǒu yìsi.) \\ Structure : Sujet + \zh{认为} + A + \zh{比} + B + Adj. \textbf{Note :} \zh{足球} (zúqiú) = football.
        \item \zh{喀麦隆的木雕很有特色。} (Kāmàilóng de mùdiāo hěn yǒu tèsè.) \\ \zh{木雕} (mùdiāo) = sculpture sur bois (bois \zh{木} + sculpter \zh{雕}). \zh{很有特色} = « très typique / a beaucoup de caractère ».
    \end{itemize}
\end{enumerate}
\textbf{Conclusion :} La maîtrise des radicaux (\zh{艹, 田, 石}) permet de deviner le sens et d'écrire correctement. Le choix de l'adjectif dépend du registre (esthétique \zh{优美} vs culturel \zh{有特色} vs ludique \zh{有意思}).
\end{exoresolu}

\begin{center}
\begin{tabularx}{\linewidth}{|X|X|X|X|}
\hline \rowcolor{popblueL} Caractères & Pinyin & Nature & Traduction \\ \hline
书法 & shūfǎ & nom & calligraphie \\ \hline
绘画 & huìhuà & nom & peinture \\ \hline
雕刻 & diāokè & nom/verbe & sculpture, graver \\ \hline
木雕 & mùdiāo & nom & sculpture sur bois \\ \hline
面具 & miànjù & nom & masque \\ \hline
京剧 & jīngjù & nom & opéra de Pékin \\ \hline
古筝 & gǔzhēng & nom & cithare chinoise \\ \hline
巴拉风 & bālāfēng & nom & balafon \\ \hline
有意思 & yǒu yìsi & adj & intéressant \\ \hline
有特色 & yǒu tèsè & adj & typique, distinctif \\ \hline
优美 & yōuměi & adj & gracieux, esthétique \\ \hline
充满生命力 & chōngmǎn shēngmìnglì & loc. adj & plein de vitalité \\ \hline
培养 & péiyǎng & verbe & cultiver, développer \\ \hline
专注力 & zhuānzhùlì & nom & concentration \\ \hline
文化内涵 & wénhuà nèihán & nom & contenu culturel, profondeur culturelle \\ \hline
\end{tabularx}
\end{center}

\courssub{Structure du texte et connecteurs logiques}

\begin{methode}[Maîtriser les structures de comparaison et d'argumentation (比较与论证句式)]
\textbf{Objectif :} Produire des phrases correctes pour comparer (A vs B) et justifier (Cause $\rightarrow$ Conséquence / Concession).
\begin{enumerate}
    \item \textbf{Structure comparative \cle{比} (bǐ) :} \zh{主语 A + 比 + 主语 B + Adj (sans 很)}.
    \begin{itemize}
        \item Négation : \zh{A 没有 B + Adj} (A n'est pas aussi Adj que B).
        \item Égalité : \zh{A 跟/和 B + 一样 + Adj}.
        \item Superlatif : \zh{最 + Adj} (le plus) -- pas de \zh{比}.
        \item Degré : \zh{A 比 B Adj + 得多 / 多了} (beaucoup plus), \zh{A 比 B Adj + 一点儿 / 些} (un peu plus).
    \end{itemize}
    \item \textbf{Structure d'opinion nuancée :} \zh{主语 + 觉得 / 认为 / 觉得 + 从句 (Sujet + Verbe/Adj)}.
    \item \textbf{Connecteurs logiques indispensables (ordre du texte) :}
    \begin{center}
    \begin{tabularx}{\linewidth}{|X|X|X|}
    \hline \rowcolor{popblueL} Fonction & Connecteurs (Caractères + Pinyin) & Place dans la phrase \\ \hline
    Énumérer & \zh{首先} (shǒuxiān), \zh{其次} (qícì), \zh{最后} (zuìhòu) & Début de phrase (sujet après) \\ \hline
    Ajouter & \zh{而且} (érqiě), \zh{也} (yě), \zh{还} (hái) & Devant V/Adj (après sujet) \\ \hline
    Cause & \zh{因为} (yīnwèi) & Début phrase ou après sujet \\ \hline
    Conséquence & \zh{所以} (suǒyǐ) & Début phrase (après \zh{因为}...) \\ \hline
    Concession & \zh{虽然} (suīrán) ... \zh{但是/可是} (dànshì/kěshì) & \zh{虽然} clause 1, \zh{但是} clause 2 \\ \hline
    Conclusion & \zh{总之} (zǒngzhī), \zh{因此} (yīncǐ) & Début phrase \\ \hline
    \end{tabularx}
    \end{center}
    \item \textbf{Règle d'or ordre des mots (Compléments) :} \zh{主语 + 时间 + 地点 + 方式/工具 + 动词 + 补语}. \textbf{Jamais} : \zh{我去北京学书法明年} $\rightarrow$ \zh{我明年去北京学书法}.
\end{enumerate}
\textbf{Réflexe :} Pour chaque argument, posez la question « Pourquoi ? » $\rightarrow$ ajoutez \zh{因为}...\zh{所以}...
\end{methode}

\begin{exoresolu}[Transformation de phrases et Mini-composition guidée]
\textbf{Énoncé :}
\begin{enumerate}
    \item \textbf{Transformez en utilisant \zh{比} (corrigez l'erreur s'il y a) :}
    \begin{itemize}
        \item \zh{中国书法很难，喀麦隆雕刻也很难。} (Égalité)
        \item \zh{我觉得中国音乐比喀麦隆音乐更好听很。} (Corrigez)
        \item \zh{我不喜欢现代艺术，我喜欢传统艺术。} (Préférence avec \zh{比})
    \end{itemize}
    \item \textbf{Reliez les deux phrases par le connecteur indiqué :}
    \begin{itemize}
        \item Phrase 1: \zh{学书法很安静。} Phrase 2: \zh{我心静。} (Cause/Conséquence : \zh{因为}...\zh{所以}...)
        \item Phrase 1: \zh{汉字很难写。} Phrase 2: \zh{我每天练习。} (Concession : \zh{虽然}...\zh{但是}...)
    \end{itemize}
    \item \textbf{Rédigez un paragraphe de 5 phrases (Introduction + 2 arguments + Concession + Conclusion) sur le thème :} « Préférez-vous la musique traditionnelle chinoise (guzheng, erhu) ou camerounaise (balafon, tam-tam) ? Pourquoi ? » Utilisez au moins 3 connecteurs du tableau.
\end{enumerate}
\tcblower
\textbf{Solution détaillée :}
\begin{enumerate}
    \item \textbf{Comparaison :}
    \begin{itemize}
        \item \zh{中国书法跟喀麦隆雕刻一样难。} (\zh{跟...一样} + Adj, pas de \zh{很}).
        \item \zh{我觉得中国音乐比喀麦隆音乐好听。} (\textbf{Erreur corrigée :} \zh{更} et \zh{很} sont interdits après l'adjectif dans la structure \zh{比}. On dit juste \zh{好听}).
        \item \zh{我比较喜欢传统艺术，不太喜欢现代艺术。} (Ou : \zh{我觉得传统艺术比现代艺术有意思。}) -- \zh{比较} (bǐjiào) = relativement/plutôt, adverbe de degré devant le verbe.
    \end{itemize}
    \item \textbf{Connecteurs :}
    \begin{itemize}
        \item \zh{因为学书法很安静，所以我心静。} (Ou : \zh{学书法很安静，所以我心静。} -- \zh{因为} optionnel à l'oral, obligatoire à l'écrit formel).
        \item \zh{虽然汉字很难写，但是我每天练习。} (\textbf{Ordre fixe :} \zh{虽然} clause 1, \zh{但是} clause 2. Sujet \zh{我} répété ou omis si identique ? Ici sujet différent \zh{汉字} vs \zh{我}, donc garder les deux).
    \end{itemize}
    \item \textbf{Production écrite modèle (Corrigé type Bac) :}
    \begin{textebox}[Production : 我更喜欢喀麦隆音乐]
    我叫保罗，我学中文两年了。\\
    \textbf{首先}，喀麦隆的传统音乐很热闹，\textbf{比如}巴拉风和鼓，节奏感强，让人想跳舞。\\
    \textbf{其次}，中国古筝音乐很优美，很安静，适合冥想。\\
    \textbf{虽然}古筝很优美，\textbf{但是}我更喜欢喀麦隆的巴拉风，\textbf{因为}它代表我的文化根源，充满生命力。\\
    \textbf{总之}，两种音乐都有特色，但我偏爱喀麦隆音乐。
    \end{textebox}
    \textit{Analyse :} Vocabulaire HSK 3 (\zh{节奏感}, \zh{冥想}, \zh{文化根源}, \zh{偏爱}). Connecteurs : \zh{首先/其次/虽然/但是/因为/总之}. Structure \zh{比} implicite via \zh{更喜欢}. Respect ordre mots (Temps/Lieu/Manière).
\end{enumerate}
\end{exoresolu}

\courssub{Écriture des caractères du thème}

\begin{methode}[Écrire les caractères du thème : Ordre des traits, Radicaux, Différenciation (规范书写)]
\textbf{Objectif :} Écrire correctement et rapidement les 15--20 caractères clés de la leçon (HSK 2-3) pour l'épreuve écrite.
\begin{enumerate}
    \item \textbf{Respecter l'ordre des traits universel :} 
    \begin{enumerate}
        \item Horizontal avant Vertical (\zh{十} : \zh{一} puis \zh{丨}).
        \item Gauche avant Droite (\zh{好} : \zh{女} puis \zh{子}).
        \item Haut avant Bas (\zh{三} : trait haut, milieu, bas).
        \item Extérieur avant Intérieur (\zh{国} : \zh{囗} puis \zh{玉}).
        \item Fermer en dernier (\zh{回} : \zh{囗} ferme à la fin).
        \item Trait central avant ailes symétriques (\zh{小}, \zh{水}).
    \end{enumerate}
    \item \textbf{Identifier le radical (clé de dictionnaire) et le phonétique :}
    \begin{itemize}
        \item \zh{书} (shū) : Radical \zh{乙} (yǐ, 2 traits). Picto-idéogramme.
        \item \zh{画} (huà) : Radical \zh{田} (tián, 5 traits). Structure Haut-Bas.
        \item \zh{雕} (diāo) : Radical \zh{石} (shí, 5 traits). Structure Gauche-Droite. Phonétique \zh{周} (zhōu).
        \item \zh{艺} (yì) : Radical \zh{艹} (cǎo, 3 traits). Structure Haut-Bas. \textbf{Attention :} Écrire \zh{艹} (3 traits : horiz, horiz, vertical) pas \zh{草} (6 traits) en haut.
        \item \zh{乐} (yuè/yào) : Radical \zh{丿} (piě). 5 traits. \textbf{Piège :} Ne pas confondre avec \zh{木} (mù, 4 traits) ou \zh{未} (wèi).
    \end{itemize}
    \item \textbf{Exercer les paires de caractères proches (易混字) :}
    \begin{itemize}
        \item \zh{艺} (yì, art) vs \zh{培} (péi, cultiver) -- Radicaux \zh{艹} vs \zh{土}.
        \item \zh{雕} (diāo, sculpter) vs \zh{吊} (diào, pendre) -- Radicaux \zh{石} vs \zh{口}.
        \item \zh{特} (tè, spécial) vs \zh{持} (chí, tenir) -- \zh{特} = \zh{牛} + \zh{寺}, \zh{持} = \zh{扌} + \zh{寺}.
        \item \zh{优} (yōu, excellent) vs \zh{忧} (yōu, soucis) -- Radical \zh{亻} (personne) pour \zh{优} vs Radical \zh{忄} (cœur) pour \zh{忧}.
    \end{itemize}
\end{enumerate}
\textbf{Entraînement :} Écrire 5 fois chaque caractère en dictée silencieuse (dire le nom des traits : \zh{横、竖、撇、捺、折}).
\end{methode}

\begin{exoresolu}[Dictée de caractères, Correction d'erreurs d'écriture, Recherche au dictionnaire]
\textbf{Énoncé :}
\begin{enumerate}
    \item \textbf{Ordre des traits :} Numérotez les traits pour \zh{艺} (yì) et \zh{雕} (diāo).
    \item \textbf{Correction :} Un élève a écrit : \zh{我最喜欢中国书\underline{泛}。} et \zh{喀麦隆的木\underline{吊}很有特色。} Corrigez les deux caractères soulignés et expliquez l'erreur (radical).
    \item \textbf{Recherche dictionnaire :} Pour trouver \zh{雕} (diāo) dans un dictionnaire papier (classement par radicaux) : Quel radical cherchez-vous ? Combien de traits pour le radical ? Combien de traits pour le reste du caractère (corps) ?
    \item \textbf{Mini-dictée (Pinyin $\rightarrow$ Caractères) :} Écrivez en caractères : \zh{wǒ rènwéi zhōngguó shūfǎ bǐ kāmàilóng de mùdiāo nán.} (Je pense que la calligraphie chinoise est plus difficile que la sculpture sur bois camerounaise.)
\end{enumerate}
\tcblower
\textbf{Solution détaillée :}
\begin{enumerate}
    \item \textbf{Ordre des traits (description) :}
    \begin{itemize}
        \item \zh{艺} (7 traits) : 1.\zh{艹} (横、横、竖) -- 4.\zh{云} (横、撇、横折弯钩、点). \textit{Total 7 traits.}
        \item \zh{雕} (16 traits) : Gauche \zh{石} (5 traits : 横、撇、竖、横折、横). Droite \zh{周} (11 traits : 竖、横折、横、横、竖、横折、横、横、撇、点、横). \textit{Total 16 traits.}
    \end{itemize}
    \item \textbf{Corrections :}
    \begin{itemize}
        \item \zh{泛} (fàn, déborder) $\rightarrow$ \zh{法} (fǎ, loi/méthode). \textbf{Erreur :} Même radical \zh{氵} (eau). Différence partie droite : \zh{去} (qù) pour \zh{法} vs \zh{凡} (fán) pour \zh{泛}. Confusion phonétique (fǎ vs fàn) et graphique.
        \item \zh{吊} (diào, pendre) $\rightarrow$ \zh{雕} (diāo, sculpter). \textbf{Erreur :} Radical \zh{口} (bouche) au lieu de \zh{石} (pierre). Même phonétique \zh{周}, sens radical différent. \zh{木雕} (sculpture bois) correct.
    \end{itemize}
    \item \textbf{Dictionnaire :} Radical \zh{石} (Shí, 5 traits). Corps \zh{周} (11 traits). Chercher dans l'index 5 traits $\rightarrow$ \zh{石} $\rightarrow$ compter 11 traits restants.
    \item \textbf{Dictée :} \zh{我认为中国书法比喀麦隆的木雕难。} \\
    \textit{Points de vigilance :} \zh{认为} (pas \zh{为}), \zh{书法} (pas \zh{泛}), \zh{比}, \zh{喀麦隆} (noms propres), \zh{木雕} (pas \zh{吊}), \zh{难} (nán, 2e ton).
\end{enumerate}
\end{exoresolu}

\courssub{Thème et version guidés}

\begin{methode}[Réussir le Thème (Français $\rightarrow$ Chinois) et la Version (Chinois $\rightarrow$ Français) (翻译策略)]
\textbf{Objectif :} Appliquer une méthode rigoureuse pour traduire des phrases simples ou un court paragraphe sans calque du français.
\begin{enumerate}
    \item \textbf{THÈME (Fr $\rightarrow$ Zh) -- Méthode en 4 étapes :}
    \begin{enumerate}
        \item \textbf{Analyser la phrase française :} Identifier Sujet, Verbe (temps), Compléments (COI, CCL, CCM), Négation, Modalité.
        \item \textbf{Poser l'ossature chinoise (Ordre SUJET -- TEMPS -- LIEU -- MANIÈRE -- VERBE -- COMPLÉMENT) :} Ne pas garder l'ordre français (SVO mais compléments avant V).
        \item \textbf{Sélectionner le vocabulaire précis :} Attention aux verbes « avoir » (\zh{有} vs \zh{是} vs \zh{在}), « être » (\zh{是} vs \zh{很} + Adj), « aimer » (\zh{喜欢} vs \zh{爱} vs \zh{感兴趣}).
        \item \textbf{Vérifier :} Tons (pinyin), Caractères (radicaux), Mots-outils (\zh{了} aspect accompli, \zh{过} expérience, \zh{的/得/地}), Classificateurs (\zh{个/幅/首/支}).
    \end{enumerate}
    \item \textbf{VERSION (Zh $\rightarrow$ Fr) -- Méthode en 3 étapes :}
    \begin{enumerate}
        \item \textbf{Segmenter (切分) :} Couper en groupes syntaxiques (Sujet / Prédicat / Compléments) -- repérer les \zh{的} (nominalisation), \zh{得} (complément de degré), \zh{了/过/着}.
        \item \textbf{Traduire mot à mot (littéral) :} Garder l'ordre chinois pour voir la logique.
        \item \textbf{Reformuler en français correct :} Adapter temps, prépositions, articles. \textbf{Ne pas traduire \zh{很} par « très » si c'est juste un liant !} (Ex: \zh{很好} = « c'est bien / il est bien », pas « il est très bien »).
    \end{enumerate}
    \item \textbf{Points de grammaire pièges (Thème/Version) :}
    \begin{itemize}
        \item \zh{比} structure : Fr « A est plus adj que B » $\rightarrow$ Zh « A \zh{比} B Adj » (pas de \zh{更/很} après Adj).
        \item \zh{把} structure : Fr « Je prends le livre » (action sur objet défini) $\rightarrow$ Zh \zh{我把书拿走}.
        \item \zh{被} passif : Fr « Le livre a été lu par moi » $\rightarrow$ Zh \zh{书被我看了}.
        \item Aspect \zh{了} (accompli/passé proche) vs \zh{过} (expérience vécue) vs \zh{着} (duratif).
    \end{itemize}
\end{enumerate}
\end{methode}

\begin{exoresolu}[Thème et Version guidés : Textes courts « Culture \& Arts »]
\textbf{Énoncé :}
\begin{enumerate}
    \item \textbf{THÈME :} Traduisez en chinois (caractères + pinyin).
    \begin{enumerate}
        \item « L'année dernière, je suis allé à Yaoundé voir une exposition de masques. »
        \item « La calligraphie chinoise est plus difficile que la peinture, mais je l'aime beaucoup. »
        \item « Si vous avez le temps, venez au Cameroun goûter le ndolé et écouter le balafon. »
    \end{enumerate}
    \item \textbf{VERSION :} Traduisez en français naturel.
    \begin{enumerate}
        \item \zh{虽然学中文不容易，但是我很有成就感。}
        \item \zh{这幅山水画画得很有意思，颜色也很鲜艳。}
        \item \zh{我想请你去北京的画展，那里有很多著名的艺术家的作品。}
    \end{enumerate}
\end{enumerate}
\tcblower
\textbf{Solution détaillée :}
\begin{enumerate}
    \item \textbf{THÈME :}
    \begin{enumerate}
        \item \textbf{Analyse :} Sujet \zh{我}, Temps \zh{去年} (l'an dernier), Lieu \zh{雅温得} (Yaoundé), Verbe \zh{去} + But \zh{看} + Objet \zh{展览} (exposition) + \zh{面具} (masques). Classificateur \zh{个} ou \zh{场} pour expo.
        \textbf{Chinois :} \zh{去年我去雅温得看了一场面具展览。} (Qùnián wǒ qù Yǎwēndé kàn le yī chǎng miànjù zhǎnlǎn.) \\
        \textit{Note :} \zh{了} après \zh{看} (action accomplie). \zh{一场} (classificateur événement). \zh{面具展览} (nom composé).
        \item \textbf{Analyse :} Comparaison \zh{比}. A=\zh{中国书法}, B=\zh{绘画}, Adj=\zh{难}. Concession \zh{虽然}...\zh{但是}... Degré \zh{很喜欢}.
        \textbf{Chinois :} \zh{虽然中国书法比绘画难，但是我非常喜欢。} (Suīrán Zhōngguó shūfǎ bǐ huìhuà nán, dànshì wǒ fēicháng xǐhuan.) \\
        \textit{Attention :} Pas de \zh{很} après \zh{难} dans structure \zh{比}. \zh{非常} (fēicháng) = très (adverbe degré devant V/Adj).
        \item \textbf{Analyse :} Condition \zh{如果}...\zh{就}... Verbes série : \zh{来} (venir) + \zh{尝} (goûter) + \zh{听} (écouter). Objets : \zh{恩多莱} (ndolé), \zh{巴拉风} (balafon).
        \textbf{Chinois :} \zh{如果你有时间，就来喀麦隆尝尝恩多莱，听听巴拉风吧。} (Rúguǒ nǐ yǒu shíjiān, jiù lái Kāmàilóng chángchang Ēnduōlái, tīngtīng Bālāfēng ba.) \\
        \textit{Reduplication :} \zh{尝尝} (goûter un peu), \zh{听听} (écouter un peu) -- adoucit le ton. \zh{吧} (particule suggestion).
    \end{enumerate}
    \item \textbf{VERSION :}
    \begin{enumerate}
        \item \zh{虽然} (Bien que) \zh{学中文} (apprendre le chinois) \zh{不容易} (pas facile), \zh{但是} (mais) \zh{我} (je) \zh{很有成就感} (ai beaucoup sentiment d'accomplissement).
        \textbf{Français :} « Bien qu'apprendre le chinois ne soit pas facile, j'ai un grand sentiment d'accomplissement. » (Ou : « ...je me sens très accompli. »)
        \item \zh{这幅} (Ce [classif tableau]) \zh{山水画} (peinture paysage) \zh{画得} (peint - \zh{得} complém. degré) \zh{很有意思} (très intéressant), \zh{颜色} (couleurs) \zh{也} (aussi) \zh{很鲜艳} (très vives/éclatantes).
        \textbf{Français :} « Ce tableau de paysage est peint de manière très intéressante, les couleurs sont aussi très vives. » (\zh{画得} structure V+\zh{得}+Adj : manière de peindre).
        \item \zh{我} (Je) \zh{想请} (veux inviter) \zh{你} (toi) \zh{去} (aller) \zh{北京的画展} (expo de Pékin), \zh{那里} (là-bas) \zh{有} (il y a) \zh{很多} (beaucoup) \zh{著名的} (famos) \zh{艺术家} (artistes) \zh{的} (de) \zh{作品} (œuvres).
        \textbf{Français :} « Je voudrais t'inviter à aller à l'exposition de Pékin, il y a là-bas beaucoup d'œuvres d'artistes célèbres. »
    \end{enumerate}
\end{enumerate}
\end{exoresolu}

\courssub{Production écrite et évaluation}

\begin{methode}[Produire une expression écrite complète (Épreuve type Bac : 80--100 caractères) (书面表达全流程)]
\textbf{Objectif :} Rédiger un texte cohérent, structuré, lexicalement riche et grammaticalement correct en 20--25 minutes.
\begin{enumerate}
    \item \textbf{Analyse du sujet (2 min) :} 
    \begin{itemize}
        \item Sujet : « 你喜欢喀麦隆还是中国艺术？为什么？ » (Comparaison + Choix + Justification).
        \item Contraintes : \textbf{Tu} (je), \textbf{Arts} (vocabulaire précis), \textbf{Pourquoi} (arguments + \zh{因为}), \textbf{Longueur} (~80--100 caractères).
    \end{itemize}
    \item \textbf{Planification / Mindmap (5 min) :} 
    \begin{center}
    \begin{tikzpicture}[mindmap, grow cyclic, every node/.style=concept, concept color=popblue!80,
        level 1/.append style={level distance=3.5cm, sibling angle=120},
        level 2/.append style={level distance=2.5cm}]
    \node {Sujet : Art Cameroun vs Chine}
        child[concept color=popgreen!80] {node {Intro : Moi + Contexte}
            child {node {Nom, Classe, Pays}}
        }
        child[concept color=poporange!80] {node {Choix : 中国书法}
            child {node {Arg 1: 静心 \zh{安静}}}
            child {node {Arg 2: Culture \zh{汉字}}}
        }
        child[concept color=poppink!80] {node {Concession: 喀麦隆雕刻 \zh{有特色}}
            child {node {Mais mon cœur: 中国}}
        };
    \end{tikzpicture}
    \end{center}
    \item \textbf{Rédaction (12 min) :} Écrire en suivant le plan. \textbf{Compter les caractères} (viser 90). Utiliser les connecteurs préparés (\zh{首先}, \zh{其次}, \zh{虽然}...\zh{但是}, \zh{因为}...\zh{所以}, \zh{总之}).
    \item \textbf{Relecture / Auto-correction (6 min) -- Check-list « Zéro Faute » :}
    \begin{itemize}
        \item [$\square$] \textbf{Tons / Pinyin :} Pas de ton neutre mal placé, \zh{一} (yī/yí/yì), \zh{不} (bù/bú).
        \item [$\square$] \textbf{Caractères :} Radicaux corrects (\zh{艺} vs \zh{培}, \zh{雕} vs \zh{吊}, \zh{法} vs \zh{泛}).
        \item [$\square$] \textbf{Ordre des mots :} Temps/Lieu avant Verbe. \zh{的/得/地} distingués.
        \item [$\square$] \textbf{Grammaire :} \zh{比} sans \zh{很}. \zh{了/过} aspect. Classificateurs (\zh{幅画}, \zh{首曲}, \zh{个雕塑}).
        \item [$\square$] \textbf{Ponctuation :} 。 ， ？ ！ (pleine chasse). Pas d'espace.
    \end{itemize}
\end{enumerate}
\textbf{Astuce temps :} Si bloqué sur un caractère $\rightarrow$ écrire le pinyin entre parenthèses et continuer. Revenir à la fin.
\end{methode}

\begin{exoresolu}[Production écrite intégrale : Sujet type Bac + Correction annotée]
\textbf{Énoncé :} « Vous êtes un élève de 1ère A. Écrivez un court texte (80--100 caractères chinois) pour répondre à la question : \textbf{你喜欢喀麦隆还是中国艺术？为什么？} »
\tcblower
\textbf{Production élève (Modèle visé ~95 caractères) :}
\begin{textebox}[Ma réponse : 我爱中国书法]
我叫张伟，在雅温得上学。我喜欢喀麦隆雕刻，\textbf{但是}我更爱中国书法。\\
\textbf{首先}，写书法很安静，能让我专心。\\
\textbf{其次}，汉字很有文化内涵，一笔一画都有意思。\\
\textbf{虽然}喀麦隆木雕很有生命力，\textbf{但是}书法能培养耐心。\\
\textbf{因为}我想提高专注力，\textbf{所以}我每天练习书法。\\
\textbf{总之}，中国书法是我的最爱。
\end{textebox}
\textbf{Analyse détaillée (Grille d'évaluation) :}
\begin{center}
\begin{tabularx}{\linewidth}{|X|X|}
\hline \rowcolor{popblueL} Critère & Évaluation \\ \hline
\textbf{Contenu / Structure} & Intro (Nom/Lieu) $\rightarrow$ Thèse claire (\zh{更爱}) $\rightarrow$ 2 Args (\zh{首先/其次}) $\rightarrow$ Concession (\zh{虽然}...\zh{但是}) $\rightarrow$ Cause/Effet (\zh{因为}...\zh{所以}) $\rightarrow$ Conclu (\zh{总之}). \textbf{Parfait.} \\ \hline
\textbf{Vocabulaire (HSK 3+)} & \zh{文化内涵} (contenu culturel), \zh{一笔一画} (trait par trait, chengyu), \zh{培养} (cultiver), \zh{专注力} (concentration), \zh{生命力} (vitalité). \textbf{Riche \& Précis.} \\ \hline
\textbf{Grammaire} & \zh{比} implicite (\zh{更爱}). \zh{能} (pouvoir/capacité). \zh{的} (attributif \zh{文化内涵}, \zh{木雕}). \zh{得} absent (pas de complém. degré). \zh{了} absent (habitude \zh{每天}). \textbf{Correct.} \\ \hline
\textbf{Caractères / Ponctuation} & \zh{伟} (radical \zh{亻}), \zh{内} (radical \zh{冂}), \zh{涵} (radical \zh{氵}), \zh{雕} (radical \zh{石}). Ponctuation pleine chasse. \textbf{Standard.} \\ \hline
\textbf{Nombre de caractères} & $\approx$ 92 caractères (hors ponctuation). \textbf{Dans la fourchette.} \\ \hline
\end{tabularx}
\end{center}
\textbf{Variante « Moins fort » (niveau passage) :} Phrases plus courtes, vocabulaire HSK 2 (\zh{有意思}, \zh{漂亮}, \zh{难}), 1 seul argument, concession omise, \zh{因为}...\zh{所以} basique. $\approx$ 60 caractères. \textbf{Conseil :} Viser le modèle « Riche » pour l'excellence.
\end{exoresolu}

\begin{attention}[Les 5 erreurs qui coûtent des points au Bac (ZHA22 Spécial)]
\begin{enumerate}
    \item \textbf{Tons \& Pinyin :} Confondre \zh{书} (shū, 1) et \zh{树} (shù, 4) / \zh{熟} (shú, 2). Écrire \zh{yi si} au lieu de \zh{yìsi} (意识 vs 意思). \textbf{Règle :} \zh{一} (yī) devient \zh{yí} (2e ton) devant 4e ton (\zh{一起} qǐ), \zh{yì} (4e ton) devant 1/2/3e ton (\zh{一样} yàng). \zh{不} (bù) devient \zh{bú} (2e ton) devant 4e ton (\zh{不是} shì).
    \item \textbf{Caractères « Jumeaux » (Radicaux) :} \zh{艺} (art, \zh{艹}) $\neq$ \zh{培} (cultiver, \zh{土}). \zh{雕} (sculpter, \zh{石}) $\neq$ \zh{吊} (pendre, \zh{口}). \zh{法} (méthode, \zh{氵}+\zh{去}) $\neq$ \zh{泛} (déborder, \zh{氵}+\zh{凡}). \zh{特} (spécial, \zh{牛}+\zh{寺}) $\neq$ \zh{持} (tenir, \zh{扌}+\zh{寺}). \textbf{Réflexe :} Identifier le radical SÉMANTIQUE.
    \item \textbf{Ordre des mots (Compléments) :} \textbf{Français :} « Je vais à Pékin \underline{demain} \underline{apprendre} la calligraphie. » $\rightarrow$ \textbf{Chinois ERRONÉ :} \zh{我去北京明年学书法。} $\rightarrow$ \textbf{CORRECT :} \zh{我\underline{明年}去北京学书法。} (Temps avant Verbe). \textbf{Lieu avant Verbe :} \zh{我在喀麦隆学中文。} (Pas \zh{我学中文在喀麦隆}).
    \item \textbf{Mots-outils \zh{的 / 得 / 地} :} 
    \begin{itemize}
        \item \zh{的} (DE) : Devant Nom (Adj + \zh{的} + N). \textit{Ex: \zh{漂亮的画}.}
        \item \zh{得} (DE) : Après Verbe (V + \zh{得} + Adj/Degré). \textit{Ex: \zh{画得很好}.}
        \item \zh{地} (DE) : Devant Verbe (Adv + \zh{地} + V). \textit{Ex: \zh{认真地写}.}
    \end{itemize}
    \textbf{Piège Bac :} \zh{书法很难写} (pas de \zh{得}) vs \zh{书法写得很难} (V+\zh{得}+Adj).
    \item \textbf{Structure \zh{比} + Adverbes de degré :} \textbf{Interdit :} \zh{A 比 B 很/非常/更 Adj}. \textbf{Correct :} \zh{A 比 B Adj}. Pour « beaucoup plus » : \zh{A 比 B Adj \underline{得多} / 多了}. Pour « un peu plus » : \zh{A 比 B Adj \underline{一点儿/些}}. \textbf{Calque français :} « Plus difficile » $\neq$ \zh{更难} dans structure \zh{比}.
\end{enumerate}
\end{attention}

\courssub{Éléments socioculturels : Cameroun et Chine}

\begin{savaistu}[Arts traditionnels : regards croisés Cameroun -- Chine]
\textbf{Objectif :} Situer les pratiques artistiques dans leur contexte culturel pour enrichir l'argumentation (niveau B1/B2).
\begin{itemize}
    \item \textbf{Chine -- 书法 (Shūfǎ) \& 绘画 (Huìhuà) :} Arts lettrés (文人画), liés au confucianisme (文化), au taoïsme (自然) et au bouddhisme (禅). La calligraphie = « peinture du cœur » (心画). Outils : \zh{文房四宝} (pinceau, encre, papier, pierre à encre). \textit{Parallèle :} Comme le \zh{ndop} (tissu bamiléké) ou les masques, l'art chinois code l'identité et le statut.
    \item \textbf{Cameroun -- 雕刻 (Diāokè) \& 面具 (Miànjù) :} Diversité ethnique (Bamiléké, Bamoun, Béti, Sawa...). Sculpture sur bois (\zh{木雕}), bronze (Foumban), terre cuite. Fonctions : rituel, royauté, initiation, décoration. \textit{Exemple :} Le \zh{管} (pipe) bamiléké, les masques \zh{Ngondo} (Sawa), les perles \zh{珠子}.
    \item \textbf{Musique :} Chine -- \zh{古筝} (cithare), \zh{二胡} (violine à 2 cordes), \zh{琵琶} (luth). Pentatonique, ornementation, évocation nature. Cameroun -- \zh{巴拉风} (balafon, xylophone à calebasses), \zh{鼓} (tambours parlants), \zh{姆 vet} (harpe). Rythmique complexe, polyphonie, lien langue/parole (tonal).
    \item \textbf{Patrimoine vivant :} Les deux pays classent des éléments à l'UNESCO (ex: \zh{京剧} Jīngjù, \zh{剪纸} Jiǎnzhǐ / \zh{Ngondo}, \zh{Patrimoine Foumban}). Transmission : maître-élève (口传心授) vs initiation societé secrète / atelier familial.
\end{itemize}
\textbf{Piste débat :} « La modernité menace-t-elle l'authenticité de l'art traditionnel ? » (Arguments : marché touristique, perte de sens, mais aussi innovation, visibilité mondiale).
\end{savaistu}

\courssub{Activité orale : Débat guidé}

\begin{experience}[Débat : Quel art pour moi ? (课堂辩论 : 我pick哪一个？)]
\textbf{Durée :} 15 min. \textbf{Format :} Groupes de 4 (2 pour Chine, 2 pour Cameroun) + 1 modérateur par groupe.
\textbf{Consigne :} « Défendez votre préférence en 3 arguments structurés (Vocabulaire leçon + Connecteurs). Le modérateur note les connecteurs utilisés sur une grille. »
\begin{enumerate}
    \item \textbf{Préparation (3 min) :} Chaque côté remplit une fiche : \zh{观点} (Thèse), \zh{论据1/2/3} (Arguments + \zh{因为}...), \zh{反驳预判} (Anticiper concession \zh{虽然}...).
    \item \textbf{Débat (8 min) :} Tour 1 : Thèse. Tour 2 : Argument 1 (Chine) / Arg 1 (Cameroun). Tour 3 : Argument 2... Tour 4 : Concession + Réaffirmation. Tour 5 : Conclusion.
    \item \textbf{Synthèse (4 min) :} Modérateur présente au tableau : « Les 3 meilleurs arguments de chaque côté » + « Le connecteur le plus utilisé / le plus oublié ».
\end{enumerate}
\textbf{Grille modérateur (extrait) :}
\begin{center}
\begin{tabularx}{\linewidth}{|X|X|X|}
\hline \rowcolor{popblueL} Connecteur visé & Entendu (Cocher) & Phrase exemple notée \\ \hline
首先 / 其次 & \square & \\ \hline
因为 / 所以 & \square & \\ \hline
虽然 / 但是 & \square & \\ \hline
而且 / 也 & \square & \\ \hline
总之 & \square & \\ \hline
\end{tabularx}
\end{center}
\textbf{Différenciation :} Niveau A2 : phrases toutes faites fournies. Niveau B1 : notes autorisées. Niveau B2 : improvisation totale.
\end{experience}

\courssub{Entraînement personnel}

\begin{exercice}[Exercices de révision et d'approfondissement (ZHA22)]
\textbf{Exercice 1 -- Vocabulaire \& Caractères (10 min)}
\begin{enumerate}
    \item Complétez par le bon caractère (radical entre parenthèses) :
    \begin{itemize}
        \item Cette \_\_\_ (艹) est très profonde. (艺术 / 艺术家)
        \item Il pratique la \_\_\_ (石) sur bois depuis 5 ans. (雕刻 / 吊刻)
        \item Le \_\_\_ (田) de ce peintre est célèbre. (画 / 画家)
    \end{itemize}
    \item Associez l'adjectif à l'œuvre (une seule réponse possible par ligne) :
    \begin{center}
    \begin{tabular}{ll}
    1. Musique de Guzheng & A. 热闹 (rènao, animé) \\
    2. Masque Ngondo & B. 优美 (yōuměi, gracieux) \\
    3. Balafon / Tam-tam & C. 有特色 (yǒu tèsè, typique) \\
    4. Calligraphie & D. 安静 (ānjìng, calme) \\
    \end{tabular}
    \end{center}
    \item Traduisez en chinois (caractères + pinyin) :
    \begin{itemize}
        \item « La sculpture sur bois est plus difficile que la peinture. »
        \item « Je trouve l'opéra de Pékin très intéressant. »
    \end{itemize}
\end{enumerate}

\textbf{Exercice 2 -- Grammaire \& Structures (15 min)}
\begin{enumerate}
    \item Reformulez les phrases en utilisant la structure \zh{比} (supériorité) :
    \begin{itemize}
        \item \zh{中国书法难，喀麦隆雕刻不难。} $\rightarrow$ \_\_\_\_\_\_\_\_\_\_
        \item \zh{我觉得喀麦隆音乐有意思，中国音乐没意思。} $\rightarrow$ \_\_\_\_\_\_\_\_\_\_
    \end{itemize}
    \item Complétez avec \zh{的 / 得 / 地} :
    \begin{itemize}
        \item 他画画画\_\_\_很好看。
        \item 这是一幅\_\_\_很有特色\_\_\_画。
        \item 她认真\_\_\_练习书法。
    \end{itemize}
    \item Reliez par \zh{虽然}...\zh{但是}... :
    \begin{itemize}
        \item \zh{汉字很难写。} + \zh{我喜欢写汉字。}
        \item \zh{喀麦隆雕刻很有生命力。} + \zh{我想学中国书法。}
    \end{itemize}
\end{enumerate}

\textbf{Exercice 3 -- Expression écrite (20 min -- Épreuve type)}
\begin{itemize}
    \item Sujet : \textbf{« 你喜欢中国书法还是喀麦隆雕刻？请写一段话 (80--100 caractères) pour expliquer ton choix. »}
    \item Contraintes : Utiliser \zh{首先/其次} (ou \zh{而且}), \zh{虽然}...\zh{但是}, \zh{因为}...\zh{所以}, \zh{至少 3 adjectifs d'appréciation}, \zh{respecter l'ordre des mots}.
\end{itemize}
\end{exercice}

\begin{corrige}[Exercices d'entraînement ZHA22]
\textbf{Exercice 1}
\begin{enumerate}
    \item \zh{艺术} (yìshù), \zh{雕刻} (diāokè), \zh{画} (huà) / \zh{画家} (huàjiā).
    \item 1-B (优美), 2-C (有特色), 3-A (热闹), 4-D (安静).
    \item \zh{木雕比绘画难。} (Mùdiāo bǐ huìhuà nán.) / \zh{我觉得京剧很有意思。} (Wǒ juéde Jīngjù hěn yǒu yìsi.)
\end{enumerate}
\textbf{Exercice 2}
\begin{enumerate}
    \item \zh{中国书法比喀麦隆雕刻难。} / \zh{我觉得喀麦隆音乐比中国音乐有意思。}
    \item \zh{得} (V+得+Adj), \zh{的} (Adj+的+N), \zh{的} (Adj+的+N -- \zh{很有特色的画}), \zh{地} (Adv+地+V).
    \item \zh{虽然汉字很难写，但是我喜欢写汉字。} / \zh{虽然喀麦隆雕刻很有生命力，但是我想学中国书法。}
\end{enumerate}
\textbf{Exercice 3 -- Proposition de corrigé (Modèle)}
\begin{textebox}[Corrigé type]
我叫李明，在杜阿拉上学。我喜欢喀麦隆雕刻，\textbf{但是}我更爱中国书法。\\
\textbf{首先}，写书法很安静，能培养专注力。\\
\textbf{其次}，汉字有文化内涵，一笔一画很有意思。\\
\textbf{虽然}木雕很有生命力，\textbf{但是}书法让我心静。\\
\textbf{因为}我想提高耐心，\textbf{所以}我每天练习。\\
\textbf{总之}，中国书法是我的最爱。
\end{textebox}
\textit{Note :} 94 caractères. Structure complète, vocabulaire HSK 3, connecteurs variés, caractères ciblés (\zh{雕, 艺, 法, 内, 涵, 伟}) corrects.
\end{corrige}

\newpage

\coursec{Exercices}

\courssub{Je vérifie mes connaissances}

\begin{exercice}[QCM : arts et opinions]
Pour chaque question, choisis la bonne réponse (une seule).
\begin{enumerate}
\item \zh{书法} signifie : a) la peinture \quad b) la calligraphie \quad c) la sculpture
\item \zh{面具} signifie : a) le masque \quad b) le tableau \quad c) la danse
\item Pour exprimer une opinion, on dit : a) \zh{我觉得…} \quad b) \zh{我有…} \quad c) \zh{我去…}
\item Dans une question offrant un choix entre deux options, on utilise : a) \zh{吗} \quad b) \zh{还是} \quad c) \zh{了}
\item \zh{因为…所以…} sert à exprimer : a) le temps \quad b) la cause et la conséquence \quad c) la comparaison
\end{enumerate}
\end{exercice}

\begin{exercice}[Vrai ou faux — justifie ta réponse]
Réponds par vrai ou faux, puis justifie en une phrase en français.
\begin{enumerate}
\item \zh{我喜欢中国艺术。} signifie « Je n'aime pas l'art chinois ».
\item On peut dire \zh{我很觉得书法很美。}
\item Le caractère \zh{术} contient un point en haut à droite du caractère \zh{木}.
\item \zh{还是} peut s'employer dans une phrase affirmative.
\item \zh{首先} signifie « d'abord / en premier lieu ».
\end{enumerate}
\end{exercice}

\begin{exercice}[Vocabulaire : complète le tableau]
Recopie et complète le tableau avec les mots manquants (caractères, pinyin ou traduction).

\begin{center}
\begin{tabularx}{\linewidth}{|X|X|X|X|}
\hline
\rowcolor{popblueL} Caractères & Pinyin & Nature & Traduction \\
\hline
艺术 & \ldots\ldots & nom & l'art \\
\hline
\ldots\ldots & shūfǎ & nom & la calligraphie \\
\hline
画 & \ldots\ldots & nom / verbe & \ldots\ldots \\
\hline
\ldots\ldots & miànjù & nom & le masque \\
\hline
觉得 & \ldots\ldots & verbe & \ldots\ldots \\
\hline
\ldots\ldots & yīnwèi & conjonction & parce que \\
\hline
\end{tabularx}
\end{center}
\end{exercice}

\begin{exercice}[Associe le pinyin aux caractères]
Associe chaque mot en caractères (colonne A) à son pinyin (colonne B).

\begin{center}
\begin{tabularx}{\linewidth}{|X|X|}
\hline
\rowcolor{popblueL} A — Caractères & B — Pinyin \\
\hline
1. 喜欢 & a) yǒumíng \\
\hline
2. 历史 & b) xǐhuan \\
\hline
3. 有名 & c) chuántǒng \\
\hline
4. 传统 & d) lìshǐ \\
\hline
5. 漂亮 & e) piàoliang \\
\hline
\end{tabularx}
\end{center}
\end{exercice}

\courssub{Je m'entraîne}

\begin{exercice}[Thème : traduis en chinois]
Traduis les phrases suivantes en chinois (caractères + pinyin).
\begin{enumerate}
\item J'aime l'art camerounais.
\item La calligraphie chinoise est très belle.
\item Je pense que les masques camerounais sont très intéressants.
\item Tu aimes l'art chinois ou l'art camerounais ?
\item Parce que l'art chinois a une longue histoire, il est très célèbre.
\end{enumerate}
\end{exercice}

\begin{exercice}[Version : traduis en français]
Traduis les phrases suivantes en français.
\begin{enumerate}
\item \zh{我觉得中国的书法很美。} (Wǒ juéde Zhōngguó de shūfǎ hěn měi.)
\item \zh{喀麦隆的面具很有名。} (Kāmàilóng de miànjù hěn yǒumíng.)
\item \zh{你喜欢喀麦隆还是中国艺术？} (Nǐ xǐhuan Kāmàilóng háishi Zhōngguó yìshù?)
\item \zh{中国艺术有几千年的历史。} (Zhōngguó yìshù yǒu jǐ qiān nián de lìshǐ.)
\end{enumerate}
\end{exercice}

\begin{textebox}[Texte à compléter — pour l'exercice suivant]
我\_\_\_\_\_\_喜欢中国艺术。\_\_\_\_\_\_，中国的书法很美，\_\_\_\_\_\_中国的画很有意思。\_\_\_\_\_\_，中国艺术有几千年的历史，\_\_\_\_\_\_我觉得中国艺术很有名。\\
Wǒ \_\_\_\_\_\_ xǐhuan Zhōngguó yìshù. \_\_\_\_\_\_, Zhōngguó de shūfǎ hěn měi, \_\_\_\_\_\_ Zhōngguó de huà hěn yǒu yìsi. \_\_\_\_\_\_, Zhōngguó yìshù yǒu jǐ qiān nián de lìshǐ, \_\_\_\_\_\_ wǒ juéde Zhōngguó yìshù hěn yǒumíng.
\end{textebox}

\begin{exercice}[Texte à trous : les connecteurs]
Complète le texte ci-dessus avec les connecteurs : \zh{首先}、\zh{然后}、\zh{还有}、\zh{所以}、\zh{因为} (chaque connecteur est utilisé une seule fois).
\end{exercice}

\begin{exercice}[Remets les phrases dans l'ordre]
Les phrases suivantes forment un petit texte argumentatif. Remets-les dans l'ordre logique : introduction, opinion, deux arguments, conclusion. Écris simplement les lettres dans l'ordre.
\begin{enumerate}[label=\Alph*.]
\item \zh{所以，我觉得喀麦隆艺术很有意思。} (Suǒyǐ, wǒ juéde Kāmàilóng yìshù hěn yǒu yìsi.)
\item \zh{首先，喀麦隆的面具很漂亮。} (Shǒuxiān, Kāmàilóng de miànjù hěn piàoliang.)
\item \zh{你喜欢喀麦隆还是中国艺术？} (Nǐ xǐhuan Kāmàilóng háishi Zhōngguó yìshù?)
\item \zh{然后，喀麦隆的音乐和舞蹈也很有名。} (Ránhòu, Kāmàilóng de yīnyuè hé wǔdǎo yě hěn yǒumíng.)
\item \zh{我很喜欢喀麦隆艺术。} (Wǒ hěn xǐhuan Kāmàilóng yìshù.)
\end{enumerate}
\end{exercice}

\courssub{Je me prépare au Bac}

\begin{textebox}[L'art chinois et l'art camerounais — texte support]
我叫玛丽，是雅温得一所中学的学生。我学习中文两年了。我很喜欢艺术。中国艺术有几千年的历史：书法很美，画也很有名。可是我也喜欢喀麦隆艺术，因为喀麦隆的面具非常漂亮，音乐和舞蹈也很有意思。我觉得中国艺术和喀麦隆艺术都很美，但是它们不一样。\\
Wǒ jiào Mǎlì, shì Yǎwēndé yì suǒ zhōngxué de xuésheng. Wǒ xuéxí Zhōngwén liǎng nián le. Wǒ hěn xǐhuan yìshù. Zhōngguó yìshù yǒu jǐ qiān nián de lìshǐ: shūfǎ hěn měi, huà yě hěn yǒumíng. Kěshì wǒ yě xǐhuan Kāmàilóng yìshù, yīnwèi Kāmàilóng de miànjù fēicháng piàoliang, yīnyuè hé wǔdǎo yě hěn yǒu yìsi. Wǒ juéde Zhōngguó yìshù hé Kāmàilóng yìshù dōu hěn měi, dànshì tāmen bù yíyàng.\\
\emph{Traduction de référence : Je m'appelle Marie, je suis élève dans un lycée de Yaoundé. J'apprends le chinois depuis deux ans. J'aime beaucoup l'art. L'art chinois a plusieurs milliers d'années d'histoire : la calligraphie est très belle et la peinture est aussi très célèbre. Mais j'aime aussi l'art camerounais, parce que les masques camerounais sont très beaux, et la musique et la danse sont aussi très intéressantes. Je pense que l'art chinois et l'art camerounais sont tous les deux très beaux, mais ils sont différents.}
\end{textebox}

\begin{exercice}[Compréhension de texte — type Bac]
Lis le texte support ci-dessus, puis réponds aux questions en français.

\textbf{Questions de compréhension}
\begin{enumerate}
\item Qui est Marie et où étudie-t-elle ?
\item Depuis combien de temps apprend-elle le chinois ?
\item Quels sont les deux éléments de l'art chinois cités dans le texte ?
\item Pourquoi Marie aime-t-elle aussi l'art camerounais ?
\item Selon elle, quelle est la différence entre les deux arts ?
\end{enumerate}

\textbf{Questions de langue}
\begin{enumerate}
\item Relève dans le texte deux connecteurs et explique leur fonction.
\item Dans la phrase \zh{我觉得中国艺术和喀麦隆艺术都很美}, quel mot exprime l'opinion ? Traduis la phrase.
\item Recopie les caractères \zh{艺术} et \zh{面具} en respectant l'ordre des traits, et nomme le radical principal de chaque caractère.
\end{enumerate}
\end{exercice}

\begin{exercice}[Thème et version — type Bac]
\textbf{A. Thème.} Traduis en chinois (caractères + pinyin) :
\begin{enumerate}
\item « J'aime l'art camerounais parce que les masques sont très beaux. Et toi, qu'est-ce que tu aimes ? »
\item « La calligraphie chinoise a une longue histoire, donc elle est très célèbre. »
\end{enumerate}

\textbf{B. Version.} Traduis en français le texte suivant :

\zh{中国有很美的书法和画。中国艺术有几千年的历史，所以很有名。}\\
(Zhōngguó yǒu hěn měi de shūfǎ hé huà. Zhōngguó yìshù yǒu jǐ qiān nián de lìshǐ, suǒyǐ hěn yǒumíng.)
\end{exercice}

\begin{exercice}[Expression écrite — type Bac]
\textbf{Sujet :} \zh{你喜欢喀麦隆还是中国艺术？为什么？} (Nǐ xǐhuan Kāmàilóng háishi Zhōngguó yìshù? Wèishénme?) — « Aimes-tu l'art camerounais ou l'art chinois ? Pourquoi ? »

\textbf{Consignes :}
\begin{itemize}
\item Rédige un texte en chinois de \textbf{8 phrases minimum} (environ 60 à 80 caractères).
\item Utilise \textbf{au moins 3 connecteurs} parmi : \zh{首先}、\zh{然后}、\zh{还有}、\zh{因为}、\zh{所以}、\zh{可是}.
\item Structure attendue : introduction (reprise de la question), ton opinion avec \zh{我觉得}, deux arguments, conclusion.
\item Soigne l'écriture des caractères du thème : \zh{艺术、画、书法、面具、喜欢、觉得} (radical correct + ordre des traits).
\item Barème indicatif : contenu et arguments 4 points ; correction grammaticale 4 points ; connecteurs 2 points ; écriture des caractères 2 points.
\end{itemize}
\end{exercice}

\begin{methode}[Réussir ta production écrite]
\begin{enumerate}
\item \textbf{Reprends la question} en introduction : \zh{你喜欢喀麦隆还是中国艺术？我很喜欢……}
\item \textbf{Annonce ton opinion} avec \zh{我觉得} ou \zh{我喜欢}.
\item \textbf{Développe deux arguments} avec \zh{首先}… \zh{然后}… ; ajoute un élément avec \zh{还有}.
\item \textbf{Justifie} avec \zh{因为}… \zh{所以}… ; nuance éventuellement avec \zh{可是}.
\item \textbf{Conclus} en reprenant ton opinion : \zh{所以，我觉得……}
\item \textbf{Relis-toi} : ordre Sujet + Verbe + Complément ; \zh{很} devant l'adjectif ; pas de \zh{是} devant un adjectif ; \zh{还是} uniquement dans les questions.
\end{enumerate}
\textbf{Modèle de réponse (niveau attendu) :}\\
\zh{你喜欢喀麦隆还是中国艺术？我很喜欢喀麦隆艺术。我觉得喀麦隆艺术很有意思。首先，喀麦隆的面具非常漂亮。然后，喀麦隆的音乐和舞蹈也很有名。还有，喀麦隆艺术有很长的历史。因为喀麦隆艺术很美，所以我很喜欢它。}\\
(Nǐ xǐhuan Kāmàilóng háishi Zhōngguó yìshù? Wǒ hěn xǐhuan Kāmàilóng yìshù. Wǒ juéde Kāmàilóng yìshù hěn yǒu yìsi. Shǒuxiān, Kāmàilóng de miànjù fēicháng piàoliang. Ránhòu, Kāmàilóng de yīnyuè hé wǔdǎo yě hěn yǒumíng. Háiyǒu, Kāmàilóng yìshù yǒu hěn cháng de lìshǐ. Yīnwèi Kāmàilóng yìshù hěn měi, suǒyǐ wǒ hěn xǐhuan tā.)\\
« Aimes-tu l'art camerounais ou l'art chinois ? J'aime beaucoup l'art camerounais. Je trouve que l'art camerounais est très intéressant. D'abord, les masques camerounais sont très beaux. Ensuite, la musique et la danse camerounaises sont aussi très célèbres. De plus, l'art camerounais a une longue histoire. Parce que l'art camerounais est très beau, je l'aime beaucoup. »
\end{methode}

\begin{attention}[Erreurs fréquentes des francophones]
\begin{itemize}
\item Ne dis pas \zh{我是很喜欢} pour « j'aime beaucoup » : \zh{是} ne se met pas devant un verbe d'affection. Dis \zh{我很喜欢}.
\item Ne mets pas \zh{吗} dans une question avec \zh{还是} : \zh{你喜欢中国艺术还是喀麦隆艺术？} (et non \zh{…还是…吗？}).
\item \zh{很} devant l'adjectif n'est pas toujours « très » : il sert souvent simplement de lien entre le sujet et l'adjectif (\zh{面具很漂亮} = « les masques sont beaux »).
\item N'oublie pas le point de \zh{术} et le radical \zh{艹} de \zh{艺} ; dans \zh{面}, ferme le cadre en dernier.
\item En français on évite « parce que… donc… » dans la même phrase ; en chinois, \zh{因为…所以…} est la structure normale et attendue.
\end{itemize}
\end{attention}

\newpage

\coursec{Corrigés des exercices}

\begin{corrige}[Exercice 1 — QCM : arts et opinions]
\begin{enumerate}
\item \textbf{b) la calligraphie.} \zh{书法} (shūfǎ) = « l'art d'écrire » : \zh{书} (livre, écrire) + \zh{法} (méthode, loi). La peinture se dit \zh{画} (huà), la sculpture \zh{雕塑} (diāosù).
\item \textbf{a) le masque.} \zh{面具} (miànjù) : \zh{面} (visage) + \zh{具} (instrument, objet).
\item \textbf{a)} \zh{我觉得…} (wǒ juéde…) = « je pense que…, je trouve que… ». \zh{我有} = « j'ai », \zh{我去} = « je vais » : aucun n'exprime une opinion.
\item \textbf{b)} \zh{还是} (háishi) = « ou » dans une question alternative. \zh{吗} sert aux questions fermées (oui/non) ; \zh{了} est une particule d'aspect, pas un outil de question.
\item \textbf{b)} \zh{因为…所以…} (yīnwèi… suǒyǐ…) = « parce que… donc… » : cause et conséquence.
\end{enumerate}
\end{corrige}

\begin{corrige}[Exercice 2 — Vrai ou faux]
\begin{enumerate}
\item \textbf{Faux.} \zh{我喜欢中国艺术。} (Wǒ xǐhuan Zhōngguó yìshù.) signifie « J'aime l'art chinois ». La négation serait \zh{我不喜欢中国艺术。} : la négation \zh{不} se place \emph{avant} le verbe.
\item \textbf{Faux.} \zh{很} (hěn) se place devant un \emph{adjectif} (\zh{很美} « très beau »), jamais devant le verbe \zh{觉得}. Phrase correcte : \zh{我觉得书法很美。} (Wǒ juéde shūfǎ hěn měi.)
\item \textbf{Vrai.} \zh{术} (shù) = \zh{木} (mù, bois) + un point en haut à droite, tracé en dernier. Piège d'écriture classique : sans ce point, on écrit \zh{木} (bois), pas \zh{术}.
\item \textbf{Faux.} \zh{还是} ne s'emploie que dans les \emph{questions} alternatives : \zh{你喜欢茶还是咖啡？} Dans une phrase affirmative, « ou » se dit \zh{或者} (huòzhě).
\item \textbf{Vrai.} \zh{首先} (shǒuxiān) = « d'abord, en premier lieu » ; c'est le connecteur qui ouvre l'argumentation.
\end{enumerate}
\end{corrige}

\begin{corrige}[Exercice 3 — Vocabulaire : complète le tableau]
\begin{center}
\begin{tabularx}{\linewidth}{|X|X|X|X|}
\hline
\rowcolor{popblueL} Caractères & Pinyin & Nature & Traduction \\
\hline
艺术 & yìshù & nom & l'art \\
\hline
书法 & shūfǎ & nom & la calligraphie \\
\hline
画 & huà & nom / verbe & le tableau, la peinture / peindre, dessiner \\
\hline
面具 & miànjù & nom & le masque \\
\hline
觉得 & juéde & verbe & penser, trouver (que) \\
\hline
因为 & yīnwèi & conjonction & parce que \\
\hline
\end{tabularx}
\end{center}
\textbf{Points de vigilance :} le deuxième caractère de \zh{觉得} porte le ton neutre (\emph{de}) ; \zh{画} est à la fois nom et verbe, comme en français « dessin / dessiner ».
\end{corrige}

\begin{corrige}[Exercice 4 — Associe le pinyin aux caractères]
1 -- b) \zh{喜欢} xǐhuan (aimer) ; 2 -- d) \zh{历史} lìshǐ (l'histoire) ; 3 -- a) \zh{有名} yǒumíng (célèbre) ; 4 -- c) \zh{传统} chuántǒng (la tradition) ; 5 -- e) \zh{漂亮} piàoliang (beau, joli).\\
\textbf{Astuce :} repère les syllabes caractéristiques : \zh{历} et \zh{史} sont tous deux en \emph{li/shi} avec 4\up{e} ton ; \zh{传统} contient deux 3\up{es} tons (le premier se prononce alors comme un 2\up{e} ton : \emph{chuántǒng}).
\end{corrige}

\begin{corrige}[Exercice 5 — Thème : traduis en chinois]
\begin{enumerate}
\item \zh{我喜欢喀麦隆艺术。} (Wǒ xǐhuan Kāmàilóng yìshù.) — le \zh{的} n'est pas obligatoire entre le pays et \zh{艺术} ; on peut aussi dire \zh{喀麦隆的艺术}.
\item \zh{中国的书法很美。} (Zhōngguó de shūfǎ hěn měi.) — pas de \zh{是} devant un adjectif : c'est \zh{很} qui porte la prédication.
\item \zh{我觉得喀麦隆的面具很有意思。} (Wǒ juéde Kāmàilóng de miànjù hěn yǒu yìsi.) — « intéressant » se dit \zh{有意思} (litt. « avoir du sens »).
\item \zh{你喜欢中国艺术还是喀麦隆艺术？} (Nǐ xǐhuan Zhōngguó yìshù háishi Kāmàilóng yìshù?) — question alternative avec \zh{还是}, \textbf{sans} \zh{吗} final.
\item \zh{因为中国艺术有很长的历史，所以很有名。} (Yīnwèi Zhōngguó yìshù yǒu hěn cháng de lìshǐ, suǒyǐ hěn yǒumíng.) — en chinois, \zh{因为} et \zh{所以} se combinent dans la même phrase, contrairement au français ; le sujet n'est pas répété dans la seconde proposition.
\end{enumerate}
\end{corrige}

\begin{corrige}[Exercice 6 — Version : traduis en français]
\begin{enumerate}
\item « Je pense que la calligraphie chinoise est très belle. »
\item « Les masques camerounais sont très célèbres. »
\item « Tu aimes l'art camerounais ou l'art chinois ? » — \zh{还是} signale l'alternative.
\item « L'art chinois a plusieurs milliers d'années d'histoire. » — \zh{几千年} (jǐ qiān nián) = « quelques milliers d'années » ; \zh{的} relie la durée à \zh{历史}.
\end{enumerate}
\end{corrige}

\begin{corrige}[Exercice 7 — Texte à trous : les connecteurs]
Texte complété :\\
\zh{我很喜欢中国艺术。首先，中国的书法很美，还有中国的画很有意思。因为中国艺术有几千年的历史，所以我觉得中国艺术很有名。}\\
(Wǒ hěn xǐhuan Zhōngguó yìshù. Shǒuxiān, Zhōngguó de shūfǎ hěn měi, háiyǒu Zhōngguó de huà hěn yǒu yìsi. Yīnwèi Zhōngguó yìshù yǒu jǐ qiān nián de lìshǐ, suǒyǐ wǒ juéde Zhōngguó yìshù hěn yǒumíng.)\\
Traduction : « J'aime beaucoup l'art chinois. D'abord, la calligraphie chinoise est très belle, et il y a aussi la peinture chinoise qui est très intéressante. Parce que l'art chinois a plusieurs milliers d'années d'histoire, je pense donc qu'il est très célèbre. »\\
\textbf{Justification :} le premier trou, devant le verbe \zh{喜欢}, appelle un adverbe de degré : le sens exige « beaucoup » (\zh{很}) ; \zh{首先} ouvre l'argumentation ; \zh{还有} ajoute un élément ; \zh{因为} introduit la cause et \zh{所以} la conséquence — les deux vont obligatoirement ensemble dans la même phrase. Chaque connecteur n'est utilisé qu'une fois.
\end{corrige}

\begin{corrige}[Exercice 8 — Remets les phrases dans l'ordre]
Ordre correct : \textbf{C -- E -- B -- D -- A}.\\
\textbf{C} pose la question (introduction) ; \textbf{E} donne l'opinion (\zh{我很喜欢喀麦隆艺术。}) ; \textbf{B} apporte le premier argument (\zh{首先}) ; \textbf{D} le deuxième (\zh{然后}) ; \textbf{A} conclut (\zh{所以}).\\
\textbf{Méthode :} repère les connecteurs — ils sont la colonne vertébrale du texte argumentatif : \zh{首先} ne peut pas ouvrir le texte (il faut d'abord une opinion), \zh{所以} ne peut que conclure.
\end{corrige}

\begin{corrige}[Exercice 9 — Compréhension de texte — type Bac]
\textbf{Questions de compréhension}
\begin{enumerate}
\item Marie est une élève d'un lycée (une école secondaire) de Yaoundé, au Cameroun. (\zh{雅温得一所中学的学生})
\item Elle apprend le chinois depuis deux ans. \zh{两年了} : \zh{了} marque ici une durée qui continue jusqu'au moment présent.
\item La calligraphie (\zh{书法}) et la peinture / les tableaux (\zh{画}).
\item Parce que les masques camerounais sont très beaux et que la musique et la danse sont aussi très intéressantes. (\zh{因为喀麦隆的面具非常漂亮，音乐和舞蹈也很有意思。})
\item Selon elle, les deux arts sont beaux tous les deux, mais ils sont différents (\zh{都很美，但是它们不一样}).
\end{enumerate}
\textbf{Questions de langue}
\begin{enumerate}
\item Exemples attendus : \zh{可是} (kěshì, « mais ») introduit une opposition ; \zh{因为} (yīnwèi, « parce que ») introduit la cause ; on peut aussi citer \zh{但是} (mais) et l'adverbe \zh{也} (aussi), qui marque l'addition.
\item Le mot d'opinion est \zh{觉得} (juéde, « penser, trouver »). Traduction : « Je pense que l'art chinois et l'art camerounais sont tous les deux très beaux. » \zh{都} (dōu) = « tous les deux » ; \zh{和} (hé) = « et ».
\item \zh{艺} : radical \zh{艹} (clé de l'herbe) en haut, puis \zh{乙} en dessous — on écrit de haut en bas. \zh{术} : on trace \zh{木} (horizontal, vertical, puis les deux obliques) et on ajoute le point en haut à droite \emph{en dernier}. \zh{面} : trait horizontal du haut, puis le cadre extérieur, puis l'intérieur, et on ferme le cadre en dernier. \zh{具} : on écrit d'abord \zh{目} (extérieur, traits intérieurs, fermeture), puis les deux traits du bas (radical \zh{八}).
\end{enumerate}
\textbf{Barème indicatif :} compréhension 5 points (1 par question) ; langue 3 points (1 par question) ; exactitude des réponses en français 2 points.
\end{corrige}

\begin{corrige}[Exercice 10 — Thème et version — type Bac]
\textbf{A. Thème}
\begin{enumerate}
\item \zh{我喜欢喀麦隆艺术，因为面具很漂亮。你呢？你喜欢什么？} (Wǒ xǐhuan Kāmàilóng yìshù, yīnwèi miànjù hěn piàoliang. Nǐ ne? Nǐ xǐhuan shénme?) — \zh{你呢？} (nǐ ne?) est la manière naturelle de dire « et toi ? » ; « qu'est-ce que tu aimes ? » = \zh{你喜欢什么？}.
\item \zh{中国的书法有很长的历史，所以很有名。} (Zhōngguó de shūfǎ yǒu hěn cháng de lìshǐ, suǒyǐ hěn yǒumíng.) — « donc » = \zh{所以} ; le sujet n'est pas répété dans la deuxième proposition.
\end{enumerate}
\textbf{B. Version}\\
« La Chine a une très belle calligraphie et de très beaux tableaux. L'art chinois a plusieurs milliers d'années d'histoire, c'est pourquoi il est très célèbre. »\\
\textbf{Points de vigilance :} \zh{有很长的历史} = « avoir une longue histoire » (\zh{长} cháng = long) ; ne pas traduire \zh{所以} par « parce que » — c'est la conséquence, pas la cause.\\
\textbf{Barème indicatif :} thème 2 phrases $\times$ 2 points (ordre des mots, particules, caractères) ; version 3 points (fidélité, français correct) ; pinyin 1 point.
\end{corrige}

\begin{corrige}[Exercice 11 — Expression écrite — type Bac]
\textbf{Proposition de rédaction modèle :}\\
\zh{你喜欢喀麦隆还是中国艺术？我很喜欢喀麦隆艺术。我觉得喀麦隆艺术很有意思。首先，喀麦隆的面具非常漂亮。然后，喀麦隆的音乐和舞蹈也很有名。还有，喀麦隆艺术有很长的历史。因为喀麦隆艺术很美，所以我很喜欢它。中国艺术也很美，可是我更喜欢喀麦隆艺术。}\\
(Nǐ xǐhuan Kāmàilóng háishi Zhōngguó yìshù? Wǒ hěn xǐhuan Kāmàilóng yìshù. Wǒ juéde Kāmàilóng yìshù hěn yǒu yìsi. Shǒuxiān, Kāmàilóng de miànjù fēicháng piàoliang. Ránhòu, Kāmàilóng de yīnyuè hé wǔdǎo yě hěn yǒumíng. Háiyǒu, Kāmàilóng yìshù yǒu hěn cháng de lìshǐ. Yīnwèi Kāmàilóng yìshù hěn měi, suǒyǐ wǒ hěn xǐhuan tā. Zhōngguó yìshù yě hěn měi, kěshì wǒ gèng xǐhuan Kāmàilóng yìshù.)\\
« Aimes-tu l'art camerounais ou l'art chinois ? J'aime beaucoup l'art camerounais. Je trouve que l'art camerounais est très intéressant. D'abord, les masques camerounais sont très beaux. Ensuite, la musique et la danse camerounaises sont aussi très célèbres. De plus, l'art camerounais a une longue histoire. Parce que l'art camerounais est très beau, je l'aime beaucoup. L'art chinois est aussi très beau, mais je préfère l'art camerounais. »\\
\textbf{Barème indicatif (12 points) :} contenu et arguments 4 points ; correction grammaticale (ordre des mots, \zh{很}, \zh{的}, absence de \zh{是} devant l'adjectif) 4 points ; connecteurs (au moins 3, bien employés) 2 points ; écriture des caractères du thème 2 points.\\
\textbf{Points de vigilance :} ne pas écrire \zh{我是很喜欢} ; pas de \zh{吗} avec \zh{还是} ; ne pas oublier le point de \zh{术} ni le radical \zh{艹} de \zh{艺} ; \zh{因为…所以…} peut et doit se combiner ; compter ses phrases (8 minimum) et vérifier que chaque connecteur n'apparaît qu'aux bons endroits.
\end{corrige}

\newpage

\coursec{Fiche bilan}

\begin{popfigure}
\begin{tikzpicture}[
    mindmap,
    concept color=popblue,
    level 1/.style={level distance=4.5cm, sibling angle=360/7},
    level 2/.style={level distance=3cm, sibling angle=45},
    every node/.style={font=\small, align=center},
    branch1/.style={concept color=poporange, grow=30},
    branch2/.style={concept color=popgreen, grow=90},
    branch3/.style={concept color=poppurple, grow=150},
    branch4/.style={concept color=poppink, grow=210},
    branch5/.style={concept color=popgold, grow=270},
    branch6/.style={concept color=popblue, grow=330},
    branch7/.style={concept color=popmuted, grow=-30},
]
\node[concept, text width=3.2cm, minimum size=3.2cm] {Culture \& Art\\Cameroun vs Chine\\ (ZHA22)}
    child[branch1] { node[concept, text width=2.2cm] {Lexique\\Arts \& Opinions} }
    child[branch2] { node[concept, text width=2.2cm] {Structures\\比 \ 还是 \ 因为} }
    child[branch3] { node[concept, text width=2.2cm] {Connecteurs\\Logiques} }
    child[branch4] { node[concept, text width=2.2cm] {Caractères\\Clés \& Traits} }
    child[branch5] { node[concept, text width=2.2cm] {Thème / Version\\Guidés} }
    child[branch6] { node[concept, text width=2.2cm] {Production\\Écrite Modèle} }
    child[branch7] { node[concept, text width=2.2cm] {Oral / Débat\\Socio-culturel} };
\end{tikzpicture}
\legende{Carte mentale de la leçon ZHA22 — Expression écrite : Culture chinoise et camerounaise}
\end{popfigure}

\begin{bilan}

\textbf{1. Lexique clé — Arts, Culture et Opinions (HSK 2--3)}

\begin{center}
\begin{tabularx}{\linewidth}{|X|X|X|X|}
\hline
\rowcolor{popblueL} \textbf{Caractères} & \textbf{Pinyin} & \textbf{Nature} & \textbf{Traduction} \\
\hline
艺术 & yìshù & n. & art \\
音乐 & yīnyuè & n. & musique \\
舞蹈 & wǔdǎo & n. & danse \\
绘画 & huìhuà & n. & peinture / dessin \\
书法 & shūfǎ & n. & calligraphie \\
传统 & chuántǒng & n./adj. & tradition / traditionnel \\
现代 & xiàndài & n./adj. & modernité / moderne \\
文化 & wénhuà & n. & culture \\
喀麦隆 & Kāmàilóng & n.p. & Cameroun \\
中国 & Zhōngguó & n.p. & Chine \\
比较 & bǐjiào & adv. & relativement / assez \\
更 & gèng & adv. & encore plus \\
最 & zuì & adv. & le plus \\
喜欢 & xǐhuan & v. & aimer \\
觉得 & juéde & v. & trouver / penser (impression) \\
认为 & rènwéi & v. & considérer / estimer (opinion) \\
因为 & yīnwèi & conj. & parce que \\
所以 & suǒyǐ & conj. & donc / c'est pourquoi \\
但是 & dànshì & conj. & mais \\
而且 & érqiě & conj. & de plus / en plus \\
\hline
\end{tabularx}
\end{center}

\textbf{2. Règles de grammaire essentielles (à connaître par cœur)}

\begin{center}
\begin{tabularx}{\linewidth}{|X|X|X|}
\hline
\rowcolor{popblueL} \textbf{Structure / Règle} & \textbf{Exemple (Caractères + Pinyin)} & \textbf{Traduction} \\
\hline
Comparatif \textbf{比} (bǐ) : A 比 B + Adj (sans 很) & 我认为中国书法比绘画难。 & Je trouve la calligraphie chinoise plus difficile que la peinture. \\
\hline
Choix / Alternative \textbf{还是} (háishì) en question & 你喜欢喀麦隆音乐还是中国音乐？ & Préfères-tu la musique camerounaise ou la musique chinoise ? \\
\hline
Cause \textbf{因为}... \textbf{所以}... (yīnwèi... suǒyǐ...) & 因为传统文化很有意思，所以我喜欢书法。 & Comme la culture traditionnelle est intéressante, j'aime la calligraphie. \\
\hline
Opinion : \textbf{觉得} (impression) vs \textbf{认为} (jugement réfléchi) & 我觉得舞蹈累。 / 我认为书法重要。 & Je trouve la danse fatigante. / J'estime que la calligraphie est importante. \\
\hline
Degré superlatif \textbf{最} (zuì) & 喀麦隆的足球最有名。 & Le football du Cameroun est le plus célèbre. \\
\hline
Degré comparatif supériorité \textbf{更} (gèng) & 中国书法更难。 & La calligraphie chinoise est encore plus difficile. \\
\hline
Structure \textbf{Subject + Adj + 的 + Noun} (attribut) & 我喜欢传统的音乐。 & J'aime la musique traditionnelle. \\
\hline
\end{tabularx}
\end{center}

\textbf{3. Méthodes express}

\end{bilan}

\begin{methode}[Rédiger un paragraphe d'opinion (Modèle « Sandwich »)]
\begin{enumerate}
\item \textbf{Accroche / Sujet} : 你喜欢喀麦隆还是中国艺术？(Nǐ xǐhuan Kāmàilóng háishì Zhōngguó yìshù ?)
\item \textbf{Choix clair} : 我喜欢... (Wǒ xǐhuan...) / 我更喜欢... (Wǒ gèng xǐhuan...)
\item \textbf{Argument 1 (Cause)} : 因为... (Yīnwèi...) — vocabulaire précis (traditionnel, moderne, rythme, trait).
\item \textbf{Argument 2 (Nuance / Contraste)} : 而且... (Érqiě...) / 但是... (Dànshì...)
\item \textbf{Conclusion (Conséquence)} : 所以... (Suǒyǐ...) / 总之... (Zǒngzhī...)
\item \textbf{Vérification} : Tons ? Ordre des mots (Sujet - Temps - Lieu - Verbe - Complément) ? Caractères (traits/radicaux) ? Ponctuation chinoise ?
\end{enumerate}
\end{methode}

\begin{methode}[Traduction Thème (Fr $\rightarrow$ Zh) \& Version (Zh $\rightarrow$ Fr)]
\begin{enumerate}
\item \textbf{Analyser} : Identifier le sujet, le verbe, les compléments, le temps/aspect (了, 过, 在).
\item \textbf{Structurer} : Appliquer l'ordre chinois (Sujet + Temps + Lieu + Manière + Verbe + Complément).
\item \textbf{Lexique} : Choisir le mot juste (ex: \emph{art} $\rightarrow$ 艺术 ; \emph{danse} $\rightarrow$ 舞蹈 ; \emph{calligraphie} $\rightarrow$ 书法).
\item \textbf{Grammaire} : Insérer 比, 还是, 因为/所以, 的/得/地 correctement.
\item \textbf{Relire} : Vérifier la nature du mot (adj vs verbe), les classificateurs, la ponctuation pleine chasse.
\end{enumerate}
\end{methode}



\begin{aretenir}[Checklist avant le Bac — Leçon ZHA22]
$\square$ Je connais par cœur le vocabulaire des arts (音乐, 舞蹈, 绘画, 书法, 艺术) et des opinions (喜欢, 觉得, 认为, 比较, 更, 最). \\
$\square$ J'écris correctement les 10 caractères clés de la leçon (ordre des traits, radicaux : \zh{艺} \emph{yì} [艹], \zh{书} \emph{shū} [乙], \zh{法} \emph{fǎ} [氵], \zh{传} \emph{chuán} [亻], \zh{统} \emph{tǒng} [纟], \zh{现} \emph{xiàn} [王], \zh{代} \emph{dài} [人], \zh{麦} \emph{mài} [麦], \zh{隆} \emph{lóng} [阝]). \\
$\square$ Je maîtrise la structure comparative \textbf{A 比 B Adj} (sans \zh{很}) et l'alternative \textbf{还是} dans les questions. \\
$\square$ J'utilise fluidement la corrélation \textbf{因为... 所以...} pour justifier mon avis. \\
$\square$ Je distingue \zh{觉得} (impression sensorielle) et \zh{认为} (opinion intellectuelle). \\
$\square$ J'emploie au moins 3 connecteurs logiques (\zh{而且}, \zh{但是}, \zh{也}, \zh{然后}) pour structurer mon paragraphe. \\
$\square$ Je respecte l'ordre des mots chinois (STLVO : Sujet - Temps - Lieu - Verbe - Objet/Complément). \\
$\square$ Je gère les particules structurelles \zh{的} (nom), \zh{得} (verbe/degré), \zh{地} (verbe/manière) sans confusion. \\
$\square$ Ma production écrite fait 80--100 caractères, ponctuation chinoise (。 ， ？ ！) incluse, écriture lisible. \\
$\square$ Je suis capable de traduire 5 phrases types (Thème \& Version) sans dictionnaire en 10 minutes. \\
$\square$ Je peux présenter oralement mon opinion (1 min) avec tons corrects et vocabulaire du thème. \\
$\square$ Je cite au moins un élément culturel précis (ex: \zh{狮子舞} \emph{shīziwǔ}, \zh{恩东} \emph{Ēndòng}, \zh{京剧} \emph{jīngjù}, \zh{面具} \emph{miànjù}) pour comparer Cameroun et Chine.
\end{aretenir}

\findecours

\end{document}
